% concatenated from futurecover.tex


































































\documentclass[12pt,oneside,english]{amsart}


%\theoremstyle{definition}
%\newtheorem{example}{Example}

























\usepackage{tcolorbox}
\usepackage{graphicx}
\usepackage{caption}
\newcommand{\comm}[2]{\begin{tcolorbox}[colback=gray!40]\textbf{Comment} (#1) : #2\end{tcolorbox}}
%\include{Macros.tex}
\newcommand{\B}{\euscript{B}}
\textwidth 6.5in \oddsidemargin 0in \textheight 9in \topmargin
-0.5in
\newcommand{\La}{\mathcal L}
\renewcommand{\baselinestretch}{1}
\renewcommand{\arraystretch}{1}


\usepackage{amsmath}
\usepackage{amssymb}
%\usepackage{amsthm}
\usepackage{hyperref}
\usepackage{enumerate}
\usepackage{graphicx}
\usepackage{color}
\usepackage{xcolor}
\usepackage{mathtools}
\usepackage{float}
\usepackage{tikz}
\usetikzlibrary{arrows}
\usepackage{bm}
\usepackage[title]{appendix}
\usepackage{mathrsfs}
\DeclareMathAlphabet\euscript{U}{eus}{m}{n}
%\tikzset{my style/. style=circle, draw, font=Large}
\usepackage{hyperref}
\hypersetup{
	colorlinks = true, %Colours links instead of ugly boxes
	urlcolor = blue, %Colour for external hyperlinks
	linkcolor = red, %Colour of internal links
	citecolor = blue %Colour of citations
}

%\newcommand{\O}{\mbox{O}}

%the following are copied from Tom's writeup
\newcommand{\N}{\mathbb{N}}
\newcommand{\Z}{\mathbb{Z}}
\newcommand{\Q}{\mathbb{Q}}
\newcommand{\R}{\mathbb{R}}
\newcommand{\C}{\mathbb{C}}
\newcommand{\cO}{\mathcal{O}}
\newcommand{\cL}{\mathcal{L}}
\DeclareMathOperator {\Fol}{Fol}
%the above are copied from Tom's writeup


%the following are copied from Klaus's weiteup
\usepackage{mathdots}
\usepackage{amsthm}
\usepackage{tikz,amsmath}
\usetikzlibrary{arrows}
\usetikzlibrary{calc}
\usepackage{tikz-cd}
\usetikzlibrary{decorations.pathreplacing}
\tikzset{%
	show curve controls/.style={
		postaction={
			decoration={
				show path construction,
				curveto code={
					\draw [blue,-] 
					(\tikzinputsegmentfirst) -- (\tikzinputsegmentsupporta)
					(\tikzinputsegmentlast) -- (\tikzinputsegmentsupportb);
					\fill [red, opacity=0.5] 
					(\tikzinputsegmentsupporta) circle [radius=.2ex]
					(\tikzinputsegmentsupportb) circle [radius=.2ex];
				}
			},
			decorate
		}
	},
	scc/.style={
		% show curve controls
	},
}
\usepackage{epsfig}
\usepackage[normalem]{ulem}
\usepackage[all,line,
%dvips % dlf, ikke noedvendig mere, kraever nyere LaTeX installation
]{xy} 
\CompileMatrices

\newcommand{\A}{{\mathbb A}}
\newcommand{\fg}{{finitely generated }}
\newcommand{\Aff}{{\operatorname{Aff}}}
\newcommand{\dist}{{\operatorname{dist}}}
%  Frequently used functions, functors, etc.
\newcommand{\Spec}{\operatorname{Spec}}
\newcommand{\gr}{\operatorname{gr}}
\newcommand{\Hom}{\operatorname{Hom}} 
\newcommand{\Diag}{\operatorname{Diag}}  
\newcommand{\Pext}{\operatorname{Pext}}
\newcommand{\Ad}{\operatorname{Ad}}
\newcommand{\CH}{\operatorname{CH}} 
\newcommand{\Log}{\operatorname{Log}} 
\newcommand{\id}{\operatorname{id}} 
\newcommand{\Aut}{\operatorname{Aut}}
\newcommand{\eq}[1]{\begin{equation}#1\end{equation}}
\newcommand{\eQ}[1]{\begin{equation*}#1\end{equation*}}
\newcommand{\spl}[1]{\begin{split}#1\end{split}}
\newcommand{\Ext}{\operatorname{Ext}} 
\newcommand{\diag}{\operatorname{diag}}
\newcommand{\cpah}{completely positive asymptotic homomorphism}
\newcommand{\dcpah}{discrete completely positive asymptotic homomorphism}
\newcommand{\ah}{asymptotic homomorphism }
\newcommand{\cone}{\operatorname{cone}}
\newcommand{\Bott}{\operatorname{Bott}}
\newcommand{\im}{\operatorname{im}}
\newcommand{\co}{\operatorname{co}}
\newcommand{\Span}{\operatorname{Span}}
\newcommand{\Tr}{\operatorname{Tr}}
\newcommand{\diam}{\operatorname{diam}}
\newcommand{\Int}{\operatorname{Int}}
\newcommand{\Per}{\operatorname{Per}}
\newcommand{\Pec}{\operatorname{Pec}}
\newcommand{\Depth}{\operatorname{Depth}} 
\newcommand{\period}{\operatorname{period}} 
\newcommand{\mit}{\operatorname{ MIN\left(\complement L\right)   }}
\newcommand{\ins}{\operatorname{int}} 
\newcommand{\Rec}{\operatorname{Rec}}
\newcommand{\supp}{\operatorname{supp}}
\newcommand{\con}{\operatorname{context}}
\newcommand{\gd}{\operatorname{gcd}}
%\newcommand{\R}{\operatorname{\mathcal R^-}}
\newcommand{\Det}{\operatorname{Det}}
\newcommand{\Dim}{\operatorname{Dim}}
\newcommand{\Rank}{\operatorname{Rank}}
\newcommand{\Is}{\operatorname{Is}}
\newcommand{\alg}{\operatorname{alg}}
\newcommand{\plus}{\operatorname{+}}
\newcommand{\Prim}{\operatorname{Prim}}
\newcommand{\tr}{\operatorname{tr}}
\newcommand{\inv}{^{-1}}
%\newcommand{\N}{{\mathbb{N}}}
%\newcommand{\Z}{{\mathbb{Z}}}
\newcommand{\T}{{\mathbb{T}}}
\newcommand{\cip}{{\mathcal{M}}}
\newcommand{\cipp}{{\cip_{\per}}}
\newcommand{\cipa}{{\cip_{\aper}}}
\newcommand{\ev}{\operatorname{ev}}
\newcommand{\Sec}{\operatorname{Sec}}
\newcommand{\sign}{\operatorname{sign}}
\newcommand{\pt}{\operatorname{point}}
\newcommand{\cof}{\operatorname{cof}}
\newcommand{\rec}{\operatorname{rec}}
\newcommand{\pec}{\operatorname{pec}}
\newcommand{\p}{\operatorname{per}}
\newcommand{\G}{\operatorname{[[G]]}}
\newcommand{\F}{\operatorname{Fol}}

\newcommand{\val}{\operatorname{val}}
\newcommand{\INT}{\operatorname{int}}
\newcommand{\Ro}{\operatorname{RO}}
\newcommand{\Orb}{\operatorname{Orb}}
\newcommand{\Sp}{\operatorname{Sp}}
\newcommand{\Cay}{\operatorname{Cay}}
\newcommand{\Geo}{\operatorname{Geo}}
%\newcommand{\N}{\mathbb{N}}
\newcommand{\words}[1][]{\mathbb{W}_{#1}(S)}
\newcommand{\Br}{\operatorname{Br}}
\newcommand{\Wan}{\operatorname{Wan}}
\newcommand{\Ray}{\operatorname{Ray}}
\newcommand{\Mray}{\operatorname{Mray}}
\newcommand{\cl}{\operatorname{Cl}}
\newcommand{\Rray}{\operatorname{Rray}}
\newcommand{\Rr}{\operatorname{Rr}}
\newcommand{\Id}{\operatorname{Id}}

%  Shortcuts for finite sums and products
\newcommand{\nsum}{{\sum_{i=1}^n}}
\newcommand{\ncap}{{\cap_{i=1}^n}}
\newcommand{\bncap}{{\bigcap_{i=1}^n}}
\newcommand{\bnoplus}{{\bigoplus_{i=1}^n}}
\newcommand{\noplus}{{\oplus_{i=1}^n}}
\newcommand{\nprod}{{\prod_{i=1}^n}}
\newcommand{\ncup}{{\cup_{i=1}^n}}
\newcommand{\bncup}{{\bigcup_{i=1}^n}}

%\DeclareMathOperator{\supp}{supp}
\DeclareMathOperator{\orb}{Orb}
\DeclareMathOperator{\per}{Per}
\DeclareMathOperator{\aper}{Aper}

%the above are copied from Klaus's writeup


\newtheorem{thm}{Theorem}[section]
\newtheorem{corollary}[thm]{Corollary}
\newenvironment{cor}{\begin{corollary}\rm}{\end{corollary}}
\newtheorem{lemma}[thm]{Lemma}
\newtheorem{ques}[thm]{Question}
\newtheorem{prop}[thm]{Proposition}
\newtheorem{conj}[thm]{Conjecture}
\newtheorem{assumption}[thm]{Assumption}
\newenvironment{assum}{\begin{assumption}\rm}{\end{assumption}}
\newtheorem{pr}[thm]{Proposition}
\newtheorem{crit}[thm]{Criterion}
\newtheorem{prob}[thm]{Problem}
\newtheorem{definition}[thm]{Definition}
\newenvironment{de}{\begin{definition}\rm}{\end{definition}}



















\newtheorem{example}[thm]{Example}
\newenvironment{exmp}{\begin{example}\rm}{\end{example}}





\newtheorem{question}[thm]{Question}
\newenvironment{qustion}{\begin{question}\rm}{\end{question}}









\newtheorem{remark}[thm]{Remark}
\newenvironment{rem}{\begin{remark}\rm}{\end{remark}}
\newtheorem{algorithm}[thm]{Algorithm}
\newenvironment{algo}{\begin{algorithm}\rm}{\end{algorithm}}
\newtheorem{observation}[thm]{Observation}
\newenvironment{ob}{\begin{observation}\rm}{\end{observation}}
%\newtheorem{problem}[thm]{Problem}
%\newenvironment{prob}{\begin{problem}\rm}{\end{problem}}
\newtheorem{claim}[thm]{Claim}
\newenvironment{clai}{\begin{claim}\rm}{\end{claim}}
\newtheorem{fact}[thm]{Fact}
\newenvironment{fa}{\begin{fact}\rm}{\end{fact}}
\newenvironment{ex}{\begin{example}\rm}{\end{example}}
\newenvironment{re}{\begin{remark}\rm}{\end{remark}}
\DeclareMathOperator*{\argmin}{argmin}
\DeclareMathOperator*{\argmax}{argmax}
%\DeclareMathOperator*{\tr}{tr}

\newcommand\blfootnote[1]{%
	\begingroup
	\renewcommand\thefootnote{}\footnote{#1}%
	\addtocounter{footnote}{-1}%
	\endgroup
}
%\newcommand{\cor}{\color{red}}
%\newcommand{\cob}{\color{black}}

\newcommand{\blue}{\color{blue}}
\newcommand{\orange}{\color{orange}}
\newcommand{\red}{\color{red}}
\newcommand{\black}{\color{black}}
\newcommand{\vertiii}[1]{{\left\vert\kern-0.25ex\left\vert\kern-0.25ex\left\vert #1
    \right\vert\kern-0.25ex\right\vert\kern-0.25ex\right\vert}}
    %\newcommand\answerbox{%%
    %\fbox{\rule{0.5in}{0pt}\rule[-0.5ex]{0pt}{1.8ex}}}
    
    
%\usepackage{showkeys}    
    
    
\title{On the future cover of a sofic shift}



\author{Klaus Thomsen}


\subjclass[2020]{Primary ; Secondary }
\keywords{Sofic shifts, canonical covers}
\begin{document}

\begin{abstract} The paper contains a new proof of the theorem by Krieger which establishes the canonicity of the future cover of a sofic shift. In addition the paper describes a method to produce a new canonical cover from a given one, resulting in a canonical cover related to, but generally different from the future cover.
\end{abstract}
%\today
\maketitle\thispagestyle{empty}


\section{Introduction}
 
 A cover of a sofic shift $Y$ is a subshift of finite type $X$ together with a factor map $\pi : X \to Y$ from $X$ onto $Y$. Every sofic shift has a wealth of different covers, so the least one can do to simplify considerations regarding covers is to identify two covers $\pi :  X \to Y$ and $\pi' : X' \to Y$ of the same sofic shift $Y$ when they are isomorphic in the sense that there is a conjugacy $\chi : X \to X'$ such that $\pi' \circ \chi = \pi$. Recall that a subshift of finite type can be realized as the edge shift $X_G$ of a finite directed graph $G$ and that a labeling of the set of edges $E_G$ is a map $L_G : E_G \to A$, where $A$ is a finite set. If $Y = L_G(X_G)$, where $L_G : X_G \to Y \subseteq A^{\mathbb Z}$ is the map obtained by reading the labels of elements in $X_G$, we say that $(G,L_G)$ is a presentation of $Y$. Then any cover of $Y$ is isomorphic to one coming from a presentation of $Y$, see Theorem 3.2.1 in \cite{LM}. But even after having reduced covers to labeled graphs, every sofic shift has an uncanny host of different presentations. Fortunately the future cover of Wolfgang Krieger, \cite{Kr1}, and its associated labeled graph which we denote by $(\mathbb K(Y),L_{\mathbb K(Y)})$, stands out because it is canonical in the sense that a conjugacy $\psi : Y \to Z$ of sofic shifts lifts to a conjugacy $\psi_{\mathbb K} : X_{\mathbb K(Y)} \to X_{\mathbb K(Z)}$ such that
\begin{equation*}
\xymatrix{
 X_{\mathbb K(Y)} \ar[d]_-{L_{\mathbb K(Y)}} \ar[r]^-{\psi_{\mathbb K}} &  X_{\mathbb K(Z)} \ar[d]^-{L_{\mathbb K(Z)}} \\
Y \ar[r]_-{\psi}  & Z }
\end{equation*}
commutes. Thus the subshift of finite type $X_{\mathbb K(Y)}$ and the factor map $L_{\mathbb K(Y)}$ are both conjugacy invariants for $Y$. The same is of course also true for any other cover which is canonical in the same way. It must be emphasized that also components of a canonical cover that are not of maximal entropy as well as the wandering parts of the cover, which often arise even when $Y$ is mixing, also give rise to conjugacy invariants for $Y$. While any cover of a sofic shift is a window to its internal structure, a canonical cover guarantees that what you see is genuinely intrinsic. In this way canonical covers offer one of the few tools we have for the study of sofic shifts; in particular the structure that set them apart from subshifts of finite type. This observation is the main motivation for the work and the results that are presented below. 


We start by giving a new proof of Krieger's theorem. While his proof is a tour de force in handling sliding block codes, the proof given here is more conceptual and leads to an abstract characterization of the future cover analogous to the abstract characterization of the canonical cover of an irreducible sofic shift known as the minimal right-resolving cover or Fischer cover, \cite{F}. Alternative proofs of parts of Krieger's theorem have been given before; first by Boyle, Kitchens and Marcus in \cite{BKM} and later by Nasu in \cite{N}. The proof presented here is also different from those. %One virtue of the proof we give is that it allows to identify universal properties of the future cover that are used to characterize it up to isomorphism. 

The main new result of the paper is a construction which produces a new canonical cover out of a given one. To appreciate this, recall that Krieger proved the uniqueness of the conjugacy $\psi_\mathbb K$ in the diagram above. Without this additional requirement it is obvious that a canonical cover is not unique up to isomorphism; for example the disjoint union of two copies of the future cover will again be canonical in the sense used so far, but not with unique lifts of conjugacies. In the following we shall therefore refer to a construction of covers as \emph{weakly canonical} when it gives lifts of conjugacies, but possibly lifts that are not unique, while a construction which gives rise to unique lifts of conjugacies, as the future cover does, will be called a \emph{strongly canonical} cover.   The main new result of the paper, summarized by Theorem \ref{31-10-25hx}, describes a method to produce a strongly canonical cover out of a given weakly canonical one. 

  Up to now only two strongly canonical covers for sofic shifts have been known; at least when we for simplicity ignore the doubling coming from the passage from right-resolving to left-resolving covers that can be achieved by 'time-reversal'. Besides the future cover it is the minimal right-resolving cover, and only for irreducible sofic shifts. With the construction in the present paper we have now at least four different ways to construct a strongly canonical cover of an irreducible sofic shift and at least two for a general sofic shift. We finish the paper by describing the results of the constructions for the even shift, showing that all four covers are mutually different in this case.
  

 

\bigskip

{\em Acknowledgement:} I thank Brian Marcus, Tom Meyerovitch and Chengyu Wu for the collaboration that produced \cite{MMTW}. I'm particularly grateful to Brian for renewing my interest in sofic shifts, and I thank him also for showing me his solution to Exercise 3.2.8 (e) in \cite{LM}.
 
\section{Krieger's future cover and its canonicity}



\subsection{Setup and notation} 

Let $A$ be a finite set. The set $A^\mathbb Z$ of bi-infinite sequences 
$$
x = (x_i)_{i \in \mathbb Z} = \cdots x_{-3}x_{-2}x_{-1}x_0x_1x_2x_3 \cdots
$$
of elements from $A$ is a compact metric space; a metric can be defined for example as in Example 6.1.10 of \cite{LM}. The shift $\sigma$ is the homeomorphism of $A^\mathbb Z$ defined such that
$$
\sigma(x)_i := x_{i+1}.
$$
A closed subset $X \subseteq A^\mathbb Z$ is shift-invariant when $\sigma(X) = X$ and it is then called a \emph{subshift}. We shall mainly consider \emph{subshifts of finite type} (abbreviated to SFT) and \emph{sofic shifts}. See Definition 2.1.1 and Definition 3.1.3 in \cite{LM}. 




The set $\mathbb W(X)$ of \emph{words} in a subshift $X$ consists of the finite strings $a_1a_2 \cdots a_n \in A^n$ with the property that $a_1a_2 \cdots a_n  = x_1x_2\cdots x_n$ for some element $x = (x_i)_{i \in \mathbb Z} \in X$. Given an element $x \in X$ and an integer $j \in \mathbb Z$ we denote by $x_{(-\infty,j]}$ the element 
$$
\cdots x_{k}x_{k+1} x_{k+2} \cdots x_j \in A^{(-\infty,j]},
$$
and by $X(-\infty,j]$ the set
$$
X(-\infty,j] := \left\{ x_{(-\infty,j]}: \ x \in X \right\}.
$$
The symbols $x_{[j,\infty)}$ and $X[j,\infty)$ have a similar meaning. The same goes for expressions like $x_{[i,j)}$ and $x_{(i,j]}$ when $i < j$ or $x_{[i,j]}$ when $i \leq j$. The set of $\sigma$-periodic points of a subshift $X$ will be denoted by $\Per(X)$.

 
 Let $G$ be a finite directed graph with edges (or arrows) $E_G$ and vertexes $V_G$. For $e \in E_G$ let $s_G(e) \in V_G$ be the start vertex and $t_G(e)$ the terminal vertex of $e$. We shall always assume that $G$ does not contain sinks or sources; i.e. for every vertex $v$ there are arrows $e_1,e_2 \in E_G$ such that $t_G(e_1) = v = s_G(e_2)$. We extend the definition of $t_G$ to finite and left-infinite paths in $G$, and the definition of $s_G$ to finite and right-infinite paths in $G$, in the obvious way. We denote the edge shift of $G$ by $X_G$. Thus $X_G$ consists of the bi-infinite sequences 
 $$
 (e_i)_{i \in \mathbb Z} = \cdots e_{-3}e_{-2}e_{-1}e_0e_1e_2e_3 \cdots \ \in \ E_G^\mathbb Z
 $$ 
 of edges in $G$ such that $s_G(e_{i+1}) = t_G(e_i)$ for all $i \in \mathbb Z$, and the words in $\mathbb W(X_G)$ are the finite paths in $G$. We shall sometimes refer to the elements of $X_G$ as  \emph{rays} in $G$. The set $X_G$ is an SFT and every SFT is conjugate to $X_G$ for some graph $G$ by Theorem 2.3.2 in \cite{LM}.
 
Let $A$ be a finite set (the alphabet) and $L_G : E_G \to A$ a map considered as a labelling of the edges in $G$. We extend $L_G$ to finite or infinite paths in the natural way. In particular, when $x=(e_i)_{i \in \mathbb Z} \in X_G$,
\begin{align*}
&L_G(x) := \left(L_G(e_i)\right)_{i \in \mathbb Z} = \cdots L_G(e_{-3}) L_G(e_{-2}) L_G(e_{-1}) L_G(e_0)L_G(e_1)L_G(e_2) L_G(e_3) \cdots \\
&\in A^\mathbb Z .
\end{align*}
Then $(G,L_G)$ is a \emph{labeled graph}, $Y:= L_G(X_G) \subseteq A^\mathbb Z$ is a sofic shift and $(G,L_G)$ is a \emph{presentation} of $Y$, cf. \S 3.1 in \cite{LM}.
 
The first lemma deals with factor codes, right-resolving graphs and right-closing factor codes. We refer to Definition 3.3.1 and Definition 8.1.8 in \cite{LM}, respectively. The book \cite{LM} is the bible of symbolic dynamics and we shall sometimes use notions and results from \cite{LM} without explicit reference. 

\begin{lemma}\label{05-05-25x} Let $X$ be an SFT and $\pi : X \to Y$ a right-closing factor code. There is a right-resolving labeled graph $(G,L_G)$ and a conjugacy $\psi : X \to X_G$ such that

\begin{equation*}
\begin{xymatrix}{
 X \ar[rr]^\psi \ar[dr]_\pi & & X_G \ar[dl]^{L_G} \\
 & Y &}
\end{xymatrix}
\end{equation*}
commutes.
\end{lemma}
\begin{proof} Recode, first using the proof of Theorem 3.2.1 in \cite{LM} to get a labeled graph $(H,L_H)$ and a conjugacy $\psi_0 : X \to X_H$ such that

\begin{equation*}
\begin{xymatrix}{
 X \ar[rr]^{\psi_0} \ar[dr]_\pi & & X_H \ar[dl]^{L_H} \\
 & Y &}
\end{xymatrix}
\end{equation*}
commutes. Note that $L_H : X_H \to Y$ is a right-closing factor code since $\pi$ is. It follows therefore from Proposition 8.1.9 in \cite{LM} that $(H,L_H)$ is right-closing with delay as defined in Definition 5.1.4 of \cite{LM}.  Consequently Proposition 5.1.11 in \cite{LM} gives us a right-resolving labeled graph $(G,L_G)$ and a conjugacy $\psi_1 : X_H \to X_G$ such that
\begin{equation*}
\begin{xymatrix}{
 X_H \ar[rr]^{\psi_1} \ar[dr]_{L_H} & & X_G \ar[dl]^{L_G} \\
 & Y &}
\end{xymatrix}
\end{equation*}
commutes. Set $\psi :=  \psi_1 \circ \psi_0$.

\end{proof}



\subsection{Regular vertexes and rays}\label{rays} 

Let $(H,L_H)$ be a labeled graph and $Y = L_H(X_H)$ the sofic shift it presents. A vertex $v \in V_H$ is \emph{regular} when there is an element $z \in X_H$ such that $t_H(z_{(-\infty,-1]}) = v$ and 
\begin{equation}\label{11-09-25}
 \left\{ L_H(x) : \ x \in X_H{[0,\infty)}, \ s_H(x) = v \right\} = \left\{y \in Y[0,\infty) : \ L_H(z_{(-\infty,-1]})y \in Y \right\} .
\end{equation}
For any element $z \in X_H$ with $t_H(z_{(-\infty,-1]}) = v$, the inclusion 
$$
\left\{ L_H(x) : \ x \in X_H{[0,\infty)}, \ s_H(x) = v \right\} \subseteq \left\{y \in Y[0,\infty) : \ L_H(z_{(-\infty,-1]})y \in Y \right\}
$$ 
always holds, so the real condition in \eqref{11-09-25} is that the opposite inclusion,
$$
\left\{ L_H(x) : \ x \in X_H{[0,\infty)}, \ s_H(x) = v \right\} \supseteq \left\{y \in Y[0,\infty) : \ L_H(z_{(-\infty,-1]})y \in Y \right\},
$$ 
also holds. 


In order to work with regularity of vertexes we introduce the following notation. Given an element $y \in Y$, set
\begin{equation}\label{23-09-25a}
F(y) := \left\{w \in Y[0,\infty): \ y_{(-\infty,-1]}w \in Y \right\} ,
\end{equation}
and given a vertex $v \in V_H$, set
$$
f_H(v) : = \left\{L_H(x): \ x \in X_H[0,\infty), \ s_H(x) = v \right\} .
$$
Then $v \in V_H$ is regular when there is an element $z \in X_H$ such that $t_H(z_{(-\infty,-1]}) =v$ and 
\begin{equation}\label{15-09-12}
f_H(v) = F(L_H(z)) .
\end{equation}
The set $f_H(v)$ will be called the \emph{follower set} of $v$ and $F(y)$ the \emph{follower set} of $y$.


















\begin{lemma}\label{31-08-25x} Assume that $(H,L_H)$ is right-resolving. Let $\gamma$ be a finite path in $H$, and assume that $s_H(\gamma)$ is regular. Then $t_H(\gamma)$ is regular.
\end{lemma}
\begin{proof} Let $z \in X_H$ be an element such that $t_H(z_{(-\infty,-1]}) = s_H(\gamma)$ and
\begin{equation}\label{03-11-25}
 \left\{ L_H(\mu) : \ \mu \in X_H{[0,\infty)}, \ s_H(\mu) =s_H(\gamma)   \right\} = \left\{y \in Y[0,\infty) : \ L_H(z_{(-\infty,-1]})y \in Y \right\} .
\end{equation}
 Set $\gamma' := z_{(-\infty,-1]}\gamma$ and note that $\gamma' = z'_{(-\infty,-1]}$ for some $z' \in X_H$ with $t_H(z'_{(-\infty,-1]}) = t_H(\gamma)$. If $y\in Y[0,\infty)$ and $L_H(z'_{(-\infty,-1]})y = L_H( z_{(-\infty,-1]})L_H(\gamma)y \in Y$, then \eqref{03-11-25} implies that $L_H(\gamma)y = L_H(\nu)$ for some $\nu \in X_H[0,\infty)$ with $s_H(\nu) = s_H(\gamma)$. Since $(H,L_H)$ is right-resolving, $\nu = \gamma \nu'$ for some $\nu' \in X_H[0,\infty)$ with $s_H(\nu') = t_H(\gamma)$. Note that $L_H(\nu') = y$.
\end{proof}


A labeled graph will be called \emph{regular} when all vertexes are regular. Adapting Definition 3.3.7 in \cite{LM} we say that a labeled graph $(H,L_H)$ is \emph{follower-separated} when $f_H(v) = f_H(w) \Rightarrow v =w$.


\begin{ex}\label{03-10-25} The following labeled graph $(G,L_G)$ is right-resolving and follower-separated. But it is not regular; the vertex a is not regular.

\begin{figure}[H]
\begin{equation*}
\begin{tikzpicture}[node distance={25mm}, thick, main/.style = {draw, circle}] 
\node[main] (1) {$a$}; 
\node[main] (2) [right of=1] {$b$}; 
\node[main] (3) [right of=2] {$c$};
\draw[->] (1) to [out=210,in=310,looseness=7] node[below ,pos=0.4] {$1$} (1);
\draw[->] (1) to [out=0,in=180,looseness=1] node[above ,pos=0.4] {$2$} (2);
\draw[->] (2) to [out=110,in=70,looseness=15] node[above ,pos=0.4] {$1$} (2);
\draw[->] (3) to [out=210,in=270,looseness=10] node[below ,pos=0.4] {$1$} (3);
\draw[->] (3) to [out=290,in=330,looseness=15] node[below ,pos=0.4] {$4$} (3);
\draw[->] (3) to [out=180,in=0,looseness=1] node[above ,pos=0.4] {$3$} (2);
\end{tikzpicture} 
\end{equation*}
\caption{A right-resolving follower-separated labeled graph which is not regular.}
\label{04-09-25gxx}

\end{figure}
\end{ex}

\begin{re}\label{26-09-25x} 
Every synchronizing graph in the sense of Jonoska, cf. \S 4 in \cite{Jo}, is regular. By Lemma 3.3.15, Proposition 3.3.16 and Corollary 3.3.19 in \cite{LM} the minimal right-resolving presentation of an irreducible sofic shift is synchronizing in the sense of Jonoska and therefore also regular. The labeled graph in Figure \ref{11-09-25cx} depicts the future cover of the even shift and it is regular, but it is not synchronizing. To give an irreducible example consider the following labeled graph which is strongly connected, right-resolving, and presents the even shift.
\begin{figure}[H]
\begin{equation*}
\begin{tikzpicture}[node distance={25mm}, thick, main/.style = {draw, circle}] 
\node[main] (1) {$a$}; 
\node[main] (2) [right of=1] {$b$}; 
\node[main] (3) [right of=2] {$c$};
\node[main] (4) [right of=3] {$d$};
\draw[->] (1) to [out=340,in=210,looseness=1.0] node[below ,pos=0.5] {$0$} (2);
\draw[->] (2) to [out=110,in=60,looseness=1.0] node[above,pos=0.5] {$0$} (1);
\draw[->] (2) to [out=60,in=110,looseness=1.0] node[above,pos=0.5] {$1$} (3);
\draw[->] (3) to [out=200,in=340,looseness=1.0] node[below,pos=0.5] {$1$} (2);
\draw[->] (3) to [out=60,in=110,looseness=1.0] node[above,pos=0.5] {$0$} (4);
\draw[->] (4) to [out=200,in=340,looseness=1.0] node[below,pos=0.5] {$0$} (3);
\end{tikzpicture} 
\end{equation*}
\caption{A strongly connected labeled graph which is right-resolving and regular, but not synchronizing.}
\label{26-09-25}
\end{figure}
%However, labeled graphs that are right-resolving and follower-separated are regular if and only if they are synchronizing. This can be shown using Lemma \ref{01-10-25a} below. ????


%To see that this graph is regular, note that there is unique left-infinite ray $\gamma$ with label $0^\infty 1$ which terminates at the vertex b. There is also a unique left-infinite path with the same label as $\gamma$ which terminates at the vertex c, and these two paths are the only left-infinite paths in $G$ with that label. Now note that the vertexes b and c have the same follower sets. It follows that if $z \in X_G$ is such that $z_{(-\infty,-1]} = \gamma$, we have that \begin{equation}\label{26-09-25a}
 %\left\{ L_G(x) : \ x \in X_G{[0,\infty)}, \ s_G(x) = b \right\} = \left\{y \in Y[0,\infty) : \ L_G(z_{(-\infty,-1]})y \in Y \right\} ,
%\end{equation}
%where $Y = L_G(X_G)$ is the even shift. Thus c is regular and it follows from Lemma \ref{31-08-25x} that so is $G$. To see that there is no magic word for $G$, in the sense of Definition 9.1.4 of \cite{LM}, note that every word of $Y$ is the label of at least two paths with distinct terminal vertexes in $G$. Hence $G$ is not synchronizing.
\end{re}


Let $x \in X_H$. We are interested in the condition that 
\begin{equation}\label{02-09-25ax}
\left\{ L_H(\mu) : \ \mu \in X_H[0,\infty), \ s_H(\mu)  = 
t_H(x_{(-\infty, k]})\right\} = \left\{ y \in Y[0,\infty): \ L_H(x_{(-\infty,k]})y \in Y \right\}, 
\end{equation}
or 
\begin{equation}\label{23-09-25}
f_H(t_H(x_{(-\infty,k]})) = F(\sigma^{k+1}(L_H(x))),
\end{equation}
for all $k \in \mathbb Z$. Let $X$ be a subshift and $d$ a metric for the topology of $X$. For $x \in X$, set 
\begin{align*}
&\mathbb U(x) := \left\{ z \in X: \ \lim_{i \to - \infty} d\left(\sigma^i(z),\sigma^i(x)\right) = 0\right\} \\
&= \left\{z \in X: \ z_i = x_i, \ i \leq N, \ \text{for some} \ N \in \mathbb Z \right\} .
\end{align*} 
This set is sometimes called 'the unstable manifold' of $x$.

\begin{lemma}\label{02-09-25x}  Assume that $(H,L_H)$ is right-resolving. Let $x \in X_H$. Then $L_H : \mathbb U(x) \to \mathbb U(L_H(x))$ is surjective if and only if \eqref{02-09-25ax} holds for all $k \in \mathbb Z$.




\end{lemma}
\begin{proof} 

 Assume first that \eqref{02-09-25ax} holds for all $k \in \mathbb Z$ and consider an element $y \in \mathbb U(L_H(x))$. Then $y_{(-\infty,L]} = L_H(x_{(-\infty,L]})$ for some $L \in \mathbb Z$. Thanks to \eqref{02-09-25ax} this implies that $y_{[L+1,\infty)} \in f_H(t_H(x_{(-\infty,L]}))$, and hence also that there is an element $z \in X_H$ such that $z_{(-\infty,L]} = x_{(-\infty,L]}$ and $L_H(z_{[L+1,\infty)}) = y_{[L+1,\infty)}$. Then $z \in \mathbb U(x)$ and $L_H(z) = y$.

Assume then that $L_H : \mathbb U(x) \to \mathbb U(L_H(x))$ is surjective, and consider $k \in \mathbb Z$ and an element $y\in Y[0,\infty)$ such that $L_H(x_{(-\infty,k]})y \in Y$. Then $L_H(x_{(-\infty,k]})y \in \mathbb U(L_H(x))$ and there is therefore an element $z \in \mathbb U(x)$ such that $L_H(z) =L_H(x_{(-\infty,k]})y$. Choose $L < k$ such that $z_{(-\infty, L]} = x_{(-\infty, L]}$. Since $(H,L_H)$ is right-resolving, it follows that $z_{(-\infty, k]} = x_{(-\infty, k]}$. Hence 
$$
y = L_H(z_{[k+1,\infty)}) \in f_H(t_H(x_{(-\infty,k]})).
$$

\end{proof}

%We note that by the argument that proved Lemma \ref{31-08-25}, if \eqref{02-09-25ax} holds for some $k \in \mathbb Z$ then it also holds for any $k' \geq k$.



An element $x \in X_H$ will be called \emph{regular} when $L_H : \mathbb U(x) \to \mathbb U(L_H(x))$ is surjective. By Lemma \ref{02-09-25x} the following clearly holds.


\begin{lemma}\label{12-09-25b} Assume that $(H,L_H)$ is right-resolving and that $x\in X_H$ is regular. Then $s_H(x_k)\in V_H$ is regular for all $k \in \mathbb Z$.
\end{lemma}
%\begin{proof} This follows from Lemma \ref{02-09-25}.
%\end{proof}

We denote in the following the set of regular rays in $X_H$ by $\mathcal R(L_H)$.  




\begin{ex}\label{03-10-25a}
In the labeled graph $H$ depicted in Figure \ref{11-09-25cx} the vertex $c$ is regular and hence all the vertexes are regular by Lemma \ref{31-08-25x}. However, $X_H$ contains non-regular rays; the two elements of $X_H$ of period 2 are not regular.

\begin{figure}[H]
\begin{equation*}
\begin{tikzpicture}[node distance={25mm}, thick, main/.style = {draw, circle}] 
\node[main] (1) {$a$}; 
\node[main] (2) [right of=1] {$b$}; 
\node[main] (3) [below of=1] {$c$};
\draw[->] (1) to [out=80,in=110,looseness=1.0] node[above ,pos=0.5] {$0$} (2); 
\draw[->] (2) to [out=260,in=310,looseness=1.0]  node[below ,pos=0.5] {$0$}(1); 
\draw[->] (1) to [out=160,in=220,looseness=15] node[above ,pos=0.4] {$1$} (1);
\draw[->] (3) to [out=90,in=270,looseness=1.0]  node[right ,pos=0.5] {$1$} (1);
\draw[->] (3) to [out=10,in=290,looseness=10]  node[below ,pos=0.5] {$0$} (3);
\end{tikzpicture} 
\end{equation*}
\caption{A regular labeled graph with non-regular rays.}
\label{11-09-25cx}
\end{figure}
\end{ex}




%Recall from Defintion 3.3.7 of \cite{LM} that a labeled graph $(X_H,L_H)$ follower- separated when $v,w \in V_H$, $f_H(v) =f_H(w) \ \Rightarrow \ v =w$.

\begin{lemma}\label{02-09-25bx} Assume that $(X_H,L_H)$ is right-resolving and follower-separated. Then $L_H$ is injective on $\mathcal R(L_H)$.
\end{lemma}
\begin{proof} If $x,y \in \mathcal R(L_H)$ and $L_H(x) = L_H(y)$, it follows first from \eqref{23-09-25} that $t_H(x_k) = t_H(y_k)$ for all $k \in \mathbb Z$ because we assume that $(X_H,L_H)$ is follower-separated. Since $(H,L_H)$ is right-resolving it follows that $x=y$.
\end{proof}








\subsection{The future cover and Krieger's theorem} 


Let $A$ be a finite set and 
$Y \subseteq A^\mathbb Z$ a sofic subshift. To define the future cover of $Y$ we use the following notation. When $D \subseteq A^\mathbb N$ is a subset of $A^{\mathbb N}$ and $a \in A$, we let $\frac{D}{a}$ denote the set
$$
\frac{D}{a} := \left\{x_1x_2x_3 \cdots \in A^\mathbb N : \ x_0x_1x_2\cdots \in D, \ x_0 =a \right\} = \left\{x \in A^\mathbb N: \ ax \in D\right\};
$$
i.e. $\frac{D}{a}$ is the set of elements of $A^{\mathbb N}$ that are obtained from the elements of $D$ that start with $a$ by deleting the first coordinate. More generally, when $w \in A^m$, we set
$$
\frac{D}{w} = \left\{x \in A^{\mathbb N} : \ wx \in D \right\}.
$$


Let $y \in Y$. Recall from \eqref{23-09-25a} that
$$
F(y)  = \left\{ z_{[0,\infty)}: \ z \in Y, \ z_j = y_j, \ j \leq -1 \right\} \subseteq A^{\mathbb N}. 
$$
Set
$$
V_{\mathbb K(Y)} := \left\{F(y): \ y \in Y\right\} .
$$
This is a finite set because $Y$ is sofic. $V_{\mathbb K(Y)}$ is the set of vertexes in a labeled graph $(\mathbb K(Y),L_{\mathbb K(Y)})$ where there is a labeled arrow
$$
F(y) \overset{a}{\to} F(z)
$$
from $F(y)$ to $F(z)$ in $\mathbb K(Y)$ precisely when 
$F(z) = \frac{F(y)}{a}$.
The labeled graph $(\mathbb K(Y),L_{\mathbb K(Y)})$ is then clearly right-resolving. Since 
$$
F(\sigma^{i+1}(y)) = \frac{F(\sigma^{i}(y))}{y_i},
$$
there is an element $\alpha_Y(y) \in X_{\mathbb K(Y)}$ determined by the two conditions:
\begin{itemize}
\item[a)] $s_{\mathbb K(Y)}(\alpha_Y(y)_i) = F(\sigma^i(y)), \ i \in \mathbb Z$, and
\item[b)] $L_{\mathbb K(Y)}(\alpha_Y(y)) = y$.
\end{itemize}
The map $\alpha_Y : Y \to X_{\mathbb K(Y)}$ commutes with the shift, i.e. $\alpha_Y \circ \sigma = \sigma \circ \alpha_Y$, but it is not continuous.

\begin{lemma}\label{04-09-25f} $L_{\mathbb K(Y)}(X_{\mathbb K(Y)}) =Y$.
\end{lemma}
\begin{proof} Let $x = (x_i)_{i \in \mathbb Z} \in X_{\mathbb K(Y)}$ and consider an $N \in \mathbb N$. By definition $s_{\mathbb K(Y)}(x_{-N}) = F(y)$ for some $y \in Y$ and
$$
L_{\mathbb K(Y)}(x_{-N})L_{\mathbb K(Y)}(x_{-N+1}) L_{\mathbb K(Y)}(x_{-N+2}) \cdots  \ \ \in F(y) .
$$
In particular,
$$
L_{\mathbb K(Y)}(x_{-N})L_{\mathbb K(Y)}(x_{-N+1}) \cdots L_{\mathbb K(Y)}(x_{N}) \in \mathbb W(Y)
$$
for all $N \in \mathbb N$, implying that $L_{\mathbb K(Y)}(X_{\mathbb K(Y)}) \subseteq Y$. The reversed inclusion, $Y \subseteq L_{\mathbb K(Y)}(X_{\mathbb K(Y)})$, follows from the observation b) above.
\end{proof}

Thus $(\mathbb K(Y),L_{\mathbb K(Y)})$ is a presentation of $Y$, and it is known as \emph{the future cover} of $Y$, cf. Exercise 3.2.8 on page 75 of \cite{LM}. \footnote{Krieger introduced the word 'future' in \cite{Kr1}, and called an equivalent version of $(\mathbb K(Y),L_{\mathbb K(Y)})$ 'the future state chain'. The terminology "future cover" was used in \cite{BKM}, but only for a component of $(\mathbb K(Y),L_{\mathbb K(Y)})$, and only for transitive sofic shifts. In Exercise 3.2.8 of \cite{LM} the quotation marks were deleted and $(\mathbb K(Y),L_{\mathbb K(Y)})$ was called the future cover for general sofic shifts. In some papers on sofic shifts the future cover is called the Krieger cover and the minimal right-resolving cover is called the Fischer cover.  }





We can now formulate Krieger's theorem on the canonicity of the future cover. In our notation Theorem 2.14 in \cite{Kr1} reads as follows: 

\begin{lemma}\label{03-02-25ax} (Theorem 2.14 of \cite{Kr1})  Let $Y$ and $Z$ be sofic shifts and $\psi : Y \to Z$ a conjugacy. There is a unique conjugacy $\psi_{\mathbb K} : X_{\mathbb K(Y)} \to X_{\mathbb K(Z)}$ such that
\begin{itemize}
\item[(a)] $\psi \circ L_{\mathbb K(Y)} = L_{\mathbb K(Z)} \circ \psi_{\mathbb K}$, and
\item[(b)] $\psi_{\mathbb K} \circ \alpha_Y = \alpha_{Z} \circ \psi$.
\end{itemize}
\end{lemma}




When $Y$ (and $Z$) are transitive, Krieger was able to partly dispense with condition (b) in Lemma \ref{03-02-25ax}, cf. Corollary 2.16 in \cite{Kr1}. And in Theorem 2.2 of \cite{Kr2} he removed condition (b) entirely by showing that it follows from (a). In this way he proved

\begin{thm}\label{Krieger} (W. Krieger, 1984, 1987)
Let $Y$ and $Z$ be sofic shifts and $\psi : Y \to Z$ a conjugacy. There is a unique conjugacy $\psi_{\mathbb K} : X_{\mathbb K(Y)} \to X_{\mathbb K(Z)}$ such that
\begin{equation*}\label{24-11-24}
\xymatrix{
 X_{\mathbb K(Y)} \ar[d]_-{L_{\mathbb K(Y)}} \ar[r]^-{\psi_{\mathbb K}} &  X_{\mathbb K(Z)} \ar[d]^-{L_{\mathbb K(Z)}} \\
Y \ar[r]_-{\psi}  & Z }
\end{equation*}
commutes. This conjugacy $\psi_{\mathbb K}$ has the property that also
\begin{equation}\label{24-11-24x}
\xymatrix{
 X_{\mathbb K(Y)} \ar[r]^-{\psi_{\mathbb K}} &  X_{\mathbb K(Z)} \\
Y \ar[r]_-\psi \ar[u]^-{\alpha_Y} & Z \ar[u]_-{\alpha_{Z}} }
\end{equation}
commutes.
\end{thm}



Note that \eqref{24-11-24x} says that the right-inverse $\alpha_Y$ of $L_{\mathbb K(Y)}$ is canonical in the same sense as $(\mathbb K(Y),L_{\mathbb K(Y)})$. It can therefore also be used to get conjugacy invariants for $Y$. 


Before we can give the proof of Theorem \ref{Krieger} we need some preparations. First of all we need the relation
\begin{equation}\label{11-09-25h}
f_{\mathbb K(Y)}(F(y)) = F(y)
\end{equation}
in which $F(y)$ is first considered as a vertex in $\mathbb K(Y)$ and next as a subset of $A^{\mathbb N}$. In more detail,


\begin{lemma}\label{02-10-24fx} $(\mathbb K(Y),L_{\mathbb K(Y)})$ is follower-separated. In fact, for all $y \in Y$,
$$
\left\{L_{\mathbb K(Y)}(w) :  \ w \in X_{\mathbb K(Y)}[0,\infty), \ s_{\mathbb K(Y)}(w) =  F(y)\right\} = F(y).
$$
\end{lemma}
\begin{proof} This follows directly from the definition of $\mathbb K(Y)$.
\end{proof}


\begin{lemma}\label{11-09-25g} $\alpha_Y(y)$ is regular in $X_{\mathbb K(Y)}$ for all $y \in Y$, i.e.
$\alpha_Y(Y) \subseteq \mathcal R(L_{\mathbb K(Y)})$.
\end{lemma}
\begin{proof} Since $L_{\mathbb K(Y)}(\alpha_Y(y)) = y$ we must check that
$$
f_{\mathbb K(Y)}\left(t_{\mathbb K(Y)}(\alpha_Y(y)_{(-\infty,k]})\right) = F(\sigma^{k+1}(y))
$$
for $k \in \mathbb Z$. This is a direct check, using \eqref{11-09-25h}: 
$$f_{\mathbb K(Y)}\left(t_{\mathbb K(Y)}(\alpha_Y(y)_{(-\infty,k]})\right) = f_{\mathbb K(Y)}\left(s_{\mathbb K(Y)}(\alpha_Y(y)_{k+1})\right)  = f_{\mathbb K(Y)}(F(\sigma^{k+1}(y))) = F(\sigma^{k+1}(y)).
$$
\end{proof}

\begin{lemma}\label{04-09-25h} $(\mathbb K(Y),L_{\mathbb K(Y)})$ is regular.
\end{lemma}
\begin{proof}  Let $y \in F(y)$. To see that $F(y)$ is a regular vertex in $\mathbb K(Y)$, recall that $(\mathbb K(Y),L_{\mathbb K(Y)})$ is right-resolving. Since $F(y) = s_{\mathbb K(Y)}(\alpha_Y(y)_0)$ the desired conclusion follows therefore by combining Lemma \ref{11-09-25g} with Lemma \ref{12-09-25b}.
\end{proof}




\begin{lemma}\label{12-09-25j} Let $\gamma : Y \to X_{\mathbb K(Y)}$ be a map such that $\gamma(Y) \subseteq \mathcal R(L_{\mathbb K(Y)})$ and $L_{\mathbb K(Y)}\circ \gamma =\id_Y$.  Then $\gamma = \alpha_Y$.   
\end{lemma}
\begin{proof} By definition and by Lemma \ref{11-09-25g} $\alpha_Y$ has the same properties as $\gamma$. Since $L_{\mathbb K(Y)}(\alpha_Y(y)) = y = L_{\mathbb K(Y)}(\gamma(y))$ for all $y \in Y$, and since $(\mathbb K(Y),L_{\mathbb K(Y)})$ is follower-separated by Lemma \ref{02-10-24fx}, it follows from Lemma \ref{02-09-25bx} that $\gamma = \alpha_Y$.
\end{proof}  



Given an element $w \in A^m$ we denote in the following by $w^\infty$ the $m$-periodic element $w^\infty$ in $\Per(A^\mathbb Z)$ with the property that $(w^\infty)_{[0,m-1]} = w$.


\begin{lemma}\label{11-09-25ix} Let $d$ be a metric for the topology of $X_{\mathbb K(Y)}$, let $x \in X_{\mathbb K(Y)}$ and let $N\in \mathbb Z$. There is a periodic point $p \in \Per(Y)$ and an element $x' \in X_{\mathbb K(Y)}$ such that $x'_i = x_i, \ i \geq N$, and $\lim_{i \to -\infty} d(\sigma^i(x'),\sigma^i(\alpha_Y(p))) = 0$. In particular, $x' \in \mathcal R(L_{\mathbb K(Y)})$.
\end{lemma}
\begin{proof} Choose $y \in Y$ such that $F(y) = s_{\mathbb K(Y)}(x_N)$. By the pigeon hole principle there are integers 
$\left\{n_i\right\}_{i \in \mathbb N}$ such that 
$$\cdots < n_3 < n_2 <n_1 < n_0 < N
$$ 
and $s_{\mathbb K(Y)}(\alpha_Y(y)_{n_j}) = s_{\mathbb K(Y)}(\alpha_Y(y)_{n_0})$ for all $j$. Let $q_j := L_{\mathbb K(Y)}(\alpha_Y(y)_{[n_j,n_0-1]}) = y_{[n_j,n_0-1]}$ and note that $p^j := q_j^{\infty} \in Y$. We claim that $F(\sigma^{n_0}(y)) \subseteq F(p^j)$ for all $j \in \mathbb N$.
To see this, let $z \in F(\sigma^{n_0}(y))$.
Since $F(\sigma^{n_0}(y)) = s_{\mathbb K(Y)}(\alpha_Y(y)_{n_j}) =s_{\mathbb K(Y)}(\alpha_Y(y)_{n_0})$ and $L_{\mathbb K(Y)}(\alpha_Y(y)_{[n_j,n_0-1]}) = q_j$, it follows that $F(\sigma^{n_0}(y)) = \frac{ F(\sigma^{n_0}(y))}{q_j}$, and hence also that
$$
F(\sigma^{n_0}(y)) = \frac{ F(\sigma^{n_0}(y))}{(q_j)^k}
$$
for all $k \in \mathbb N$. Hence
\begin{align*}
&F(\sigma^{n_0}(y)) = \left\{w \in F(\sigma^{n_0}(y)) : \ (q_j)^k w \in F(\sigma^{n_0}(y))\  \forall k \in \mathbb N \right\} ,
\end{align*} 
which implies that $p^j_{(-\infty,-1]}z \in  Y$, i.e. $z \in F(p^j)$.
 We claim next that $F(\sigma^{n_0}(y)) = F(p^l)$ for some $l \in \mathbb N$. By passing to a subsequence we may assume for this that $F(p^j) = F(p^1)$ for all $j$. If $F(p^1) \backslash F(\sigma^{n_0}(y)) \neq \emptyset$, there is an element $z \in Y[0,\infty)$ such that $y_{[n_j,n_0-1]}z \in Y[0,\infty)$ for all $j > 1$, but $y_{(-\infty,n_0-1]}z \notin Y$, which is impossible. This proves the claim. It follows that $\alpha_Y(y)_{[n_0,-1]}$ is a path in $\mathbb K(Y)$, labeled by $y_{[n_0,-1]}$, which starts at $ F(\sigma^{n_0}(y)) =F(p^l) = s_{\mathbb K(Y)}(\alpha_Y(p^l)_0)$ and ends at $s_{\mathbb K(Y)}(\alpha_Y(y)_0) = F(y)= s_{\mathbb K(Y)}(x_N)$. Set $p:= \sigma^{-n_0}(p^l)$ and 
 $$
 x':= \alpha_Y(p)_{(-\infty,n_0-1]}\alpha_Y(y)_{[n_0,-1]}x_{[N,\infty)} \in X_{\mathbb K(Y)}.
 $$ 
  Then $p$ and $x'$ have the desired properties, and $x'$ is regular because $\mathbb U(x') = \mathbb U(\alpha_Y(p))$ and $\alpha_Y(p) \in \mathcal R(L_{\mathbb K(Y)})$ by Lemma \ref{11-09-25g}. 
 



\end{proof}


\begin{corollary}\label{12-09-25} For every element $x \in X_{\mathbb K(Y)}$ and every $N \in \mathbb Z$ there is a regular ray $x' \in \mathcal R(L_{\mathbb K(Y)})$ such that $x'_i = x_i$ for $i \geq N$.
\end{corollary}
\begin{proof} This follows from Lemma \ref{11-09-25ix}.
\end{proof}

In the following we shall work with labeled-graph homomorphisms, cf. Definition 3.1.2 of \cite{LM}. Let $(G,L_G)$ and $(H,L_H)$ be labeled graphs. A \emph{labeled-graph homomorphism} $\theta : (G,L_G) \to (H,L_G)$ is a graph homomorphism $\theta: G \to H$ such that $L_H \circ \theta(e)= L_G(e)$ for all $e \in E_G$. The sliding block code $X_G \to X_H$ induced by $\theta$, which we also denote by $\theta$, has the property that $L_H \circ \theta = L_G$ on $X_G$. 


\begin{lemma}\label{02-09-25dx} Assume that $(G,L_G)$ is a right-resolving and regular presentation of $Y$. It follows that there is a labeled-graph homomorphism $\theta : (G,L_G) \to (\mathbb K(Y), L_{\mathbb K(Y)})$ such that
$$
\theta\left(\mathcal R(L_G)\right) \subseteq \mathcal R(L_{\mathbb K(Y)}).
$$
If $(G,L_G)$ also is follower-separated, $\theta : G \to \mathbb K(Y)$ is injective.
\end{lemma}
\begin{proof} Let $v \in V_G$. Since $v$ is regular there is an $x \in X_G$ such that $t_G(x_{(-\infty,-1]}) = v$ and 
$f_G(v)  = F(L_G(x))$.
Define $\theta(v) \in V_{\mathbb K(Y)}$ such that 
$$
\theta(v) := F(L_G(x)) .
$$
This is well-defined: If $x'$ has the same properties as $x$, it follows that
$F(L_G(x')) = f_G(v) = F(L_G(x))$.


Let $e \in E_G$. Choose $x \in X_G$ such that $t_G(x_{(-\infty,-1]}) = s_G(e)$ and 
$f_G(s_G(e)) = F(L_G(x))$.
Set $x'_i = x_{i+1}, \ i \leq -2$, $x'_{-1} =e$ and let $x'_i, i \geq 0$, be arbitrary subject to the condition that $x' \in X_G$. Then $t_G(x'_{(-\infty, -1]}) = t_G(e)$ and by using that $(G,L_G)$ is right-resolving it follows that
\begin{align*}
& f_G(t_G(e)) = \frac{f_G(s_G(e))}{L_G(e)}
 = \frac{\left\{y \in Y[0,\infty): \ L_G(x_{(-\infty,-1]})y \in Y \right\}}{L_G(e)} \\
 &= \left\{y \in Y[0,\infty): \ L_G(x'_{(-\infty,-1]})y \in Y \right\}.
\end{align*}
Hence $\theta(t_G(e)) = F(L_G(x'))$ and  $F(L_G(x')) = \frac{F(L_G(x))}{L_G(e)}$. There is therefore an edge $e'$ in $\mathbb K(Y)$ from $\theta(s_G(e)) =F(L_G(x)) $ to $\theta(t_G(e)) = F(L_G(x'))$ with label $L_{\mathbb K(Y)}(e') =L_G(e)$. Since $(\mathbb K(Y), L_{\mathbb K(Y)})$ is right-resolving this arrow is unique and we set $\theta(e):= e'$. Then $\theta : G \to \mathbb K(Y)$ is a labeled-graph morphism.

When $x \in \mathcal R(L_G)$, it follows from the definition of $\theta$ and the defining relation \eqref{02-09-25ax} for regular rays that $\theta(x) \in X_{\mathbb K(Y)}$ is the path
$$
\cdots  \overset{L_G(x_{k-1})}{\to}F(L_G(\sigma^k(x))) \overset{L_G(x_k)}{\to} F(L_G(\sigma^{k+1}(x))) \cdots .
$$ 
Thus $\theta(x) = \alpha_Y(L_G(x))$ and hence $\theta(x) \in \mathcal R(L_{\mathbb K(Y)})$ by Lemma \ref{11-09-25g}. 






%Since, by Lemma \ref{02-10-24fx},
%\begin{align*}
%&f_{\mathbb K(Y)}(t_{\mathbb K(Y)}(\theta(x)_{(-\infty, k]})) = f_{\mathbb K(Y)}(F(L_G(\sigma^{k+1}(x)))) = F(L_G(\sigma^{k+1}(x)) \\
%&= \left\{ y \in Y[0,\infty): \ L_G(x_{(-\infty, k]})y \in Y \right\} \\
%& = \left\{ y \in Y[0,\infty): \ L_{\mathbb K(Y)}(\theta(x)_{(-\infty, k]})y \in Y \right\}  \\
%& = F(\sigma^{k+1}(L_{\mathbb K(Y)}(\theta(x))))
%\end{align*}
%for all $k \in \mathbb Z$, we see that $\theta(x) \in \mathcal R(L_{\mathbb K(Y)})$, cf. \eqref{23-09-25}.

Finally, assume that $v,w \in V_G$ are such that $\theta(v) =\theta(w)$. There are elements $x,x' \in X_G$ such that $t_G(x_{(-\infty,-1]}) = v$ and $t_G(x'_{(-\infty,-1]}) = w$, $f_{G}(v) = F(L_G(x))$ and $f_G(w) = F(L_G(x'))$. By definition of $\theta$ this implies that $f_G(v) = \theta(v) = \theta(w) = f_G(w)$ and hence that $v=w$ when $(G,L_G)$ is follower-separated.
\end{proof}


\subsection{Proof of Krieger's theorem, Theorem \ref{Krieger}}

Let $\psi: Y \to Z$ be a conjugacy of sofic shifts. The future cover $(X_{\mathbb K(Y)},L_{\mathbb K(Y)})$ of $Y$ is isomorphic to the cover $(X_{\mathbb K(Y)},\psi \circ L_{\mathbb K(Y)})$ of $Z$ in the sense that
\begin{equation*}
\begin{xymatrix}{
X_{\mathbb K(Y)} \ar@{=}[r] \ar[d]_{L_{\mathbb K(Y)}} & X_{\mathbb K(Y)} \ar[d]^-{\psi \circ L_{\mathbb K(Y)}}  \\
Y \ar[r]^-\psi & Z }
\end{xymatrix}
\end{equation*}
commutes. Note that $\psi \circ L_{\mathbb K(Y)}$ is right-closing because $(\mathbb K(Y),L_{\mathbb K(Y)})$ is right-resolving and $\psi$ is a conjugacy. By Lemma \ref{05-05-25x} there is therefore a right-resolving labeled graph $(G,L_G)$ and a conjugacy $\phi :  X_{\mathbb K(Y)} \to X_G$ such that
\begin{equation*}
\begin{xymatrix}{
X_{\mathbb K(Y)}  \ar[r]^{\phi} \ar[d]_{\psi \circ L_{\mathbb K(Y)}} & X_{G} \ar[d]^-{L_G}  \\
Z \ar@{=}[r] & Z &}
\end{xymatrix}
\end{equation*}
commutes. Hence also
\begin{equation*}
\begin{xymatrix}{
X_{\mathbb K(Y)}   \ar[r]^{\phi} \ar[d]_{L_{\mathbb K(Y)}} & X_{G} \ar[d]^-{L_G}  \\
Y \ar[r]^-\psi & Z &}
\end{xymatrix}
\end{equation*}
commutes. If $L_{\mathbb K(Y)} : \mathbb U(x) \to  \mathbb U(L_{\mathbb K(Y)}(x))$ is surjective, then so is $L_G : \mathbb U(\phi(x)) \to   \mathbb U(L_G(\phi(x))) =\mathbb U(\psi(L_{\mathbb K(Y)}(x)))$, showing that 
\begin{equation}\label{02-09-25ex}
\phi\left(\mathcal R(L_{\mathbb K(Y)} )\right) \subseteq \mathcal R(L_G).
\end{equation}
Let $v\in V_G$ and consider an element $x \in X_G$ with $s_G(x_0) = v$. Then $x = \phi(z)$, where $z = \phi^{-1}(x)$. Since $\phi$ is a conjugacy of shift spaces, it is given by a sliding block code defined on blocks of a given size, say $D \in \mathbb N$, cf. \S 1.5 in \cite{LM}. It follows that $v = s_G(\phi(z')_0)$ for all $z' \in X_{\mathbb K(Y)}$ with the property that $z'_i = z_i$ for $i > -D$. By Corollary \ref{12-09-25} there is a regular ray $z' \in \mathcal R(L_{\mathbb K(Y)})$ such that $z'_i = z_i, \ i > -D$. Then $\phi(z') \in \mathcal R(L_G)$ by \eqref{02-09-25ex} and hence $v = s_G(\phi(z)_0) = s_G(\phi(z')_0)$ is regular by Lemma \ref{12-09-25b}. This shows that $(G,L_G)$ is regular. It follows then from Lemma \ref{02-09-25dx} that there is a labeled-graph homomorphism $\theta_0 : (G,L_G) \to (\mathbb K(Z),L_{\mathbb K(Z)})$ such that $\theta_0 (\mathcal R(L_G)) \subseteq \mathcal R(L_{\mathbb K(Z)})$ and
\begin{equation*}
\begin{xymatrix}{
X_{G}   \ar[r]^{\theta_0} \ar[d]_{L_G} & X_{\mathbb K(Z)} \ar[d]^-{L_{\mathbb K(Z)}}  \\
Z \ar@{=}[r] & Z }
\end{xymatrix}
\end{equation*}
commutes. All together we get a sliding block code $\phi' := \theta_0 \circ \phi :   X_{\mathbb K(Y)} \to X_{\mathbb K(Z)}$ such that 
\begin{equation*}\label{02-09-25gx}
\phi'\left(\mathcal R(L_{\mathbb K(Y)}) \right) \subseteq \mathcal R(L_{\mathbb K(Z)})
\end{equation*}
and
\begin{equation*}
\begin{xymatrix}{
X_{\mathbb K(Y)}   \ar[r]^{\phi'} \ar[d]_{L_{\mathbb K(Y)}} & X_{\mathbb K(Z)} \ar[d]^-{L_{\mathbb K(Z)}}  \\
Y \ar[r]_\psi & Z }
\end{xymatrix}
\end{equation*}
commutes.

Now apply this with $\psi$ replaced by $\psi^{-1}$. This gives us a sliding block code $\phi'' :   X_{\mathbb K(Z)} \to X_{\mathbb K(Y)}$ such that 
\begin{equation*}\label{02-09-25h}
\phi''\left(\mathcal R(L_{\mathbb K(Z)} )\right) \subseteq \mathcal R(L_{\mathbb K(Y)})
\end{equation*}
and
\begin{equation*}
\begin{xymatrix}{
X_{\mathbb K(Z)}   \ar[r]^{\phi''} \ar[d]_{L_{\mathbb K(Z)}} & X_{\mathbb K(Y)} \ar[d]^-{L_{\mathbb K(Y)}}  \\
Z \ar[r]_{\psi^{-1}} & Y }
\end{xymatrix}
\end{equation*}
commutes. Then
\begin{equation*}\label{02-09-25hx}
\phi'' \circ \phi'\left(\mathcal R(L_{\mathbb K(Y)}) \right) \subseteq \mathcal R(L_{\mathbb K(Y)})
\end{equation*}
and
\begin{equation*}
\begin{xymatrix}{
X_{\mathbb K(Y)}   \ar[r]^{\phi''\circ \phi'} \ar[d]_{L_{\mathbb K(Y)}} & X_{\mathbb K(Y)} \ar[d]^-{L_{\mathbb K(Y)}}  \\
Y \ar@{=}[r] & Y }
\end{xymatrix}
\end{equation*}
commutes. Let $x \in \mathcal R(L_{\mathbb K(Y)})$. Then $L_{\mathbb K(Y)}(x) = L_{\mathbb K(Y)}(\phi'' \circ \phi'(x))$ and $\phi'' \circ \phi'(x) \in \mathcal R(L_{\mathbb K(Y)})$. Since $(\mathbb K(Y),L_{\mathbb K(Y)})$ is right-resolving and follower-separated, Lemma \ref{02-09-25bx} implies that $x =\phi'' \circ \phi'(x)$. Since $\mathcal R(L_{\mathbb K(Y)})$ is dense in $X_{\mathbb K(Y)}$ by Corollary \ref{12-09-25}, it follows that $\phi'' \circ \phi' = \id$. By exchanging the roles of $Y$ and $Z$ it follows that also $\phi' \circ \phi'' = \id$. Set $\psi_{\mathbb K} := \phi'$.

Uniqueness: If $\phi_i : \ X_{\mathbb K(Y)} \to  X_{\mathbb K(Z)}$, $i =1,2$, are conjugacies such that $L_{\mathbb K(Z)} \circ \phi_i = \psi \circ L_{\mathbb K(Y)}$, then
$$
\phi_i\left(\mathcal R(L_{\mathbb K(Y)})\right) = \mathcal R(L_{\mathbb K(Z)}).
$$
Let $x \in \mathcal R(L_{\mathbb K(Y)})$. Since $L_{\mathbb K(Z)} \circ \phi_1(x) = \psi \circ L_{\mathbb K(Y)}(x) = L_{\mathbb K(Z)} \circ \phi_2(x)$ and $L_{\mathbb K(Z)}$ is injective on $\mathcal R(L_{\mathbb K(Z)})$ by Lemma \ref{02-09-25bx}, it follows that $\phi_1(x) = \phi_2(x)$. Since $\mathcal R(L_{\mathbb K(Y)})$ is dense in $X_{\mathbb K(Y)}$ by Corollary \ref{12-09-25}, it follows that $\phi_1 = \phi_2$.

Finally, to get the commuting diagram \eqref{24-11-24x}, set $\gamma:= \psi_\mathbb K^{-1} \circ \alpha_{Z} \circ \psi$. By using Lemma \ref{11-09-25g}, and that $\psi_\mathbb K^{-1}$ and $\psi$ both are conjugacies, we find that $\gamma(Y) \subseteq \mathcal R(L_{\mathbb K(Y)})$. Since
$$
L_{\mathbb K(Y)}\circ \gamma = \psi^{-1} \circ L_{\mathbb K(Z)} \circ \alpha_{Z} \circ \psi = \psi^{-1} \circ \psi = \id_Y = L_{\mathbb K(Y)} \circ \alpha_Y ,
$$
it follows from Lemma \ref{12-09-25j} shows that $\gamma = \alpha_Y$. \qed


As noted in the introduction an alternative proof of Theorem \ref{Krieger}, without the uniqueness part and the diagram \eqref{24-11-24x}, was given by Nasu in \cite{N}.

\subsection{The terminal component in the future cover of an irreducible sofic shift}

The material in this section is probably folklore to experts, but much of it I haven't found in the literature.

Let $H$ be a finite directed graph without sinks or sources. A vertex $v \in V_H$ is \emph{recurrent} when there is a path $\gamma$ in $H$ such that $s_H(\gamma) = t_H(\gamma) = v$, and it is called \emph{transient} otherwise. Let $V^r_H$ denote the set of recurrent vertexes and $V^t_H$ the set of transient vertexes. Two vertexes $v,w \in V_H$ \emph{communicate}, and we write $v \rightsquigarrow w$, when there is a path $\gamma$ in $H$ with $s_H(\gamma) = v$ and $t_H(\gamma) = w$. When $v,w \in V_H^r, \  v \rightsquigarrow w$ and $w \rightsquigarrow v$, we write $v \sim w$. This an equivalence relation $\sim$ on $V_H^r$ and an element $C \in V_H^r/\sim$ defines a strongly connected subgraph $H_C$ of $H$ with $V_{H_C} = C$ and $E_{H_C} = \left\{e \in E_H: \ s_H(e), \ t_H(e) \in C\right\}$. The elements of $ V_H^r/\sim$ and the corresponding graphs $H_C, \ C \in V_H^r$, will be referred to as the \emph{components} in $H$ and the irreducible subshifts $X_C := X_{H_C}, \ C \in V_H^r/\sim$, as the components of $X_H$.
 
 The components $C \in V_H^r/\sim$ are vertexes in a directed graph $C_H$ where there is an arrow from $C$ to $C'$ when $v \rightsquigarrow w$ for some elements $v\in C, \ w \in C'$. While there are no sinks or sources in $H$ by assumption, the graph $C_H$ of its communicating components always have both. A \emph{terminal component} (\emph{source component}) in $H$ is a sink (source) $C$ in $C_H$ and the associated subshift $X_C$ is a \emph{terminal component} (\emph{source component}) in $X_H$. 


Consider then a labeled graph $(H,L_H)$. A subset $U \subseteq V_H$ of vertexes in $H$ is a \emph{hereditary} subset when $e \in E_H, \ s_H(e) \in U \ \Rightarrow \ t_H(e) \in U$. When this holds we let $H_U$ denote the graph with $V_{H_U} := U$ and $E_{H_U} := \left\{ e\in E_H: \ s_H(e) \in U\right\}$. By restriction the labeling $L_H$ defines a labeling of $H_U$ and we denote the resulting labeled graph by $(H_U,L_H)$. We call this a \emph{heritary labeled subgraph} of $(H,L_H)$. Note that a terminal component $C \subseteq V_H$ of $H$ is a hereditary subset.





\begin{lemma}\label{01-10-25} Let $(H,L_H)$ be a right-resolving and follower-separated labeled graph and $C$ a terminal component in $H$. Then $(H_C,L_H)$ is isomorphic, as a labeled graph, to the minimal right-resolving presentation of $L_H(X_C)$
\end{lemma}
\begin{proof} Note that $(H_C,L_H)$ is  right-resolving and follower-separated. The lemma follows therefore from Corollary 3.3.19 in \cite{LM}. 
\end{proof}

Let $Y$ be a sofic subshift.
When $w \in \mathbb W(Y)$, set
$$
F(w) := \left\{z \in Y[0,\infty): \ wz \in Y[0,\infty) \right\} .
$$

A dual version of the following lemma has appeared as Lemma 4.7 of \cite{J}.

\begin{lemma}\label{27e-09-25} For every $y \in Y$ there is an $N \in \mathbb N$ such that 
$$
F(y_{[-n,-1]})= F(y)
$$
for all $n \geq N$.
\end{lemma}
\begin{proof} Note that $F(y_{[-n-1,-1]}) \subseteq F(y_{[-n,-1]})$ for all $n$ and that
$$
F(y) = \bigcap_n F(y_{[-n,-1]}) .
$$
Since $\left\{F(w) : \ w \in \mathbb W(Y)\right\}$ is a finite set there must be a sequence $n_1 < n_2 < \cdots$ in $\mathbb N$ such that $F(y_{[-n_i,-1]}) =F(y_{[-n_1,-1]})$ for all $i$. Set $N := n_1$. 
\end{proof}


 A word $w \in \mathbb W(Y)$ is \emph{synchronizing} when 
$u,v \in \mathbb W(Y), \ uw,wv \in \mathbb W(Y) \ \Rightarrow \ uwv \in \mathbb W(Y)$.
\begin{lemma}\label{01-10-25a} Let $(H,L_H)$ be a right-resolving, regular and follower-separated labeled graph and $C$ a terminal component in $H$. For each $v \in C$ there is a finite path $\gamma$ in $H$ such that $t_H(\gamma) = v$ and such that every path in $H$ with label $L_H(\gamma)$ terminates at $v$. In particular, $L_H(\gamma)$ is synchronizing for $L_H(X_H)$.
\end{lemma}
\begin{proof} Let $v \in C$. Since $(H,L_H)$ is regular there is a ray $z \in X_H$ such that $t_H(z_{(-\infty,-1]}) = v$ and $f_H(v) = F(L_H(z))$. By Lemma \ref{27e-09-25} there is an $N \in \mathbb N$ such that $F(L_H(z_{[-N,-1]})) = F(L_H(z))$. By Proposition 3.3.16 in \cite{LM} there is a word $u \in \mathbb W(L_H(X_H))$ such that the paths in $H$ with label $L_H(z_{[-N,-1]})u$ all terminate at the same vertex. Since $u$ is a prefix of an element from $F(L_H(z)) = f_H(v)$ there is a path $\gamma_1$ in $H$ such that $s_H(\gamma_1) =v$ and $L_H(\gamma_1) = u$. Note that $t_H(\gamma_1) \in C$ since $C$ is a terminal component. Since $H_C$ is strongly connected there is a path $\gamma_2$ in $H_C$ such that $s_H(\gamma_2) = t_H(\gamma_1)$ and $t_H(\gamma_2) = v$. Set $\gamma :=z_{[-N,-1]}\gamma_1\gamma_2$ and note that every path in $H$ with label $L_H(\gamma)$ must terminate at $v$ since $L_H(z_{[-N,-1]})u$ is a prefix of $L_H(\gamma)$. 
\end{proof}






\begin{prop}\label{29-09-25bx} Let $Y$ be an irreducible sofic shift. There is only one terminal component $C$ in $\mathbb K(Y)$ and $(\mathbb K(Y)_C,L_{\mathbb K(Y)})$ is isomorphic to the minimal right-resolving presentation of $Y$.
\end{prop}
\begin{proof} Let $v,w$ both be vertexes of a terminal component in $\mathbb K(Y)$. By Lemma \ref{01-10-25a} there are synchronizing words $s,t$ for $Y$ such that $v = F(s)$ and $w = F(t)$. Let $u \in \mathbb W(Y)$. Since $Y$ is irreducible there is a word $x \in \mathbb W(Y)$ such that $ u =x_{[i,j]}$ for some $i \leq j$ and $sxt \in \mathbb W(Y)$. By definition of $\mathbb K(Y)$ there is therefore a path in $\mathbb K(Y)$ with label $xt$ from $v=F(s)$ to $w = F(t)$. It follows that $v$ and $w$ lie in the same component $C$ of $\mathbb K(Y)$. Furthermore, since the word $u$ was arbitrary it follows that $L_{\mathbb K(Y)}(X_C) = Y$. It follows then from Lemma \ref{01-10-25} that $(\mathbb K(Y)_C,L_{\mathbb K(Y)})$ is isomorphic to the minimal right-resolving presentation of $Y$.
\end{proof}

In the following we let $(\mathbb F(Y),L_{\mathbb F(Y)})$ be the minimal right-resolving presentation of an irreducible sofic shift $Y$, also known as the \emph{Fischer cover} of $Y$. The conclusion of Proposition \ref{29-09-25bx} can then be written as
$$
(\mathbb F(Y),L_{\mathbb F(Y)}) =  (\mathbb K(Y)_C,L_{\mathbb K(Y)}),
$$
when $C$ denotes the terminal component of $\mathbb K(Y)$.

\begin{lemma}\label{30-09-25a} Let $G$ and $H$ be finite directed graphs with unique terminal components, $C$ and $C'$, respectively. Assume that $\chi : X_G \to X_{H}$ is a conjugacy. It follows that $\chi(X_C) = X_{C'}$.
\end{lemma}
\begin{proof} The periodic points $p$ of $X_C$ are characterized by the property that for all $x \in X_G$ there is an element $y \in X_G$ and a natural numbers $k \in \mathbb N$ such that $\lim_{i \to \infty} d(\sigma^{-i}(x),\sigma^{-i}(y)) = 0$ and $\lim_{i \to \infty} d(\sigma^{i}(y),\sigma^{i+k}(p)) = 0$. There is a similar characterization of the periodic elements of $X_{C'}$ and we conclude therefore that $\chi(\Per X_C) \subseteq \Per(X_{C'})$. Since the periodic points of $X_C$ are dense in $X_C$ it follows that $\chi(X_C) \subseteq X_{C'}$. Hence, by symmetry, $\chi(X_C) = X_{C'}$.
\end{proof}

\begin{corollary}\label{30-09-25x} Let $\psi : Y \to Z$ be a conjugacy of irreducible sofic subshifts. Let $C$ and $C'$ be the terminal components in $\mathbb K(Y)$ and $\mathbb K(Z)$, respectively, and let $\psi_{\mathbb K} : X_{\mathbb K(Y)} \to X_{\mathbb K(Z)}$ be the conjugacy from Theorem \ref{Krieger}. Then $\psi_{\mathbb K}(X_C) = X_{C'}$.
\end{corollary}


\begin{corollary}\label{30-09-25b} (Corollary (2.16) in \cite{Kr1}) Let $\psi : Y \to Z$ be a conjugacy of irreducible sofic subshifts. There is unique conjugacy $\psi_{\mathbb F}: X_{\mathbb F(Y)} \to X_{\mathbb F(Z)}$ such that
\begin{equation}\label{30-09-25c}
\xymatrix{
 X_{\mathbb F(Y)} \ar[d]_-{L_{\mathbb F(Y)}} \ar[r]^-{\psi_{\mathbb F}} &  X_{\mathbb F(Z)} \ar[d]^-{L_{\mathbb F(Z)}} \\
Y \ar[r]_-{\psi}  & Z }
\end{equation}
commutes.
\end{corollary}
\begin{proof} The existence of $\psi_{\mathbb F}$ follows from Corollary \ref{30-09-25x}, Proposition  \ref{29-09-25bx} and Theorem \ref{Krieger}. Let $\chi(Y)$ denote the set of elements $x\in X_{\mathbb F(Y)}$ for which $L_{\mathbb F(Y)}^{-1}(L_{\mathbb F(Y)}(x)) = \{x\}$. It follows from Lemma \ref{01-10-25a} that $\mathbb F(Y)$ contains a finite path $\gamma$ such that every path in $\mathbb F(Y)$ with the label $L_{\mathbb F(Y)}(\gamma)$ terminates at the same vertex as $\gamma$. Since $\mathbb F(Y)$ is strongly connected every ray in $X_{\mathbb F(Y)}$ can be approximated by elements that contain $\gamma$ infinitely often to the left, and these elements are in $\chi(Y)$ because $(\mathbb F(Y),L_{\mathbb F(Y)})$ is right-resolving. Hence $\chi(Y)$ is dense $X_{\mathbb F(Y)}$.

 If $\phi : X_{\mathbb F(Y)} \to X_{\mathbb F(Z)}$ is a conjugacy such that \eqref{30-09-25c} commutes, $\phi(\chi(Y)) = \chi(Z)$. If $\phi_1$ and $\phi_2$ are both conjugacies that make the diagram \eqref{30-09-25c} commute, and $x \in \chi(Y)$, we have that $\phi_i(x) \in \chi(Z)$ and 
$$
L_{\mathbb F(Z)}(\phi_1(x)) = \psi(L_{\mathbb F(Y)}(x)) = L_{\mathbb F(Z)}(\phi_2(x)),
$$
implying that $\phi_1(x) = \phi_2(x)$. Since $\chi(Y)$ is dense in $X_{\mathbb F(Y)}$ it follows that $\phi_1 = \phi_2$.
\end{proof}


As noted in the introduction, an alternative proof of Corollary \ref{30-09-25b}, without the uniqueness part, was given in \cite{BKM}.






\section{Additional properties of the future cover}




\subsection{A universal property of the future cover $(\mathbb K(Y),L_{\mathbb K(Y)})$.}

 Lemma \ref{02-09-25dx} can now be improved to the following.



\begin{prop}\label{25-09-25} Assume that $(G,L_G)$ is a right-resolving and regular presentation of the sofic shift $Y$. It follows that there is a labeled graph homomorphism 
$\theta : (G,L_G) \to (\mathbb K(Y), L_{\mathbb K(Y)})$ 
such that $
\theta\left(\mathcal R(L_G)\right) \subseteq \mathcal R(L_{\mathbb K(Y)})$,
and $\theta(V_G)$ is a hereditary subset of $V_{\mathbb K(Y)}$. If, furthermore, $(G,L_G)$ is also follower-separated, $\theta$ is a labeled graph isomorphism of $(G,L_G)$ onto the hereditary labeled subgraph $(\mathbb K(Y)_{\theta(V_G)}, L_{\mathbb K(Y)})$.
\end{prop}
\begin{proof} By Lemma \ref{02-09-25dx} and its proof there is a labeled-graph homomorphism $\theta : (G,L_G) \to (\mathbb K(Y), L_{\mathbb K(Y)})$ such that
$\theta(v) =F(L_G(x))$ where $x \in X_G$, $t_G(x_{(-\infty,-1]}) =v$ and $f_G(v) = F(L_G(x))$. Let $F(\theta(v)) \overset{a}{\to} F(z)$ be a labeled arrow in $\mathbb K(Y)$. Then $a$ is the first letter in an element $w \in F(L_G(x))) = f_G(v)$, and there is therefore a right-infinite path $\gamma$ in $G$ such that $s_G(\gamma) = v$ and $L_G(\gamma) = w$. The first arrow in $\gamma$ is an arrow $e \in E_G$ with $s_G(e) = v$ and $L_G(e) = a$. There is therefore an arrow $F(\theta(v)) \overset{a}{\to} \theta(t_G(e))$ in $\mathbb K(Y)$. Since $(\mathbb K(Y),L_{\mathbb K(Y)})$ is right-resolving this implies that $F(z) = \theta(t_G(e))$, proving that $\theta(V_G)$ is a hereditary subset of $V_{\mathbb K(Y)}$. The remaining assertions follow now from  Lemma \ref{02-09-25dx}.

 
\end{proof}


\begin{corollary}\label{29-09-25} Let $Y$ be a sofic shift. Any labeled graph which presents $Y$ and is right-resolving, regular and follower-separated is isomorphic to a hereditary labeled subgraph of $(\mathbb K(Y),L_{\mathbb K(Y)})$.
\end{corollary}

\begin{corollary}\label{02-10-25c} Let $Y$ be a sofic shift space. For any presentation $(G,L_G)$ of $Y$ which is right-resolving, regular and follower-separated, the number $\# V_G$ of vertexes is less than or equal to the number of follower sets in $Y$. The future cover $(\mathbb K(Y),L_{\mathbb K(Y)})$ is the only presentation of $Y$, up to isomorphism of labeled graphs, which is right-resolving, regular and follower-separated, and for which the number of vertexes is equal to the number of follower sets in $Y$.
\end{corollary}




\begin{ex}
If we denote the sofic shift presented by the labeled graph in Figure \ref{04-09-25gxx} by $Y$ the future cover $(\mathbb K(Y),L_{\mathbb K(Y)})$ of $Y$ is isomorphic to the following labeled graph.

\begin{figure}[H]
\begin{equation*}
\begin{tikzpicture}[node distance={25mm}, thick, main/.style = {draw, circle}] 
\node[main] (1) {$a$}; 
\node[main] (2) [below of=1] {$b$}; 
\node[main] (3) [right of=2] {$c$};
\draw[->] (1) to [out=10,in=80,looseness=7] node[above ,pos=0.4] {$1$} (1);
\draw[->] (1) to [out=260,in=100,looseness=1] node[left ,pos=0.4] {$2$} (2);
\draw[->] (1) to [out=280,in=70,looseness=1] node[right ,pos=0.4] {$3$} (2);
\draw[->] (1) to [out=320,in=120,looseness=1] node[above ,pos=0.4] {$4$} (3);
\draw[->] (3) to [out=180,in=0,looseness=1] node[below ,pos=0.5] {$3$} (2);
\draw[->] (3) to [out=200,in=260,looseness=5] node[below ,pos=0.4] {$1$} (3);
\draw[->] (3) to [out=280,in=340,looseness=5] node[below ,pos=0.4] {$4$} (3);
\draw[->] (2) to [out=200,in=270,looseness=5] node[below ,pos=0.4] {$1$} (2);
\end{tikzpicture} 
\end{equation*}
\caption{The future cover of the sofic shift presented by the labeled graph in Figure \ref{04-09-25gxx}.}
\label{03-10-25d}%\caption{A right-resolving follower-separated labeled graph with the same number of vertexes as in the future cover of the sofic shift it presents}
%\label{03-10-25c}
\end{figure}
The labeled graph in Figure \ref{04-09-25gxx} is not isomorphic to that in Figure \ref{03-10-25d} although the two graphs have the same number of vertexes. Thus regularity is a crucial assumption both in Corollary \ref{02-10-25c}, and indeed in Proposition \ref{25-09-25}.
\end{ex}


\begin{ex}\label{04-10-25} The following labeled graph $(G,L_G)$ is right-resolving, follower-separated and regular (synchronizing, in fact). It is the minimal right-resolving presentation of the  even shift.

\begin{figure}[H]
\begin{equation*}
\begin{tikzpicture}[node distance={25mm}, thick, main/.style = {draw, circle}] 
\node[main] (1) {$a$}; 
\node[main] (2) [right of=1] {$b$}; 
\draw[->] (1) to [out=80,in=110,looseness=1.0] node[above ,pos=0.5] {$0$} (2); 
\draw[->] (2) to [out=260,in=310,looseness=1.0]  node[below ,pos=0.5] {$0$}(1); 
\draw[->] (1) to [out=160,in=220,looseness=15] node[above ,pos=0.4] {$1$} (1);
\end{tikzpicture} 
\end{equation*}
\caption{A right-resolving, regular and follower-separated presentation of the even shift.}
\label{03-10-25e}
\end{figure}

The future cover of the even shift is depicted in Figure \ref{11-09-25cx}. It has one vertex more than the minimal right-resolving presentation, showing that the condition of maximality of the number of vertexes in Corollary \ref{02-10-25c} can not be dropped.


\end{ex} 
The future cover has one more property which seems worthwhile to point out; it is predecessor-separated: Consider a labeled graph $(H,L_H)$, and $v \in V_H$ a vertex in $H$. The \emph{predecessor set} of $v$ is the set
$$
p_H(v) := \left\{L_H(z): \ z \in X_H(-\infty,-1], \ t_H(z) = v \right\} .
$$
$(H,L_H)$ is \emph{predecessor-separated} when $p_H(v) = p_H(w) \ \Rightarrow \ v=w$.


\begin{lemma}\label{20-11-25f} Let $(H,L_H)$ be a right-resolving, regular and follower-separated labeled graph. Then $(H,L_H)$ is predecessor-separated.
\end{lemma}
\begin{proof} Assume $v,w \in V_H$ and $p_H(v) = p_H(w)$. Since $v$ is regular there is a $z \in X_H$ such that $t_H(z_{(-\infty,-1]}) =v$ and $f_H(v) = F(L_H(z))$. Then $L_H(z_{(-\infty,-1]}) \in p_H(v) = p_H(w)$, so there is a $u \in X_H(-\infty,-1]$ such that $t_H(u) = w$ and $L_H(u_{(-\infty,-1]}) = L_H(z_{(-\infty,-1]})$. It follows that $f_H(w) \subseteq F(L_H(u)) =F(L_H(z)) = f_H(v)$. By symmetry $f_H(v) \subseteq f_H(w)$ and hence $f_H(v) = f_H(w)$, which implies that $v = w$ since we assume that $(H,L_H)$ is follower-separated.
\end{proof}


\begin{corollary}\label{19-10-25a} The labeled graph $(\mathbb K(Y),L_{\mathbb K(Y)})$ is both follower-separated and predecessor separated.
\end{corollary}
\begin{proof} Combine Lemma \ref{02-10-24fx}, Lemma \ref{04-09-25h} and Lemma \ref{20-11-25f}.
\end{proof}

\subsection{A universal property of the future cover $L_{\mathbb K(Y)} : X_{\mathbb K(Y)} \to Y$.}

By Proposition 4 in \cite{BKM} the minimal right-resolving cover of an irreducible sofic shift has the following universal property:

\begin{prop}\label{20-11-25} Let $Y$ be an irreducible sofic shift. Let $X$ be an irreducible SFT and $\pi : X \to Y$ a right-closing factor map. There is a sliding block code $\rho : X \to X_{\mathbb F(Y)}$ such that
\begin{equation*}
\begin{xymatrix}{
 X \ar[r]^\rho \ar[d]_\pi  & X_{\mathbb F(Y)} \ar[dl]^{L_{\mathbb F(Y)}} \\
  Y &}
\end{xymatrix}
\end{equation*}
commutes.
\end{prop}

This property can be considered as a minimality condition with respect to irreducible and right-closing covers. In general, or more specifically when $Y$ is not of almost finite type, the minimal right-resolving cover is not minimal with respect to arbitrary irreducible covers by Corollary 13 in \cite{BKM}. It is also not minimal with respect to general right-resolving covers; for example the future cover of the even shift which is depicted in Figure \ref{11-09-25cx} does not factor through the minimal right-resolving cover of the even shift. However, its minimality property can be used to characterise it up to isomorphism.

\begin{prop}\label{20-11-25a} Let $Y$ be an irreducible SFT and $\pi : X \to Y$ a factor map with $X$ an SFT. Then $\pi$ is isomorphic to the minimal right-resolving cover $L_{\mathbb F(Y)} : X_{\mathbb F(Y)} \to Y$ of $Y$ if and only if
\begin{itemize}
\item[a)] $X$ is irreducible,
\item[b)] $\pi$ is right-closing,
\item[c)] $\pi$ is injective on the set of doubly transitive points in $X$, and
\item[d)] $\pi : X \to Y$ is minimal with respect to irreducible right-closing covers of $Y$.
\end{itemize}
\end{prop}
\begin{proof} The following arguments are all implicitly contained in \cite{BKM}. First, it is well-known from \cite{BKM} and \cite{LM} that the minimal right-resolving cover has all four properties, and hence so has any cover isomorphic to it. So assume that $\pi : X \to Y$ is a cover for which they all hold. We must show that $\pi$ is then isomorphic to the minimal right-resolving cover $L_{\mathbb F(Y)} : X_{\mathbb F(Y)} \to Y$ of $Y$. For this note that there are sliding block codes $\rho : X \to X_{\mathbb F(Y)}$ and $\rho' : X_{\mathbb F(Y)} \to X$ such that
\begin{equation*}
\begin{xymatrix}{
X_{\mathbb F(Y)} \ar[r]^-{\rho'} \ar[dr]_-{L_{\mathbb F(Y)}} &X \ar[r]^\rho \ar[d]_\pi  & X_{\mathbb F(Y)} \ar[dl]^{L_{\mathbb F(Y)}} \\
 &  Y &}
\end{xymatrix}
\end{equation*}
and
\begin{equation*}
\begin{xymatrix}{
X \ar[r]^-\rho \ar[dr]_-{\pi} & X_{\mathbb F(Y)} \ar[r]^{\rho'} \ar[d]_{L_{\mathbb F(Y)}}  & X \ar[dl]^{\pi} \\
  & Y &}
\end{xymatrix}
\end{equation*}
both commute. An application of Lemma 9.1.13 in \cite{LM} shows that $\rho$ and $\rho'$ both maps doubly transitive points to doubly transitive points, so condition c) implies that $\rho \circ \rho' = \id_{X_{\mathbb F(Y)}}$ and $\rho' \circ \rho = \id_X$, first on doubly transitive points and then all over by continuity. 
\end{proof}

We need some preparations in order to discuss a similar characterization of the future cover.


\begin{lemma}\label{17-11-25g} Let $(H,L_H)$ be a regular right-resolving presentation of a sofic shift. For $x \in X_H$ and $N \in \mathbb N$ there is a $y \in \mathcal R(L_H)$ such that $y_i = x_i$ for $i \geq -N$. 
\end{lemma}
\begin{proof} Let $N \in \mathbb N$. Since $s_H(x_{[-N,\infty)})$ is regular there is a $z \in X_H$ such that $t_H(z_{(-\infty,-1]}) = s_H(x_{[-N,\infty)})$ and $f_H(s_H(x_{[-N,\infty)})) = F(L_H(z))$. It follows from Lemma \ref{27e-09-25} that there is an $M \in \mathbb N$ such that $F(L_H(z)) = F(L_H(z_{[-M,-1]}))$. Set $w_1 :=z_{[-M,-1]}$. Since $s_H(w_1)$ is regular we can repeat the argument to get a finite path $w_2$ in $H$ such that $t_H(w_2) = s_H(w_1)$ and $F(L_H(w_2)) = f_H(s_H(w_1))$. Repeat this argument inductively and set
$$
y := \cdots w_3w_2w_1x_{[-N,\infty)} .
$$
For $j \in \left\{-N-\sum_{k =1}^n|w_k|: \ n \in \mathbb N\right\}$ we have by construction that
$$
f_H(t_H(y_{(-\infty,j-1]})) = F(\sigma^j(L_H(y))) .
$$
As in the proof of Lemma \ref{31-08-25x}, by using that $(H,L_H)$ is right-resolving it follows from this that the same holds for all $j \in \mathbb Z$; i.e. $y\in X_H$ is regular.  
\end{proof}

\begin{lemma}\label{17-11-25h} Let $(H,L_H)$ be a right-resolving presentation of a sofic shift. Then $(H,L_H)$ is regular if and only if $\mathcal R(L_H)$ is dense in $X_H$.
\end{lemma}
\begin{proof} One implication follows from Lemma \ref{17-11-25g}. The other from Lemma \ref{12-09-25b}.
\end{proof}


\begin{definition}\label{18-11-25} Let $\pi : X \to Y$ be a sliding block code between subshifts. An element $x \in X$ is \emph{regular} for $\pi$ when $\pi : \mathbb U(x) \to \mathbb U(\pi(x))$ is surjective. The set of 
elements $x \in X$ that are regular for $\pi$ will be denoted by $\mathcal R(\pi)$. We say that $\pi$ is \emph{regular} when $\mathcal R(\pi)$ is dense in $X$. 
\end{definition}



\begin{prop}\label{18-11-25a} Let $\pi : X \to Y$ be a factor code where $X$ is an SFT. Assume that $\pi$ is right-closing and regular. There is a sliding block code $\rho : X \to X_{\mathbb K(Y)}$ such that 
\begin{equation*}
\begin{xymatrix}{
 X \ar[r]^\rho \ar[d]_\pi  & X_{\mathbb K(Y)} \ar[dl]^{L_{\mathbb K(Y)}} \\
  Y &}
\end{xymatrix}
\end{equation*}
commutes.
\end{prop}
\begin{proof} By Lemma \ref{05-05-25x} there is a right-resolving labeled graph $(G,L_G)$ and a conjugacy $\psi : X \to X_G$ such that
\begin{equation*}
\begin{xymatrix}{
 X \ar[r]^\psi \ar[d]_\pi & X_G \ar[dl]^{L_G} \\
  Y &}
\end{xymatrix}
\end{equation*}
commutes. Note that $\psi(\mathcal R(\pi)) = \mathcal R(L_G)$. It follows from Lemma \ref{17-11-25h} that $(G,L_G)$ is regular since $\pi$ is. It follows then from Lemma \ref{02-09-25dx} that there is a sliding block code $\mu : X_G \to X_{\mathbb K(Y)}$ such that
\begin{equation*}
\begin{xymatrix}{
 X_G \ar[r]^\mu \ar[d]_{L_G}  & X_{\mathbb K(Y)} \ar[dl]^{L_{\mathbb K(Y)}} \\
  Y &}
\end{xymatrix}
\end{equation*}
commutes. Set $\rho := \mu \circ \psi$.
\end{proof}

%In the language of \cite{BKM} which we have adopted above, Proposition \ref{18-11-25a} shows that the future cover is minimal with respect to regular and right-closing factor maps. The regularity condition on $\pi$ can not be removed: The right-resolving labeled graph $(G,L_G)$ in Figure \ref{16-10-25c} presents a sofic shift $Y$ whose future cover $(\mathbb K(Y),L_{\mathbb K(Y)})$ is obtained by removing the vertex 10 in Figure \ref{16-10-25c}, and $L_G$ does not factor through $L_{\mathbb K(Y)}$. The future cover is therefore not minimal with respect to arbitrary right-resolving factor maps in general.



\begin{lemma}\label{18-11-25b} Let $\pi : X \to Y$, $\pi': X' \to Y$ and $\rho : X \to X'$ be sliding block codes such that $\pi' \circ \rho = \pi$. Then $\rho(\mathcal R(\pi)) \subseteq \mathcal R(\pi')$.
\end{lemma}
\begin{proof} Let $x \in \mathcal R(\pi)$ and consider an element $y \in \mathbb U(\pi'(\rho(x))) = \mathbb U(\pi(x))$. There is then an element $z \in \mathbb U(x)$ such that $\pi(z) = y$. Since $\rho(z) \in \mathbb U(\rho(x))$ and $\pi'(\rho(z)) = y$, it follows that $\rho(x) \in \mathcal R(\pi')$.
\end{proof}


\begin{prop}\label{20-11-25ax} Let $\pi : X \to Y$ be a factor map with $X$ an SFT. Then $\pi$ is isomorphic to the future cover $L_{\mathbb K(Y)} : X_{\mathbb K(Y)} \to Y$ of $Y$ if and only if
\begin{itemize}
\item[a)] $\pi$ is regular,
\item[b)] $\pi$ is right-closing,
\item[c)] $\pi$ is injective on the set $\mathcal R(\pi)$ of regular points,
and
\item[d)] $\pi : X \to Y$ is minimal with respect to regular and right-closing factor maps.
\end{itemize}
\end{prop}
\begin{proof} $L_{\mathbb K(Y)}$ is regular by Lemma \ref{17-11-25h} and Lemma \ref{04-09-25h}, and right-closing because the future cover is right-resolving. By Proposition \ref{18-11-25a} it is minimal with respect to regular and right-closing factor maps. That $L_{\mathbb K(Y)}$ is injective on $\mathcal R(L_{\mathbb K(Y)})$ follows from Lemma \ref{02-09-25bx}.

Assume $\pi : X \to Y$ is a factor map with $X$ an SFT such that a), b), c) and d) hold. It follows that there are sliding block codes $\rho : X \to X_{\mathbb K(Y)}$ and $\rho' : X_{\mathbb K(Y)} \to X$ such that
\begin{equation*}
\begin{xymatrix}{
X_{\mathbb K(Y)}\ar[r]^-{\rho'} \ar[dr]_-{L_{\mathbb K(Y)}} &X \ar[r]^\rho \ar[d]_\pi  & X_{\mathbb K(Y)} \ar[dl]^{L_{\mathbb K(Y)}} \\
 &  Y &}
\end{xymatrix}
\end{equation*}
and
\begin{equation*}
\begin{xymatrix}{
X \ar[r]^-\rho \ar[dr]_-{\pi} & X_{\mathbb K(Y)}\ar[r]^{\rho'} \ar[d]_-{L_{\mathbb K(Y)}}  & X \ar[dl]^{\pi} \\
  & Y &}
\end{xymatrix}
\end{equation*}
both commute. It follows from Lemma \ref{18-11-25b} that $\rho$ and $\rho'$ both map regular points to regular points. Since $\pi$ and $L_{\mathbb K(Y)}$ are regular it follows that $\rho \circ \rho' = \id_{X_G}$ and $\rho' \circ \rho = \id_X$, first on regular points and then all over by continuity. 
\end{proof}




\subsection{The relation between the future cover and the follower set graph}\label{REG}

In this section we show how the future cover is related to another commonly used presentation of sofic shifts, \emph{the follower set graph}, cf. page 73 in \cite{LM}. 


Given a right-resolving labeled graph $(H,L_H)$, the set of its regular vertexes is a hereditary subset by Lemma \ref{31-08-25x} and we denote the corresponding hereditary labeled subgraph by $(H_{reg},L_H)$. As shown by the example depicted in Figure \ref{04-09-25gxx}, the sofic shift presented by $(H_{reg},L_H)$ can be smaller than that presented by $(H,L_H)$.


\begin{prop}\label{16-10-25e}
Let $Y$ be a sofic shift and $(G,\mathcal L)$ the follower set graph of $Y$. The future cover $(\mathbb K(Y),L_{\mathbb K(Y)})$ of $Y$ is isomorphic to $(G_{reg},\mathcal L)$ as a labeled graph.
\end{prop}
\begin{proof} The proof uses the observation that two words $w_1,w_2 \in \mathbb W(Y)$ have the same follower sets as defined in \cite{LM} if and only if $F(w_1) = F(w_2)$. Let $y \in Y$. By Lemma \ref{27e-09-25} there is an $N \in \mathbb N$ such that  $F(y) = F(y_{[-n,-1]})$ for all $n \geq N$. If $z \in Y$ and $F(y) = F(z)$ there is an $M \in \mathbb N$ such that $F(z_{[-n,-1]}) = F(y_{[-n,-1]})$ for all $n \geq M$. It follows that we can define a map $\chi: V_{\mathbb K(Y)} \to V_G$ such that $\chi(F(y)) = F(y_{[-n,-1]})$ for all large $n \in \mathbb N$. This map, $\chi$, is clearly injective.  For $y,z \in Y$ there is a labeled arrow $F(y) \overset{a}{\to} F(z)$ in $(\mathbb K(Y),L_{\mathbb K(Y)})$ if and only if $F(z) = \frac{F(y)}{a}$ if and only if $F(z_{[-n,-1]}) = \frac{F(y_{[-n,-1]})}{a}$ for all large $n$ if and only if there is a labeled arrow $\chi(F(y)) \overset{a}{\to} \chi(F(z))$ in $G$. This shows that $\chi$ is an embedding of labeled graphs. To show that $\chi(V_{\mathbb K(Y)})$ is a hereditary subset, let $w \in \mathbb W(Y)$ and assume that $\chi(F(y)) \overset{a}{\to} F(w)$ in $G$. Then $F(w) = \frac{F(y)}{a}$, and hence $F(w) = F(z)$, where $z \in Y$ is any element with $z_{(-\infty,-1]} = y_{(-\infty,-1]}a$. It follows that $F(w) = \chi(F(z))$. 

Since $(\mathbb K(Y),L_{\mathbb K(Y)})$ is regular by Lemma \ref{04-09-25h} and since the range of $\chi$ is a hereditary subgraph of $(G,\mathcal L)$ it follows from Lemma \ref{31-08-25x} that $\chi(V_{\mathbb K(Y)})$ consist only of regular vertexes in $G$. Conversely, if $w \in \mathbb W(Y)$ and $F(w)$ is a regular vertex in $G$, it follows from the defining relation \eqref{11-09-25} that $F(w) = F(z)$ for some $z \in Y$, proving that $F(w) = \chi(F(z)) \in \chi(V_{\mathbb K(Y)})$.




\end{proof}


\begin{ex} Consider the following labeled graph.

\begin{figure}[H]
\begin{equation*}
\begin{tikzpicture}[node distance={25mm}, thick, main/.style = {draw, circle}] 
\node[main] (1) {$a$}; 
\node[main] (2) [right of=1] {$b$}; 
\node[main] (3) [right of=2] {$c$};
\node[main] (4) [below of=3] {$d$};
\node[main] (5) [below of=2] {$e$};
%\node[main] (6) [below of=1] {$f$};
\draw[->] (1) to [out=120,in=70,looseness=8] node[above ,pos=0.4] {$0$} (1);
\draw[->] (1) to [out=0,in=180,looseness=1] node[above ,pos=0.4] {$2$} (2);
\draw[->] (2) to [out=120,in=70,looseness=8] node[above,pos=0.4] {$0$} (2);
\draw[->] (2) to [out=0,in=180,looseness=1] node[above ,pos=0.4] {$3$} (3);
\draw[->] (3) to [out=270,in=90,looseness=1] node[right ,pos=0.4] {$1$} (4);
\draw[->] (4) to [out=290,in=250,looseness=8] node[below ,pos=0.4] {$0$} (4);
\draw[->] (4) to [out=180,in=0,looseness=1] node[below ,pos=0.4] {$4$} (5);
\draw[->] (5) to [out=120,in=300,looseness=1] node[left ,pos=0.4] {$1$} (1);
\end{tikzpicture} 
\end{equation*}
\caption{The minimal right-resolving presentation of a mixing sofic shift whose follower set graph is larger than its future cover.}
\label{16-10-25}
\end{figure}


The follower sets of the vertexes $a,b,c,d,e$ are also follower sets of (synchronizing) words in $\mathbb W(Y)$. Besides these the follower sets of words from $\mathbb W(Y)$ include only the follower sets of the words $10$ and $0$, resulting in the following labeled graph depicting the follower set graph of $Y$. 


\begin{figure}[H]
\begin{equation*}
\begin{tikzpicture}[node distance={25mm}, thick, main/.style = {draw, circle}] 
\node[main] (1) {$10$};
\node[main] (2) [right of=1] {$a$}; 
\node[main] (3) [right of=2] {$b$}; 
\node[main] (4) [right of=3] {$c$};
\node[main] (5) [below of=4] {$d$};
\node[main] (6) [below of=3] {$e$};
\node[main] (7) [right of=5] {$0$};
\draw[->] (2) to [out=170,in=250,looseness=8] node[below ,pos=0.4] {$0$} (2);
\draw[->] (2) to [out=0,in=180,looseness=1] node[above ,pos=0.4] {$2$} (3);
\draw[->] (3) to [out=120,in=70,looseness=8] node[above,pos=0.4] {$0$} (3);
\draw[->] (3) to [out=0,in=180,looseness=1] node[above ,pos=0.4] {$3$} (4);
\draw[->] (4) to [out=270,in=90,looseness=1] node[right ,pos=0.4] {$1$} (5);
\draw[->] (5) to [out=70,in=350,looseness=8] node[right ,pos=0.4] {$0$} (5);
\draw[->] (5) to [out=180,in=0,looseness=1] node[above ,pos=0.4] {$4$} (6);
\draw[->] (6) to [out=120,in=300,looseness=1] node[left ,pos=0.4] {$1$} (2);
\draw[->] (1) to [out=40,in=150,looseness=1] node[above ,pos=0.4] {$2$} (3);
\draw[->] (1) to [out=270,in=180,looseness=1] node[below ,pos=0.4] {$4$} (6);
\draw[->] (7) to [out=100,in=45,looseness=2] node[above,pos=0.4] {$2$} (3);
\draw[->] (7) to [out=110,in=340,looseness=1] node[left,pos=0.4] {$3$} (4);
\draw[->] (7) to [out=210,in=340,looseness=1] node[below,pos=0.4] {$4$} (6);
\draw[->] (7) to [out=30,in=330,looseness=5] node[right,pos=0.4] {$0$} (7);
\draw[->] (1) to [out=180,in=110,looseness=4] node[left,pos=0.4] {$0$} (1);
\end{tikzpicture} 
\end{equation*}
\caption{The follower set graph of the sofic shift presented by Figure \ref{16-10-25}.}
\label{16-10-25c}
\end{figure}
The vertex $10$ is not regular while all the other are, and hence the future cover of $Y$ is the labeled graph obtained from the graph in Figure \ref{16-10-25c} by deleting the vertex $10$ and the arrows it emits.


\end{ex}


\section{An addendum to Krieger's theorem}


\subsection{Subset constructions }

In this section by a \emph{dynamical system} we mean a pair $(X,\phi)$ where $X$ is a compact metrizable space and $\phi : X \to X$ is a homeomorphism of $X$ onto itself. 

Let $(X,\phi)$ be a dynamical system and $d$ a metric for the topology of $X$. The set of non-empty compact subsets $K$ of $X$ can be equipped with a metric, the \emph{Hausdorff metric} $d_{Haus}$, defined by
$$
d_{Haus}(K,K') := \max \left\{ \sup_{x \in K}\left(\inf_{y \in K'}d(x,y)\right), \  \sup_{y \in K'}\left(\inf_{x \in K}d(x,y)\right)\right\}.
$$
This turns the set of non-empty compact subsets of $X$ into a compact metric space 
$$(X^{\#},d_{Haus}).
$$
See for example Section 2.6 in \cite{B}.

When $(X,d)$ and $(X',d')$ are two compact metric spaces and $\psi : X \to X'$ is a continuous map we get a continuous map $\psi^{\#} : X^{\#} \to X'^{\#}$ defined such that
$$
\psi^{\#}(K) := \psi(K) 
$$ 
for $K \in X^\#$. In particular, we get from $(X,\phi)$ a new dynamical system $(X^{\#},\phi^\#)$. More generally, when $(X,\phi)$ and $(Y,\phi')$ are dynamical systems and $\pi : X \to Y$ is a map of dynamical systems in the sense that $\pi \circ \phi = \phi' \circ \pi$, the set
\begin{equation*}\label{23-06-25b}
X^\pi := \left\{ K \in X^{\#}: \ \pi(x) = \pi(y)\ \forall x,y \in K \right\} 
\end{equation*}
is a closed $\phi^\#$-invariant subset of $X^\#$ and $(X^\pi,\phi^\#)$ is also a dynamical system. When $\pi$ is surjective and hence a factor map, then so is $\pi^\# : (X^\pi,\phi^\#) \to (Y,\phi')$.


%For each $n \in \mathbb N$ the set
%$$
%X^{(n)} := \left\{ K \in X^\# : \ \# K \leq n\right\}
%$$
%is closed and $\phi^\#$-invariant and hence so is %$$
%X^\pi \cap X^{(n)}.
%$$
%It follows that $(X^\pi \cap X^{(n)},\phi^\#)$ is a dynamical system and $\pi^\# :(X^\pi \cap X^{(n)},\phi^\#) \to (X,\phi)$ a factor map when $\pi$ is. 

The following is now a straightforward observation.
\begin{lemma}\label{25-06-25ax}
Consider two conjugacies $\psi' : (X,\phi) \to (X',\phi')$ and $\psi : (Y,\mu) \to (Y',\mu')$ of dynamical systems, and factor maps $\pi : (X,\phi) \to (Y,\mu)$ and $\pi' : (X',\phi') \to (Y',\mu')$ such that
\begin{equation*}\label{23-06-25x}
\xymatrix{
 X \ar[r]^-{\psi'} \ar[d]_\pi &  X' \ar[d]^{\pi'} \\
Y  \ar[r]^-\psi & Y'}
\end{equation*}
commutes. Then
$\psi'^\# : X^\pi \to {X'}^{\pi'}$
is a conjugacy, and % and the same is
%$$
%\psi'^\# : X^\pi \cap X^{(n)} \to X'^{\pi'}\cap %X'^{(n)}
%$$
%for each $n \in \mathbb N$. Furthermore, the diagram 
\begin{equation*}\label{25-06-25bx}
\xymatrix{
 X^\pi \ar[r]^-{\psi'^\#} \ar[d]_{\pi^\#} &  X'^{\pi'} \ar[d]^{\pi'^\#} \\
Y  \ar[r]^-\psi & Y'}
\end{equation*}
%and 
%\begin{equation}\label{25-06-25c}
%\xymatrix{
% X^\pi \cap X^{(n)}  \ar[r]^-{\psi'^\#} \ar[d]_{\pi^\#} &  X'^{\pi'}\cap X'^{(n)}  \ar[d]^{\pi'^\#} \\
%Y  \ar[r]^-\psi & Z}
%\end{equation}
commutes.
\end{lemma}





\subsection{Subset constructions for sofic shifts}\label{subset}

Let $(H,L_H)$ be a labeled graph and $Y:= L_H(X_H)$ the presented sofic shift. We define a labeled graph $(H'', L_{H''})$ as follows. The set $V_{H''}$ of vertexes consists of the non-empty subsets $F$ of vertexes from $V_H$. For $F,F' \in  V_{H''}$ there is an edge $e' \in E_{H''}$ with $s_{H''} (e') = F$ and $t_{H''} (e') = F'$ when there is a symbol $a \in A$, the alphabet of $Y$, such that
\begin{equation*}\label{23-06-25ay}
F \subseteq \left\{ s_H(e): \ e \in E_H, \ L_H(e) = a \right\} 
\end{equation*}
and
$$
F' = \left\{ t_H(e) : \ e \in E_H, \ s_H(e) \in F, \ L_H(e) = a \right\} .
$$ 
We set then $L_{H''} (e') = a$. While the graph $H''$ in principle contains all subsets of $V_H$ as vertexes, it is clear from the definition of arrows in $E_{H''}$ that typically only a much smaller collection of such subsets will occur in elements of $X_{H''}$. In particular, when we consider $H''$ in the following it will be assumed that $H''$ has been trimmed by repeatedly removing sources and sinks to obtain a graph that does not contain neither.




\begin{lemma}\label{23-06-25dyx} Assume that $(H,L_H)$ is right-resolving. There is a conjugacy $\kappa : (X_H^{L_H}, \sigma^\#) \to (X_{H''} , \sigma)$ such that
\begin{equation*}\label{23-06-25cy}
\xymatrix{
 X^{L_H}_H \ar[rr]^-\kappa \ar[dr]_-{L_H^\#} & &  X_{H''}  \ar[dl]^-{L_{H''} } \\
 & Y & }
\end{equation*}
commutes.
\end{lemma}
\begin{proof} Let $ K \in X^{L_H}_H$. Note that $K$ is a finite set since $(H,L_H)$ is right-resolving. Define 
$$
F_i := \left\{ s_H(k_i) : \ k \in K \right\}
$$
for $i \in \mathbb Z$. For each $i \in \mathbb Z$ there is a unique edge $e_i(K)$ in $E_{H''} $ from $F_i$ to $F_{i+1}$ with label $L_{H''} (e_i(K)) = L_H(k_i)$ for all $k \in K$, and we set
$$
\kappa(K) := \left(e_i(K)\right)_{i \in \mathbb Z} \in X_{H''}  .
$$ 
Since $e_i(\sigma(K)) = e_{i+1}(K)$ we find that 
$$
\kappa \circ \sigma^\#(K) = \kappa(\sigma(K)) = \left(e_i(\sigma(K))\right)_{i \in \mathbb Z} = \left(e_{i+1}(K)\right)_{i \in \mathbb Z} = \left(\sigma(\kappa(K))_i\right)_{i \in \mathbb Z},
$$
proving that $\kappa\circ \sigma^\# = \sigma \circ \kappa$. Note also that $L_{H''}  (\kappa(K)) = L_H^\#(K)$.


To construct an inverse map $\kappa^{-1} : X_{H''}  \to X_H^{L_H}$ to $\kappa$, let $x = (x_i)_{i \in \mathbb Z} \in X_{H''} $. Since $(H,L_H)$ is right-resolving, $\# s_{H''} (x_i) \geq  \# s_{H''} (x_{i+1})$, and we choose $N \in \mathbb N$ such that $\# s_{H''} (x_i) = \# s_{H''} (x_N)$ for all $i \leq N$. For each $v \in  s_{H''} (x_N)$ there is a unique ray $\gamma^v \in X_H$ such that $s_H(\gamma^v_N) = v$, $s_H(\gamma^v_i) \in s_{H''}(x_i)$ for all $i \in \mathbb Z$ and $L_H(\gamma^v) = L_{H''} (x)$. Setting
$$
\kappa^{-1}(x) := \left\{\gamma^v : \ v \in  s_{H''} (x_N) \right\} 
$$
we obtain a well-defined map $\kappa^{-1}: X_{H''}  \to X^{L_H}_H$. It is straightforward to check that $\kappa^{-1}$ is the inverse of $\kappa$.
%\begin{color}{red}
%\begin{align*}
%& \psi^{-1}(\psi(K)) = \psi^{-1}\left( \left(e_i(K)\right)_{i \in \mathbb Z}\right) \\
%& = \left\{ \gamma_v : \ v \in s_{H''} (e_N(K)) \right\} = \left\{ \gamma_v : \ v \in F_N \right\} \\
%&  = \left\{ \gamma_v : \ v \in \left\{s_G(k_N) : \ k \in K\right\}  \right\} = K
%\end{align*}
%and since $F_i = \left\{s_H((\gamma_v)_i) : \  v \in  s_{H''} (x_N) \right\} = s_{H''} (x_i)$,
%\begin{align*}
%& \psi(\psi^{-1}(x)) = \psi\left( \left\{\gamma_v : \ v \in  s_{H''} (x_N) \right\}\right) \\
%& = (e_i( \left\{\gamma_v : \ v \in  s_{H''} (x_N) \right\}))_{i \in \mathbb Z}
% = (x_i)_{i \in \mathbb Z} =x.\\
%\end{align*}
%\end{color}


Since $X_H^{L_H}$ and $X_{H''} $ are compact Hausdorff spaces it suffices now to show that $\kappa$ is continuous. For this we use the metric $d$ on $X_H$ defined such that $d(x,y) = \frac{1}{n+1}$, where $n = \infty$ when $x = y$ and otherwise $n = \min \left\{ j\in \mathbb N \cup \{0\} : \ x_{[-j,j]} \neq y_{[-j,j]}\right\}$. Consider then two elements $K,K' \in X_H^{L_H}$. If $d_{Haus}(K,K') < \frac{1}{N}$, it follows that 
$$
\left\{k_i : \ k \in K\right\} =  \left\{k_i : \ k \in K'\right\}  
$$
for $-N \leq i \leq N$, and hence also that $\kappa(K)_i = e_i(K) = e_i(K') = \kappa(K')_i$ for $-N \leq i \leq N$. This shows that $\kappa$ is continuous.
\end{proof}


\begin{corollary}\label{17-11-25} Let $Y$ and $Z$ be sofic shifts and $(G,L_G)$, $(H,L_H)$ right-resolving presentations of $Y$ and $Y'$, respectively. Assume that there are conjugacies $\psi : Y \to Z$ and $\phi : X_G \to X_H$ such that 
\begin{equation*}\label{17-11-25b}
\xymatrix{
 X_G \ar[r]^-{\phi} \ar[d]_{L_G}&  X_H \ar[d]^{L_H} \\
 Y  \ar[r]^-\psi & Z}
\end{equation*}
commutes. There is a conjugacy $\phi'' : X_{G''} \to X_{H''}$ such that
\begin{equation*}\label{17-11-25c}
\xymatrix{
 X_{G''} \ar[r]^-{\phi''} \ar[d]_{L_{G''}}&  X_{H''} \ar[d]^{L_{H''}} \\
Y   \ar[r]^-\psi & Z}
\end{equation*}
commutes.
\end{corollary}
\begin{proof} Combine Lemma \ref{25-06-25ax} and Lemma \ref{23-06-25dyx}.
\end{proof}


Thus $({G''},L_{G''})$ is a weakly canonical cover when $(G,L_G)$ is. However, the conjugacy $\phi''$ may not be unique and there is no analogue to the right-inverse for the labeling map as in Krieger's theorem.  We continue therefore to identify a labeled subgraph of $(G'',L_{G''})$ which does give rise to a strongly canonical cover.

\bigskip
 

Since $(H,L_H)$ is right-resolving, $L_H^{-1}(y)$ is a finite set and hence an element of $X_H^{L_H}$ for each $y \in Y$. We define $\beta_H : Y \to X_{H''} $ such that
$$
\beta_H(y) := \kappa(L_H^{-1}(y)),
$$
where $\kappa : X_H^{L_H} \to X_{H''}$ is the conjugacy from Lemma \ref{23-06-25dyx}. Note that
\begin{equation}\label{22-11-25}
L_{H''} \circ \beta_H = \id_Y .
\end{equation}
Like the map $\alpha_Y$ considered in Section \ref{rays} the map $\beta_H$ commutes with the shift but is not continuous.





To proceed we need some terminology. Let $X$ be a subshift. We say that two elements $x,y \in X$ are \emph{backward asymptotic} and that \emph{$x$ is backward asymptotic to $y$} when there is an $N \in \mathbb Z$ such that $x_i=y_i$ for all $i \leq N$. When $A \subseteq X$ is a subset and $x \in X$ we say that $x$ is \emph{backward asymptotic to $A$} when $x$ is backward asymptotic to an element of $A$. When $A$ is a subshift of $X$ we say that $x \in X$ is \emph{backward transitive} to $A$ when $x$ is backward asymptotic to $A$ and every word $w \in \mathbb W(A)$ appears infinitely often to the left in $x$, in the sense that for all $N\in \mathbb Z$ there are $i \leq j \leq N$ such that $w =x_{[i,j]}$. 

Concerning the components of $H''$ we will say that a component $C'$ in $H''$ \emph{connects} to another component $C$ when there is a path $\mu$ in $H''$ such that $s_{H''}(\mu) \in C'$ and $t_{H''}(\mu) \in C$. Note that for any component $C$ in $H''$ the cardinality $\# F$ of the elements $F$ of $C \subseteq 2^{V_H}$ is the same; this number is the \emph{multiplicity} of $C$ and we denote it by $M(C)$.
 
\begin{lemma}\label{28-12-25d} Let $C$ be a component in $H''$ such that $\beta_H(y)$ is backward asymptotic to $X_C$ for some $y \in Y$. Then $X_C \subseteq \overline{\beta_H(Y)}$.
\end{lemma}
\begin{proof} Let $y \in Y$ and let $C$ be the component in $H''$ for which there is a $N \in \mathbb Z$ such that $\beta_H(y)_j \in C$ for all $j \leq N$. It follows that there is a $ R < N$ such that if $z \in X_{H''}$ and $L_{H''}(z)_{[R,\infty)} = y_{[R,\infty)}$, then $s_{H''}(z_N) \subseteq s_{H''}(\beta_H(y)_N)$. (If not, a compactness argument would produce an element $z' \in X_{H''}$ such that $L_{H''}(z') = y$ and $\# s_{H''}(z'_N) > \# s_{H''}(\beta_H(y)_N)$. A contradiction.) Let $C'$ be a component which connects to $C$ and has the largest multiplicity among the components that connect to $C$. Let $\gamma$ be a finite path in $H''_C$. We can then choose an element $z'  \in X_{H''}$ with the following properties:
\begin{itemize}
\item[i)] $z'$ is backward asymptotic to $C'$,
\item[ii)] there is a $R' < R$ such that $z'_{[R',R]}$ is a path in $H''_C$ with $\gamma \subseteq z'_{[R',R]}$, and
\item[iii)] $z'_{[R,\infty)} = \beta_H(y)_{[R,\infty)}$.
\end{itemize}
We claim that 
\begin{equation}\label{foerstejuledag1}
z' = \beta_H(L_{H''}(z')).
\end{equation}
To see this, we use that  
\begin{equation}\label{foerstejuledag2} 
\beta_H(L_{H''}(z'))_N =  \beta_H(y)_N.
 \end{equation}
Indeed, $L_{H''}(\beta_H(L_{H''}(z')))_{[R,\infty)} = L_{H''}(z')_{[R,\infty)}= y_{[R,\infty)}$ by iii), which implies that 
 $$
 s_{H''}(\beta_H(L_{H''}(z'))_N) \subseteq  s_{H''}(\beta_H(y)_N)
 $$ 
 by the choice of $R$, while the opposite inclusion follows from 
 $$
 s_{H''}(\beta_H(y)_N) = s_{H''}(z'_N) \subseteq s_{H''}(\beta_H(L_{H''}(z'))_N).
 $$ 


 Let $C''$ be the component in $H''$ with the property that $\beta_H(L_{H''}(z'))$ is backward asymptotic to $X_{C''}$. Since 
 \begin{equation}\label{29-12-25}
 s_{H''}(z'_j) \subseteq s_{H''}(\beta_H(L_{H''}(z'))_j)
 \end{equation} 
 for all $j$ and $z'$ is backward asymptotic to $X_{C'}$ while $\beta_H(L_{H''}(z'))$ is backward asymptotic to $X_{C''}$, we have that $M(C')  \leq M(C'')$. But it follows from \eqref{foerstejuledag2} that $C''$ connects to $C$ and hence $M(C'') \leq M(C')$ by definition of $C'$, and we find therefore that $M(C'') = M(C')$. It follows from this identity and \eqref{29-12-25} that $s_{H''}(z'_j) = s_{H''}(\beta_H(L_{H''}(z'))_j)$ for all $j$ 'close to $-\infty$'; a conclusion which implies the claim; \eqref{foerstejuledag1}. The inclusion $X_C \subseteq \overline{\beta_H(Y)}$ follows now from \eqref{foerstejuledag1} because the path $\gamma$ is arbitrary and $\beta_H(Y)$ is $\sigma$-invariant.
\end{proof}


\begin{definition}\label{30-10-25} We denote by $H'$ the hereditary subgraph of $H''$ generated by the components $C$ of $H''$ for which there is a point $y \in Y$ such that $\beta_H(y)$ is backward asymptotic to $X_C$. 
\end{definition}

To simplify notation, set $L_{H'} := L_{H''}|_{H'}$. Since $\beta_H(Y) \subseteq X_{G'}$ by definition of $G'$, it follows from \eqref{22-11-25} that $L_{X_{G'}}(X_{G'}) = Y$, implying that $(G',L_{G'})$ is a presentation of $Y$.





\begin{thm}\label{31-10-25hx} Let $(G,L_G)$ and $(H,L_H)$ be right-resolving presentations of the sofic shifts $Y$ and $Z$, respectively. Assume that there are conjugacies $\psi : Y \to Z$ and $\phi : X_G \to X_H$ such that 
\begin{equation*}\label{30-10-25bx}
\xymatrix{
 X_G \ar[r]^-{\phi} \ar[d]_{L_G}&  X_H \ar[d]^{L_H} \\
 Y  \ar[r]^-\psi & Z}
\end{equation*}
commutes. There is a unique conjugacy $\phi' : X_{G'} \to X_{H'}$ such that
\begin{equation*}\label{30-10-25cx}
\xymatrix{
 X_{G'} \ar[r]^-{\phi'} \ar[d]_{L_{G'}}&  X_{H'} \ar[d]^{L_{H'}} \\
Y   \ar[r]^-\psi & Z}
\end{equation*}
commutes.
\end{thm}
\begin{proof} It follows from Corollary \ref{17-11-25} that there is a conjugacy $\phi'' : X_{G''} \to X_{H''}$ such that
\begin{equation*}\label{22-11-25a}
\xymatrix{
 X_{G''} \ar[r]^-{\phi''} \ar[d]_{L_{G''}}&  X_{H''} \ar[d]^{L_{H''}} \\
 Y  \ar[r]^-\psi & Z}
\end{equation*}
commutes. When we use the same symbol $\kappa$ for the two conjugacies $\kappa : X_G^{L_G} \to X_{G''}$ and $\kappa : X_H^{L_H} \to X_{H''}$ from Lemma \ref{23-06-25dyx}, we note that $\phi'' = \kappa \circ \phi^\# \circ \kappa^{-1}$. Since $\phi^\#(L_G^{-1}(y)) = L_H^{-1}(\psi(y))$ it follows that \begin{equation}\label{28-12-25c}
\phi''(\beta_G(y)) = \beta_H(\psi(y))
\end{equation}
 for all $y \in Y$. Let $\mathcal C_G$ denote the set of components $C$ of $G''$ with the property that there is a point $y \in Y$ such that $\beta_G(y)$ is backward asymptotic to $X_C$. We denote the analogous set of components in $H''$ by $\mathcal C_H$. It follows from \eqref{28-12-25c} that when $C \in \mathcal C_G$ we have that $\phi''(X_C) = X_{C'}$ for some $C' \in \mathcal C_H$. If $x \in X_{G''}$ is backward asymptotic to $X_C$ we have that $\phi''(x)$ is backward asymptotic to $X_{C'}$. Since the elements of $X_{G''}$ that are backward asymptotic to $X_C$ for some $C \in \mathcal C_G$ constitute a dense set in $X_{G'}$, it follows that $\phi''(X_{G'}) \subseteq X_{H'}$. By symmetry we have equality; $\phi''(X_{G'}) =X_{H'}$, and we set $\phi':= \phi''|_{X_{G'}}$. This proves the existence part. 
 To prove uniqueness, consider a conjugacy $\lambda : X_{G'} \to X_{H'}$ such that
\begin{equation*}\label{30-10-25cy}
\xymatrix{
 X_{G'} \ar[r]^-{\lambda} \ar[d]_{L_{G'}}&  X_{H'} \ar[d]^{L_{H'}} \\
Y   \ar[r]^-\psi & Z}
\end{equation*}
commutes. Set 
$$
\lambda' := \kappa^{-1}\circ \lambda \circ \kappa|_{\kappa^{-1}(X_{G'})}.
$$ 
By Lemma \ref{23-06-25dyx} also
\begin{equation}\label{28-12-25e}
\xymatrix{
 X_{G}^{L_G}\cap \kappa^{-1}(X_{G'}) \ar[r]^-{\lambda'} \ar[d]_{L_{G}^{\#}}&  X_{H}^{L_H}\cap \kappa^{-1}(X_{H'}) \ar[d]^{L_{H}^{\#}} \\
Y   \ar[r]^-\psi & Z}
\end{equation}
commutes. Note that $L_G^{-1}(y)  \in  X_{G}^{L_G}\cap \kappa^{-1}(X_{G'})$ for all $y \in Y$ since $\beta_G(Y) \subseteq X_{G'}$. It follows from \eqref{28-12-25e} that $\lambda'(L_G^{-1}(y)) \subseteq L_H^{-1}(\psi(y))$ and by symmetry that $\lambda'(L_G^{-1}(y)) =L_H^{-1}(\psi(y))$ for all $y \in Y$. Hence $\lambda(\beta_G(y)) = \beta_H(\psi(y))$ for all $y \in Y$. In particular, $\lambda(\beta_G(y)) = \phi'(\beta_G(y))$ for all $y \in Y$ and hence, by Lemma \ref{28-12-25d}, $\lambda|_{X_C} = \phi'|_{X_C}$ for $C \in \mathcal C_G$. Consider then an $x \in X_{G'}$ backward asymptotic to $X_C, \ C \in \mathcal C_G$. Since $\lambda$ and $\phi'$ agree on $X_C$ it follows that there is an $N \in \mathbb Z$ such that $\lambda(x)_{(-\infty,N]} = \phi'(x)_{(-\infty,N]}$. Since $L_{H'}$ is right-resolving and $L_{H'}(\lambda(x)) = \psi(L_{G'}(x)) = L_{H'}(\phi'(x))$, we find that $\lambda(x) = \phi'(x)$. Since such $x$'s are dense in $X_{G'}$ we conclude that $\lambda = \phi'$.
 \end{proof}


We note that in analogy with \eqref{24-11-24x} in Krieger's theorem, the diagram
\begin{equation}\label{30-10-25d}
\xymatrix{
 X_{G'} \ar[r]^-{\phi'} &  X_{H'}  \\
Y \ar[u]^-{\beta_G}  \ar[r]^-\psi & Z \ar[u]_-{\beta_H}}
\end{equation}
commutes.


For a sofic shift $Y$, set
$$
(\mathbb K'(Y),L_{\mathbb K'(Y)})  := (\mathbb K(Y)', L_{\mathbb K(Y)'}).
$$
By Theorem \ref{Krieger} and Theorem \ref{31-10-25hx} this defines a strongly canonical cover for sofic shifts:


\begin{thm}\label{20-01-26m} Let $\psi : Y \to Z$ be a conjugacy of sofic shifts. There is a unique conjugacy $\psi_{\mathbb K'} : X_{\mathbb K'(Y)} \to X_{\mathbb K'(Z)}$ such that
\begin{equation*}
\xymatrix{
 X_{\mathbb K'(Y)}  \ar[r]^-{\psi_{\mathbb K'}} \ar[d]_{L_{\mathbb K'(Y)}}&  X_{\mathbb K'(Z)} \ar[d]^{L_{\mathbb K'(Z)}} \\
Y   \ar[r]^-\psi & Z}
\end{equation*}
commutes.
\end{thm}

For an irreducible sofic shift $Y$ we set
$$
(\mathbb F'(Y),L_{\mathbb F'(Y)})  := (\mathbb F(Y)', L_{\mathbb F(Y)'}).
$$
By Corollary \ref{30-09-25c} and Theorem \ref{31-10-25hx} this defines also a strongly canonical cover for irreducible sofic shifts:


\begin{thm}\label{20-01-26n} Let $\psi : Y \to Z$ be a conjugacy of irreducible sofic shifts. There is a unique conjugacy $\psi_{\mathbb F'} : X_{\mathbb F'(Y)} \to X_{\mathbb F'(Z)}$ such that
\begin{equation*}
\xymatrix{
 X_{\mathbb F'(Y)}  \ar[r]^-{\psi_{\mathbb F'}} \ar[d]_{L_{\mathbb F'(Y)}}&  X_{\mathbb F'(Z)} \ar[d]^{L_{\mathbb F'(Z)}} \\
Y   \ar[r]^-\psi & Z}
\end{equation*}
commutes.
\end{thm}




\begin{ex}\label{evenshift}
It seems to be a well-established tradition that examples of sofic shifts start with the even shift, and we have no reason to deviate from this. The minimal right-resolving presentation $(\mathbb F(E),L_{\mathbb F(E)})$ of the even shift $E$ is shown in Figure \ref{03-10-25e} and the future cover $(\mathbb K(E),L_{\mathbb K(E)})$ in Figure \ref{11-09-25cx}. Then $(\mathbb F '(E),L_{\mathbb F'(E)})$ is the following labeled graph,

\begin{equation*}
\begin{tikzpicture}[node distance={25mm}, thick, main/.style = {draw, circle}] 
\node[main] (1) {$a$}; 
\node[main] (2) [right of=1] {$b$}; 
\node[main] (3) [below of=1] {$\{a,b\}$};
\draw[->] (1) to [out=80,in=110,looseness=1.0] node[above ,pos=0.5] {$0$} (2); 
\draw[->] (2) to [out=260,in=310,looseness=1.0]  node[below ,pos=0.5] {$0$}(1); 
\draw[->] (1) to [out=160,in=220,looseness=15] node[above ,pos=0.4] {$1$} (1);
\draw[->] (3) to [out=220,in=300,looseness=5] node[below ,pos=0.5] {$0$} (3);
\end{tikzpicture} 
\end{equation*}

and $(\mathbb K'(E), L_{\mathbb K'(E)})$ is the labeled graph 

\begin{equation*}
\begin{tikzpicture}[node distance={25mm}, thick, main/.style = {draw, circle}] 
\node[main] (1) {$\{a,b,c\}$}; 
\node[main] (2) [below of=1] {$\{a,c\}$}; 
\node[main] (3) [right of=2] {$\{a\}$};
\node[main] (4) [right of=3] {$\{b\}$};
\node[main] (5) [below of=2] {$\{b,c\}$};
\draw[->] (1) to [out=110,in=70,looseness=7] node[above ,pos=0.5] {$0$} (1);
\draw[->] (2) to [out=250,in=110,looseness=1] node[left, pos=0.5] {$0$} (5);
\draw[->] (5) to [out=80,in=280,looseness=1] node[right, pos=0.5] {$0$} (2);
\draw[->] (2) to [out=10,in=170,looseness=1] node[above, pos=0.5] {$1$} (3);
\draw[->] (3) to [out=240,in=290,looseness=7] node[below, pos=0.5] {$1$} (3);
\draw[->] (3) to [out=20,in=170,looseness=1] node[above, pos=0.5] {$0$} (4);
\draw[->] (4) to [out=200,in=350,looseness=1] node[below, pos=0.5] {$0$} (3);
\end{tikzpicture} 
\end{equation*}

We see therefore, already from the even shift, that in general the canonical covers $(\mathbb F(E),L_{\mathbb F(E)})$, $(\mathbb K(E),L_{\mathbb K(E)})$, $(\mathbb F'(E),L_{\mathbb F'(E)})$ and $(\mathbb K'(E),L_{\mathbb K'(E)})$ are all different.
\end{ex}



%In order to compare the strongly canonical covers $(\mathbb K(Y),L_{\mathbb K(Y)}) $, $(\mathbb K'(Y),L_{\mathbb K'(Y)})$ and $(\mathbb F'(Y),L_{\mathbb F'(Y)})$ we need a general procedure to construct the future cover.

\begin{thebibliography}{WWWWW} %antallet af W'er er bredeste ind
%\bibitem[M]{M} S. MacDonald, {\em Gibbs measure and factor codes in symbolic dynamics}, Dissertation, January 2023.

%\bibitem[B]{B} M. Boyle, {\em Lower entropy factors of sofic systems}, Ergod. Th. Dynam. Syst. {\bf 4} (1984), 541-557.

%\bibitem[BCE]{BCE} M. Boyle, T.M. Carlsen, S. Eilers, {\em Flow equivalence of sofic shifts}, Israel J. Math. {\bf 225} (2018), 111-146.

\bibitem[B]{B} M. Barnsley, {\em Fractals everywhere}, Academic Press, 1988.

\bibitem[BKM]{BKM} M. Boyle, B. Kitchens and B.H. Marcus, {\em A note on minimal covers for sofic systems}, Proc. Amer. Math. Soc. {\bf 95} (1985), 403-411.


\bibitem[F]{F} R. Fischer, {\em Sofic systems and graphs}, Monatsh. Math. {\bf 80} (1975), 179--186.


\bibitem[J]{J} R. Johansen, {\em On the structure of covers of sofic shifts}, Documenta Mathematica {\bf 16} (2011), 111-131. 

\bibitem[Jo]{Jo} N. Jonoska, {\em Sofic shifts with synchronizing presentations}, Theoret. Comput. Sci. {\bf 158} (1996), 81--115.






\bibitem[Kr1]{Kr1} W. Krieger, {\em On sofic systems I}, Israel J. Math. {\bf 48} (1984), 305-330.

\bibitem[Kr2]{Kr2} W. Krieger, {\em On sofic systems II}, Israel J. Math. {\bf 60} (1987), 167-176.


\bibitem[LM]{LM} D. Lind and B. Marcus, {\em An introduction to symbolic dynamics and coding}, Cambridge University Press, 1995.




\bibitem[MMTW]{MMTW} B. Marcus, T. Meyerovitch, K. Thomsen and C. Wu, {\em
Factorizable embeddings and the period of an irreducible sofic shift}, arXiv:2508.02554 


\bibitem[N]{N} M. Nasu, {\em Topological conjugacy for sofic systems}, Ergod. Th. \& Dynam. Sys. {\bf 6} (1986), 265-280.


%\bibitem[T]{T} K. Thomsen, {\em On the structure of a sofic shift space}, Trans. Amer. Math. Soc. {\bf 356} (2004), 3557-3619.

\end{thebibliography}


\end{document}


\section{Constructing the future cover}

In this section we present a method to construct the future cover of a sofic shift from any given presentation of it. 


\subsection{On the merging of a labeled graph}\label{merged}

Let $(H,L_H)$ be a labeled graph. The \emph{merged graph} $([H],L_{[H]})$ of $(H,L_H)$ is then defined as follows, cf. page 78 of \cite{LM}. The set of vertexes $V_{[H]}$ is the set
$$
V_{[H]} := \left\{f_H(v): \ v \in V_H\right\}
$$
and for $a \in A$ there is a labeled arrow $f_H(v) \overset{a}{\to} f_H(w)$ when there is are vertexes $v',w' \in V_H$ such that $f_H(v') = f_H(v), \ f_H(w') = f_H(w)$ and there is a labeled arrow $v' \overset{a}{\to} w'$ in $H$. When $(H,L_H)$ is right-resolving we can define a map $f_H: E_H \to E_{[H]}$ such that $f_H(e)$ is the arrow $f_H(e) \in E_{[H]}$ with the properties
\begin{itemize}
\item $s_{[H]}(f_H(e)) = f_H\left(s_H(e)\right)$, and
\item $L_{[H]}(e) = L_H(e)$.
\end{itemize}
Then $f_H:H \to [H]$ is a labeled-graphs homomorphism. 


 
 \begin{lemma}\label{03-11-24x} Assume that $(H,L_H)$ is right-resolving. Let $v,w \in V_{H}$. There is a labeled arrow $f_H(v) \overset{a}{\to} f_H(w)$ in $[H]$ if and only if there is a labeled arrow $v \overset{a}{\to} u$ in $H$ such that $f_H(u) =f_H(w)$. 
\end{lemma} 
\begin{proof} Assume that $f_H(v) \overset{a}{\to} f_H(w)$ in $[H]$ and let $v_1, w_1 \in V_{H}$ be such that $v_1 \overset{a}{\to} w_1$ in $H$, $f_H(v_1) = f_H(v)$ and $f_H(w_1) = f_H(w)$. Since $f_H(v) = f_H(v_1)$ there is an arrow $e$ in $E_{H}$ labeled $a$ which leaves $v$. Set $u:= t_{H}(e)$. By using that $(H,L_{H})$ is right-resolving we find then that
$$
f_H(u) = \frac{f_H(v)}{a} = \frac{f_H(v_1)}{a} = f_H(w_1) .
$$
This establishes the 'only if' part. The 'if' part is trivial.
\end{proof}



\begin{lemma}\label{10-09-25} Assume that $(H,L_H)$ is right-resolving. Then $([H],L_{[H]})$ is right-resolving and follower-separated, and there is a labeled-graph homomorphism $f_H: H \to [H]$ such that $f_H: X_H \to X_{[H]}$ is a factor map. If $(H,L_H)$ is regular, then so is $([H],L_{[H]})$.
\end{lemma}
\begin{proof} $([H],L_{[H]})$ is right-resolving and follower-separated by Lemma 3.3.8 of \cite{LM}. To see that $f_H: X_H \to X_{[H]}$ is surjective, consider a finite path 
\begin{equation}\label{03-09-25d}
f_H(v_1) \overset{a_1}{\to} f_H(v_2) \overset{a_2}{\to} f_H(v_3) \overset{a_3}{\to} \cdots \cdots \overset{a_{n-1}}{\to} f_H(v_n)
\end{equation}
in $[H]$. Repeated applications of Lemma \ref{03-11-24x} give a path 
$$
v_1 \overset{a_1}{\to} u_2 \overset{a_2}{\to}  u_3 \overset{a_3}{\to} \cdots \cdots \overset{a_{n-1}}{\to} u_n
$$ 
in $H$ which is mapped to \eqref{03-09-25d} under $f_H$. The surjectivity of $f_H$ follows therefore from compactness of $ X_{H}$. 

Assume now that $(H,L_H)$ is regular and consider a vertex $v \in V_H$. Since $(H,L_H)$ is regular there is a ray $z' \in X_H$ such that $t_H(z'_{(-\infty,-1]}) = v$ and $f_H(v) = F(L_H(z'))$. Set $z:= f_H(z') \in X_{[H]}$ and note that $t_{[H]}(z_{(-\infty,-1]}) = f_H( t_H(z'_{(-\infty,-1]})) = f_H(v)$. An application of Lemma \ref{03-11-24x} as above shows that $f_{[H]}(f_H(v)) = f_H(v)$. Since $L_{[H]}(z) = L_H(z')$ we conclude first that $f_{[H]}(f_H(v)) = F(L_{[H]}(z))$ and then that $([H],L_{[H]})$ is regular.
\end{proof}


We note that the morphism $f_H: H \to [H]$ is injective on $V_H$ if and only if it is an isomorphism of labeled graphs, and this happens if and only if $(H,L_H)$ is follower-separated.



\subsection{A cover of the future cover}\label{consec}


Let $(G,L_G)$ be a labeled graph and $Y =L_G(X_G)$ the sofic shift presented by $(G,L_G)$. When $D\subseteq V_{G}$ and $a \in A$, set
$$
[D,a]: = \left\{ t_G(e): \ e \in E_G, \ s_G(e) \in D, \ L_G(e) = a \right\};
$$
the set of vertexes that can be reached from an element of $D$ using an edge in $G$ labeled by $a$. We consider the collection $2^{V_G}$ of non-empty subsets of $V_G$ as the set of vertexes in a graph $\overline{G}$ with edges 
$$
E_{\overline{G}} := \left\{(D,a) \in 2^{V_{G}} \times A : \ [D,a] \neq \emptyset \right\} .
$$
The start vertex $s_{\overline{G}}(D,a)$ of $(D,a) \in E_{\overline{G}}$ is $D \in 2^{V_G}$ and the terminal vertex $t_{\overline{G}}(D,a)$ is $[D,a]$. We define a labeling $L_{\overline{G}}$ of $\overline{G}$ such that
$$
L_{\overline{G}}(D,a) := a 
$$
when $(D,a) \in E_{\overline{G}}$. Then $(\overline{G},L_{\overline{G}})$ is a right-resolving labeled graph and 
$L_{\overline{G}}(X_{\overline{G}}) =  Y$.




For $y \in Y$, set
$$
D^y:= \left\{ t_G(x) : \ x \in X_G(-\infty,-1], \ L_G(x) = y_{(-\infty,-1]} \right\} .
$$
 
The set of vertexes $V_{\underline{G}} := \left\{D^y: \ y \in Y\right\}$ is a hereditary subset of vertexes in $V_{\overline{G}}$; that is, if $e \in E_{\overline{G}}$ and $s_{\overline{G}}(e) \in V_{\underline{G}}$, then $t_{\overline{G}}(e) \in V_{\underline{G}}$. Hence, when we set $E_{\underline{G}} := \left\{e \in E_{\overline{G}}: \ s_{\overline{G}}(e) \in V_{\underline{G}}\right\}$, we have defined a graph $\underline{G}$, and hence also a labeled graph $(\underline{G}, L_{\underline{G}})$, where $L_{\underline{G}}:= L_{\overline{G}}|_{E_{\underline{G}}}$. 

Note that we can define a shift-commuting map $\alpha_G : Y \to X_{\underline{G}}$ such that 
\begin{itemize}
\item[a)] $s_{\underline{G}}(\alpha_G(y)_k) = D^{\sigma^k(y)}$ for all $k \in \mathbb Z$, and
\item[b)] $L_{\underline{G}}\left(\alpha_G(y)\right) = y$.
\end{itemize}
In particular, $L_{\underline{G}}(X_{\underline{G}}) = Y$ and $(\underline{G}, L_{\underline{G}})$ is thus a presentation of $Y$. We show next that the future cover of $Y$ is the merged graph of $(\underline{G}, L_{\underline{G}})$.


Let $([\underline{G}],L_{[\underline{G}]})$ be the merged graph of $(\underline{G}, L_{\underline{G}})$. Thus the set of vertexes $V_{[\underline{G}]}$ consists of the set 
$$
V_{[\underline{G}]} : = \left\{ f_{\underline{G}}(v) : \ v \in V_{\underline{G}} \right\}.
$$
%where
%$$
%f_{\underline{G}}(v) := \left\{L_{\underline{G}}(x) : \ x \in X_{\underline{G}}[0,\infty), \ s_{\underline{G}}(x) = v\right\}.
%$$
By Lemma \ref{10-09-25} $([\underline{G}], L_{[\underline{G}]})$ is right-resolving, follower-separated and the morphism $f_{\underline{G}} : \underline{G} \to [\underline{G}]$ of labeled graphs induces a factor map $f_{\underline{G}} : X_{\underline{G}} \to X_{[\underline{G}]}$.











\begin{lemma}\label{04-09-25x} $f_{\underline{G}}(D^y) =F(y)$ for all $y \in Y$.
\end{lemma}
\begin{proof} It follows from the definitions that
\begin{align*}
& f_{\underline{G}}(D^y) = \left\{L_{\underline{G}}(x) : \ x \in X_{\underline{G}}[0,\infty), \ s_{\underline{G}}(x) = D^y\right\} \\
& =  \left\{ L_G(x') : \ x' \in X_G[0,\infty), \ s_G(x') \in D^y\right\} \\
& = \left\{ z_{[0,\infty)}: \ z \in Y, \ z_j = y_j, \ j \leq -1 \right\} = F(y).
\end{align*}
\end{proof}

Note that $t_{\underline{G}}(\alpha_G(y)_{(-\infty,-1]}) = s_{\underline{G}}(\alpha(y)_0) = D^y$ and that $L_{\underline{G}}(\alpha_G(y)) = y$ for all $y \in Y$. It follows therefore from Lemma \ref{04-09-25x} that $(\underline{G},L_{\underline{G}})$ is regular. Hence $([\underline{G}], L_{[\underline{G}]})$ is right-resolving, follower-separated and regular by Lemma \ref{10-09-25}. It follows therefore from Proposition \ref{25-09-25} and its proof that there is an isomorphism $\theta : (\underline{G},L_{\underline{G}}) \to (\mathbb K(Y),L_{\mathbb K(Y)})$ of labeled graphs defined such that 
$$
\theta(f_{\underline{G}}(D^y)) = F(y) ,
$$
allowing us to make the following identification which we shall use from now on:

\begin{corollary}\label{04-09-25d}  $(\mathbb K(Y),L_{\mathbb K(Y)}) = ([\underline{G}],L_{[\underline{G}]})$ as labeled graphs.
\end{corollary}

























































\section{A strongly canonical extension of the future cover}

Let $(G,L_G)$ be a presentation of the sofic shift $Y$. It follows from Section \ref{consec} that  there is a factor map $f_{\underline{G}} : X_{\underline{G}} \to X_{\mathbb K(Y)}$. Assume now that $\psi : Y \to Z$ is a conjugacy of sofic shifts and that $(G,L_G)$ and $(H,L_H)$ are right-resolving presentations of $Y$ and $Z$, respectively, such that
\begin{equation}\label{09-10-25axx}
\xymatrix{
 X_G\ar[r]^-{\phi} \ar[d]_{L_G}& X_{H} \ar[d]^-{L_H}\\
   Y \ar[r]^{\psi} & Z}
\end{equation}
commutes. We aim to prove

\begin{thm}\label{18-01-26d} There is a unique conjugacy $\tilde{\phi} : X_{\underline{G}} \to X_{\underline{H}}$ such that
\begin{equation}\label{18-01-26e}
\xymatrix{
 X_{\underline{G}} \ar[r]^-{\tilde{\phi}} \ar[d]_{L_{\underline{G}}}& X_{\underline{H}} \ar[d]^-{L_{\underline{H}}}\\
   Y \ar[r]^{\psi} & Z}
\end{equation}
commutes. This conjugacy has the property that also
\begin{equation}\label{20-01-25d}
\xymatrix{\\
X_{\underline{G}} \ar[r]^-{\tilde{\phi}} \ar[d]_-{f_{\underline{G}}} & X_{\underline{H}}  \ar[d]^-{f_{\underline{H}}}\\
 X_{\mathbb K(Y)} \ar[d]_-{L_{\mathbb K(Y)}} \ar[r]^-{\psi_{\mathbb K}} &  X_{\mathbb K(Z)} \ar[d]^-{L_{\mathbb K(Z)}} \\
Y \ar[r]_-{\psi}  & Z }
\end{equation}
commutes.
\end{thm}


The proof will proceed through a series of lemmas, using material from the previous sections. Recall that the labeled graphs $(G',L_{G'})$ and $(\underline{G},L_{\underline G})$ are both subgraphs of $(\overline{G},L_{\overline{G}})$. While $(\underline{G},L_{\underline{G}})$ is obtained from $(\overline{G},L_{\overline{G}})$ by restricting to a subset of the vertexes in $\overline{G}$, $(G',L_{G'})$ comes from a restriction both to a subset of vertexes and arrows in $\overline{G}$. Working with these graphs, the maps $\beta_G$ and $\alpha_G$ that were defined in Section \ref{subset} and Section \ref{consec}, respectively, will be important tools.



\begin{lemma}\label{08-10-25bxz} For all $y \in Y$ the element $\beta_G(y)$ is forward asymptotic to $\alpha_G(y)$.
\end{lemma}
\begin{proof} Note that
$$
s_{\underline{G}}(\alpha_G(y)_n) = \left\{t_G(x): \ x \in X_G(-\infty,n-1], \ L_G(x) = y_{(-\infty,n-1]}\right\} \supseteq  s_{G''}(\beta_G(y)_n) 
$$
for all $n \in \mathbb Z$. A compactness argument shows that there is an $N \in \mathbb Z$ such that
$$
\left\{t_G(x): \ x \in X_G(-\infty,n-1], \ L_G(x) = y_{(-\infty,n-1]}\right\} = s_{G''}(\beta_G(y)_n) 
$$
for all $n \geq N$. It follows that $\alpha_G(y)_j = \beta_G(y)_j$ for all $j \geq N$.
\end{proof}



When $y$ is periodic, both $\alpha_G(y)$ and $\beta_G(y)$ are periodic and hence Lemma \ref{08-10-25bxz} has the following

\begin{corollary}\label{09-10-25b} Let $p \in \Per (Y)$. Then $\beta_G(p) = \alpha_G(p)$.
\end{corollary}

The proof of the following lemma is similar to that of Lemma \ref{11-09-25ix}.

\begin{lemma}\label{10-02-25dxxz} Let $v\in V_{\underline{G}}$. There is a periodic point $p \in \Per (Y)$ and an element $y' \in X_{\underline{G}}$ backward asymptotic to $\alpha_G(p)$ such that $t_{\underline{G}}(y'_{(-\infty,-1]}) = v$.
\end{lemma}
\begin{proof} Since $v \in V_{\underline{G}}$ there is an element $y \in Y$ such that $v = s_{\underline{G}}(\alpha_G(y)_0)$. Set $H_j := s_{\underline{G}}(\alpha(y)_j), \ j \leq 0$. Since $(G,L_G)$ is right-resolving there is an $N <0$ such that 
\begin{equation}\label{17-04-25bxz}
\# H_j = \# H_N
\end{equation} 
for all $j \leq N$.
 By the pigeon hole principle there is a sequence $\cdots < n_3 < n_2 <n_1 < n_0 \leq N$ and a vertex $H \in V_{\underline{G}}$ such that $H =s_{\underline{G}}(\alpha_G(y)_{n_k})$ for all $k \in \mathbb N \cup \{0\}$. Consider the loops in $\underline{G}$:
$$
\gamma_k := H \overset{y_{n_k}}{\to} H_{n_k+1}  \overset{y_{n_k+1}}{\to} \cdots \cdots \overset{y_{n_0-1}}{\to} H .
$$
If $\# L_G^{-1}(L_{\underline{G}}(\gamma_k)^\infty) > \# H$ for all $k$ it follows that there is an element $x' \in X_{\overline{G}}(-\infty,{n_0}-1]$ such that $H \subsetneq t_{\overline{G}}(x')$ and $L_{\overline{G}}(x') = y_{(-\infty,n_0-1]}$. By definition of $\alpha_G(y)$ this implies that $\# s_{\underline{G}}(\alpha_G(y)_j)  = \# H_j > \# H =\# H_N $ for all $j \leq n_0$, contradicting \eqref{17-04-25bxz}. This shows that there is a $k$ such that $q:= L_{\underline{G}}((\gamma_k)^\infty) \in Y$ is periodic and $\# L_G^{-1}(q) = \#H$. Thus $\alpha_G(q) = (\gamma_k)^\infty$ and 
$$
y' := \alpha_G(\sigma^{n_0}(q))_{(-\infty,n_0-1]}\alpha_G(y)_{[n_0,\infty)}
$$
has the stated properties.
\end{proof}


\begin{lemma}\label{11-11-24axz} Let $x,y \in X_{\underline{G}}$, and assume that $x = \alpha_G(L_{\underline{G}}(x))$. If there is an $N \in \mathbb Z$ such that $x_i = y_i$ for $i \leq N$, then $y = \alpha_G(L_{\underline{G}}(y))$.
\end{lemma}
\begin{proof} Since $L_{\underline{G}}(x)_i = L_{\underline{G}}(y)_i$ for $i \leq N$, it follows from the definition of 
$\alpha_G$ that 
$$
\alpha_G(L_{\underline{G}}(x))_j = \alpha_G(L_{\underline{G}}(y))_j
$$ 
for $j \leq N-1$. Hence 
$$
t_{\underline{G}}(y_{N-1}) =t_{\underline{G}}(x_{N-1}) = t_{\underline{G}}\left(\alpha_G(L_{\underline{G}}(x))_{N-1}\right) =  t_{\underline{G}}\left(\alpha_G(L_{\underline{G}}(y))_{N-1}\right).
$$
Since $L_{\underline{G}}(y_{[N,\infty)}) = L_{\underline{G}}\left(\alpha_G(L_{\underline{G}}(y))_{[N,\infty)}\right)$ it follows that $y_{[N,\infty)} = \alpha_G(L_{\underline{G}}(y))_{[N,\infty)}$ because $(\underline{G},L_{\underline{G}})$ is right resolving.  We know already that $y_{(-\infty,N-1]} = x_{(-\infty,N-1]} = \alpha_G(L_{\underline{G}}(x))_{(-\infty,N-1]} =  \alpha_G(L_{\underline{G}}(y))_{(-\infty,N-1]}$, and conclude therefore that $y = \alpha_G(L_{\underline{G}}(y))$.
\end{proof}


\begin{lemma}\label{03-07-25axz} Let $x \in X_{\underline{G}}$. There is an element $y \in Y$ such that $\beta_G(y)$ is forward asymptotic to $x$.
\end{lemma}
\begin{proof} It follows from
Lemma \ref{10-02-25dxxz} there is an element $y' \in X_{\underline{G}}$ and a periodic point $p\in Y$ such that $y'$ is backward asymptotic to $\alpha_G(p)$ and forward asymptotic to $x$. Since $\alpha_G(p) =  \alpha_G(L_{\underline{G}}(\alpha_G(p)))$, it follows from Lemma \ref{11-11-24axz} that $y' = \alpha_G(L_{\underline{G}}(y'))$. Set $y := L_{\underline{G}}(y')$ and note that $\beta_G(y)$ is forward asymptotic to $\alpha_G(y) = y'$ by Lemma \ref{08-10-25bxz}. Hence $\beta_G(y)$ is also forward asymptotic to $x$.
\end{proof}


\begin{lemma}\label{08-10-25az}
$$
\overline{\Per X_{\underline{G}}}
 = \bigcup_{y \in Y} \bigcap_{N \geq 1} \overline{\left\{\sigma^n(\beta_G(y)) : \ n \geq N \right\}} .
$$
\end{lemma}
\begin{proof} Let $y \in Y$. It follows from Lemma \ref{08-10-25bxz} that
$$
\bigcap_{N \geq 1} \overline{\left\{\sigma^n(\beta_G(y)) : \ n \geq N \right\}} = \bigcap_{N \geq 1} \overline{\left\{\sigma^n(\alpha_G(y)) : \ n \geq N \right\}}.
$$
Since there is a component of $X_{\underline{G}}$ which contains $\sigma^n(\alpha_G(y))$ for all sufficiently large $n$, it follows that 
$$
 \bigcap_{N \geq 1} \overline{\left\{\sigma^n(\beta_G(y)) : \ n \geq N \right\}} \subseteq \overline{\Per X_{\underline{G}}} .
 $$
This proves the inclusion $\supseteq$. The obtain the reversed inclusion, let $C$ be a component in $X_{\underline{G}}$. Choose an element $x \in X_C$ which is forward transitive in $X_C$; i.e. every word from $\mathbb W(X_C)$ occurs infinitely many times to the right in $x$. It follows from Lemma \ref{03-07-25axz} that there is an element $y \in Y$ such that 
\begin{equation}\label{08-10-25dz}
\lim_{n \to \infty} d\left(\sigma^n(\beta_G(y)),\sigma^n(x)\right) = 0.
\end{equation}
For each $z \in X_C$ there is a strictly increasing sequence $\{n_i\}$ in $\mathbb N$ such that
$\lim_{i \to \infty} d(\sigma^{n_i}(x),z) = 0$.
Thanks to \eqref{08-10-25dz}, by passing to a subsequence we can then arrange that
$$
\lim_{i \to \infty} d(\sigma^{n_i}(\beta_G(y)),z) = 0.
$$
It follows that $z \in  \bigcap_{N \geq 1} \overline{\left\{\sigma^n(\beta_G(y)) : \ n \geq N \right\}}$.
\end{proof}


\begin{corollary}\label{01-11-25x}  $\overline{\Per X_{\underline{G}}} \subseteq X_{G'} \cap X_{\underline{G}}$
\end{corollary}





\begin{lemma}\label{27-01-26b} $f_{\underline{G}} \circ \alpha_G = \alpha_Y$.
\end{lemma}
\begin{proof} Let $y \in Y$. Then $L_{\mathbb K(Y)}( f_{\underline{G}} \circ \alpha_G(y)) = L_{\underline{G}}(\alpha_G(y)) = y = L_{\mathbb K(Y)}(\alpha_Y(y))$, and 
$$
s_{\mathbb K(Y)} (f_{\underline{G}} \circ \alpha_G(y)_k) = f_{\underline{G}}(D^{\sigma^k(y)}) = F(\sigma^k(y)) = s_{\mathbb K(Y)}(\alpha_Y(y)_k)
$$
for all $k \in \mathbb Z$. Hence  $f_{\underline{G}} \circ \alpha_G(y) = \alpha_Y(y)$
\end{proof}

















\begin{lemma}\label{09-10-25fz} 
$$
\phi'( \overline{\Per X_{\underline{G}}}) = \overline{\Per X_{\underline{H}}} .
$$
and 
$$
f_{\underline{H}} \circ \phi'(x) = \psi_{\mathbb K} \circ f_{\underline{G}}(x)
$$
for all $x \in \overline{\Per X_{\underline{G}}}$.
\end{lemma}



\begin{proof} The first equality follows from Lemma \ref{08-10-25az} and \eqref{30-10-25d}. Let $z \in \overline{\Per X_{\underline{G}}}$. By Lemma \ref{08-10-25az} there is an $y \in Y$ such that $z = \lim_{i \to \infty} \sigma^{n_i}(\beta_{G}(y)) $ for some sequence $\{n_i\}$ in $\mathbb N$.
Note that 
$$
\lim_{i \to \infty} \sigma^{n_i}(\beta_{H}(y)) = \lim_{i \to \infty} \sigma^{n_i}(\alpha_{H}(y))
$$
by Lemma \ref{08-10-25bxz}. From \eqref{30-10-25d} we see that 
$\phi' \circ \beta_G = \beta_H \circ \psi$. 
Hence
\begin{align*}
&\phi'( \lim_{i \to \infty} \sigma^{n_i}(\beta_{G}(y))) = \lim_{i \to \infty} \sigma^{n_i}(\beta_{H}(\psi(y))) =  \lim_{i \to \infty} \sigma^{n_i}(\alpha_{H}(\psi(y)),
\end{align*}
where the last identity follows from Lemma \ref{08-10-25bxz}.
Note that $f_{\underline{H}} \circ \alpha_H \circ \psi = \alpha_{Z} \circ \psi$ by Lemma \ref{27-01-26b}. Since $\alpha_{Z}\circ\psi = \psi_{\mathbb K} \circ \alpha_Y$ by \eqref{24-11-24x} we conclude that
\begin{align*}
&f_{\underline{H}} \circ \phi'(z) = f_{\underline{H}} \circ \phi'( \lim_{i \to \infty} \sigma^{n_i}(\beta_{G}(y))) = \lim_{i \to \infty} f_{\underline{H}}(\sigma^{n_i}(\alpha_{H}(\psi(y))) \\
&  = \lim_{i \to \infty} \sigma^{n_i}\left( f_{\underline{H}} \circ \alpha_H \circ \psi(y)\right)= \lim_{i \to \infty}  \sigma^{n_i}(\alpha_{Z}(\psi(y)))= \lim_{i \to \infty}  \sigma^{n_i}(\psi_{\mathbb K}\circ\alpha_Y(y)) \\
& =  \psi_{\mathbb K}\left(\lim_{i \to \infty} \sigma^{n_i}(\alpha_{Y}(y))\right) =  \psi_{\mathbb K}\left(\lim_{i \to \infty} \sigma^{n_i}(f_{\underline{G}} \circ \alpha_{G}(y))\right) = \psi_{\mathbb K}\circ f_{\underline{G}}(z).
\end{align*}
%For comparison 
%\begin{align*}
%&\phi \circ f_{G}\left( \lim_{i \to \infty} \sigma^{n_i}(\beta_{G}(y))\right) \\
%& = \phi\left(  \lim_{i \to \infty} \sigma^{n_i}(f_{G} \circ \beta_{G}(y))\right)\\
%& = \phi \left( \lim_{i \to \infty} \sigma^{n_i}(\alpha_Y(y))\right)\\
%&  = \phi \left( \lim_{i \to \infty} \sigma^{n_i}(f_G\circ \alpha_G(y))\right)\\
%&= \phi \circ f_G(z) 
%\end{align*}
\end{proof}


\begin{cor}\label{30-01-26c} $L_{\underline{H}} \circ \phi'(x) = \psi\circ L_{\underline{G}}(x)$ for $x \in \overline{\Per (X_G)}$.
\end{cor}
\begin{proof} For $x \in \overline{\Per (X_G)}$, 
\begin{align*} &L_{\underline{H}} \circ \phi'(x) = L_{\mathbb K(Z)} \circ f_{\underline{H}} \circ \phi'(x) = L_{\mathbb K(Z)} \circ \psi_\mathbb K \circ f_{\underline{G}}(x) \\
& = \psi\circ L_{\mathbb K(Y)} \circ f_{\underline{G}}(x) = \psi \circ L_{\underline{G}}(x).
\end{align*}
\end{proof}

%This leads to the following:

%\begin{lemma}\label{09-10-25gz} The map $\phi' : \overline{\Per X_{\underline{G}}} \to \overline{\Per X_{\underline{H}}}$ is a conjuagcy and the diagram 
%\begin{equation}\label{09-10-25xz}
%\xymatrix{\\
%\overline{\Per (X_{\underline{G}})} \ar[r]^-%{\phi'} \ar[d]_-{f_{\underline{G}}} & %\overline{\Per (X_{\underline{H}})}  \ar[d]^-%{f_{\underline{H}}}\\
% \overline{\Per X_{\mathbb K(Y)}} \ar[d]_-%{L_{\mathbb K(Y)}} \ar[r]^-{\psi_{\mathbb K}} &  %\overline{\Per X_{\mathbb K(Z)}} \ar[d]^-%{L_{\mathbb K(Z)}} \\
%\overline{\Per Y} \ar[r]_-{\psi}  & %\overline{\Per Z }}
%\end{equation}
%commutes.
%\end{lemma}



\begin{lemma}\label{20-01-26d} Let $C$ be a source component in $\underline{G}$. Then
$$
\# L_G^{-1}\left(L_{\underline{G}}(x)\right) = M(C)
$$
for all $x \in X_C$.
\end{lemma}
\begin{proof} Since $C$ is a source component it follows from Lemma \ref{10-02-25dxxz} that there is a periodic point $p \in \Per Y$ such that $\alpha_G(p) \in X_C$. Since $\alpha_G(p) = \beta_G(p)$ by Corollary \ref{09-10-25b} it follows that $\# L_G^{-1}(p) = \# s_{\underline{G}}(\alpha_G(p)_j)$ and hence that
\begin{equation}\label{20-01-26e}
 \# L_G^{-1}(p) = M(C)
\end{equation} 
for all $j \in \mathbb Z$. Assume now that there is an $x \in X_C$ such that 
\begin{equation}\label{20-01-26f}
\# L_G^{-1}(L_{\underline{G}}(x)) > M(C). 
\end{equation}
Since $\underline{G}_C$ is strongly connected there is an element $z \in X_C$ and an $m \in \mathbb N$ such that $z_i = x_i$ for $i \leq 0$ while $z_i = \alpha_G(p)_i$ for $i \geq m$. It follows from \eqref{20-01-26f} and the definition of $\alpha_G$ that 
\begin{equation}\label{20-01-26g}
s_{\underline{G}}(\alpha_G(L_{\underline{G}}(z))_i) \geq M(C)+1
\end{equation} 
for all $i < 0$. Since $L_{\underline{G}}(\alpha_G(z)_i) = p_i$ for $i \geq m$, it follows from \eqref{20-01-26e} and the periodicity of $p$ that $\#s_{\underline{G}}(\alpha_G(L_{\underline{G}}(z))_i) \leq M(C)$ for all sufficiently large $i$. However, since $s_{\underline{G}}(\alpha_G(L_{\underline{G}}(z))_i)  \supseteq s_{\underline{G}}(z_i)$ and $\# s_{\underline{G}}(z_i) = M(C)$ for all $i$, we find that $s_{\underline{G}}(\alpha_G(L_{\underline{G}}(z))_i)  =s_{\underline{G}}(z_i)$ and hence also that $\alpha_G(L_{\underline{G}}(z))_i = z_i = p_i$ for all sufficiently large $i$. Since $C$ is a source component it follows that $s_{\underline{G}}(\alpha_G(L_{\underline{G}}(z)_i) \in C$ for all $i$; a contradiction to \eqref{20-01-26g}. Hence $\# L_G^{-1}(L_{\underline{G}}(x)) = M(C)$ for all $x \in X_C$. 
\end{proof}


\begin{cor}\label{20-01-26h} Let $C,C'$ be source components in $\underline{G}$. Then
\begin{itemize}
\item[1)] $z = \beta_G(L_{\underline{G}}(z))$ for all $z \in X_C$,
\item[2)] $L_{\underline{G}} : X_C \to Y$ is injective and hence a conjugacy of $X_C$ onto $L_{\underline{G}}(X_C)$, and
\item[3)] $L_{\underline{G}}(X_C) \cap L_{\underline{G}}(X_{C'}) = \emptyset$ if $C \neq C'$.
\end{itemize}
\end{cor}
\begin{proof} This follows readily from Lemma \ref{20-01-26d}.
\end{proof}


\begin{lemma}\label{29-01-26} Let $C$ a source component in $\underline{G}$. If $z \in X_{\underline{G}}$ is backward asymptotic to $X_C$, then $z =\alpha_G\circ L_{\underline{G}}(z)$.
\end{lemma}
\begin{proof} By Lemma \ref{11-11-24axz} we may assume that $z \in X_C$. It follows then from Lemma \ref{20-01-26d} that $z = \beta_G(L_{\underline{G}}(z)$. Let $C'$ be the component in $\underline{G}$ such that $\alpha_G(L_{\underline{G}}(z))$ is backward asymptotic to $X_{C'}$. Since $\alpha_G(L_{\underline{G}}(z))$ is forward asymptotic to $\beta_G(L_{\underline{G}}(z)) =z$ by Lemma \ref{08-10-25bxz}, it follows that $C' = C$ since $C$ is a source component. Since $s_{\underline{G}}(z_j) \subseteq s_{\underline{G}}(\alpha_G(L_{\underline{G}}(z))_j)$ for all $j$, it follows that $s_{\underline{G}}(z_j) = s_{\underline{G}}(\alpha_G(L_{\underline{G}}(z))_j)$ for all $j$ 'close to $-\infty$', and hence $z = \alpha_G\circ L_{\underline{G}}(z))$ 
\end{proof}

\begin{lemma}\label{29-01-26a} Let $C$ be a component in $\underline{G}$ with the property that $z=\alpha_G(L_{\underline{G}}(z))$ for all $z \in X_C$, and let $x \in X_{\underline{G}}$ be backward asymptotic to $X_C$. Then $x =\alpha_G\circ L_{\underline{G}}(x))$.
\end{lemma}
\begin{proof} This follows from Lemma \ref{11-11-24axz}. 
\end{proof}


\begin{lemma}\label{29-01-26b} Let $C$ be a component in $\underline{G}$ with the property that $p=\beta_G(L_{\underline{G}}(p))$ for all $p  \in \Per (X_C)$. Then $z=\alpha_G(L_{\underline{G}}(z))$ for all $z \in X_C$.
\end{lemma}
\begin{proof} Let $z \in X_C$. There are $N \in \mathbb Z$ and $m \in \mathbb N$ such that $\# s_{\underline{G}}(\alpha_G(L_{\underline{G}}(z))_j) = m$ for all $j \leq N$. There are then $j_1 < j_2 \leq N$ such that  $\alpha_G(L_{\underline{G}}(z))_{j_1}=\alpha_G(L_{\underline{G}}(z))_{j_2}$ and $z_{j_1} = z_{j_2}$. Then $p:=(z_{[j_1,j_2-1]})^{\infty} \in X_C$ is a periodic point with $\# L_G^{-1}(L_{\underline{G}}(p)) \geq m$, implying that $M(C) \geq m$ since $\# L_G^{-1}(L_{\underline{G}}(p)) = M(C)$. It follows that $s_{\underline{G}}(z_j)$  and $s_{\underline{G}}(\alpha_G(L_{\underline{G}}(z))_{j})$ have the same number of elements for all $j$ 'close to $-\infty$', implying that $z=\alpha_G(L_{\underline{G}}(z))$. 
\end{proof}



\begin{lemma}\label{29-01-26c} Let $C$ be a component in $\underline{G}$ with the property that $p=\beta_G(L_{\underline{G}}(p))$ for all $p \in \Per (X_C)$, and let $x \in X_{\underline{G}}$ be backward asymptotic to $X_C$. Then $x =\alpha_G\circ L_{\underline{G}}(x))$.
\end{lemma}
\begin{proof} This follows from Lemma \ref{29-01-26b} and Lemma \ref{11-11-24axz}. 
\end{proof}



Let $\mathcal C_G$ be the set of components in $\underline{G}$ with the property that $p=\beta_G(L_{\underline{G}}(p))$ for all $p \in \Per (X_C)$, and let $\mathcal C_H$ be the corresponding set of components in $\underline{H}$. Then $\phi'(C) \in \mathcal C_H$ when $C \in \mathcal C_G$. This follows because by  Lemma \ref{09-10-25fz} and \eqref{28-12-25c}:
$$
\beta_H \circ L_{\underline H} \circ \phi' = \beta_H \circ \psi \circ L_{\underline{G}} = \phi' \circ \beta_G \circ L_{\underline{G}}
$$ 



\begin{cor}\label{29-01-26d} Let $C,C' \in \mathcal C_G$. Then
\begin{itemize}
%\item[1)] $z = \beta_G(L_{\underline{G}}(z))$ for all $z \in X_C$,
\item[1)] $L_{\underline{G}} : X_C \to Y$ is injective and hence a conjugacy of $X_C$ onto $L_{\underline{G}}(X_C)$, and
\item[2)] $L_{\underline{G}}(X_C) \cap L_{\underline{G}}(X_{C'}) = \emptyset$ if $C \neq C'$.
\end{itemize}
\end{cor}























\begin{lemma}\label{24-01-26} $(\underline{G}, L_{\underline{G}})$ is regular and $\alpha_G(Y) \subseteq \mathcal R\left(L_{\underline{G}}\right)$.
\end{lemma}
\begin{proof} Let $y \in Y$ and note that $D^y = s_{\underline{G}}(\alpha_G(y)_0)$. It follows therefore from Lemma \ref{12-09-25b} that it suffices here to show that $\alpha_G(y) \in X_{\underline{G}}$ is regular; i.e. it suffices to show that 
$$
F(\sigma^{k+1}(y)) = F\left(\sigma^{k+1}(L_{\underline{G}}(\alpha_G(y)))\right) \subseteq f_{\underline{G}}(t_{\underline{G}}(\alpha_G(y)_k))
$$
for $k \in \mathbb Z$. Let $ z \in F(\sigma^{k+1}(y))$. Define $w \in Y$ such that
$$
w_i = \begin{cases} y_i,\ i \leq k \\ z_{i-k-1}, \ i \geq k+1. \end{cases}
$$
Then $\alpha_G(w) \in X_{\underline{G}}$ and $L_{\underline{G}}(\alpha_G(w)_{[k+1,\infty)}) = z$. Hence $z \in f_{\underline{G}}(t_{\underline{G}}(\alpha_G(y)_k))$.
\end{proof}


Let $\mathcal X_G$ be the set of elements in $X_G$ that are backward asymptotic to $\bigcup_{C  \in \mathcal C_G} X_C$.

\begin{lemma}\label{30-01-26} $L_{\underline{G}} : \ \mathcal X_{\underline{G}} \to Y$ is injective and maps onto the set of elements in $Y$ that are backward asymptotic to $\bigcup_{C  \in \mathcal C_G} L_{\underline{G}}(X_C)$. In fact,
$$
\chi = \alpha_H \circ \psi \circ L_{\underline{G}} 
$$
on $\mathcal X_G$.
\end{lemma}
\begin{proof} Let $x,y \in \mathcal X_G$ and assume that $L_{\underline{G}}(x) = L_{\underline{G}}(y)$. It follows follows from Lemma \ref{29-01-26c} that $x = \alpha_G(L_{\underline{G}}(x)) =  \alpha_G(L_{\underline{G}}(y)) =y$, proving in the injectivity part. Let $z \in Y$ be backward asymptotic to $\bigcup_{C  \in \mathcal C_G} L_{\underline{G}}(X_C)$. There is a $C \in \mathcal C_G$, an element $x \in \mathcal X_C$ and a $N \in \mathbb Z$ such that $z_i = L_{\underline{G}}(x)_i$ for all $i \leq N$. Then $\alpha_G(z)_i = \alpha_G(L_{\underline{G}}(x))_i$, and hence, using Lemma \ref{29-01-26c}, $\alpha_G(z)_i = x_i$ for all $i \in N$. Hence $\alpha_G(z) \in \mathcal X_G$, completing the proof because $L_{\underline{G}}(\alpha_G(z)) = z$.
\end{proof}


\begin{lemma}\label{30-01-26a} There is a shift-commuting bijection $\chi : \mathcal X_G \to \mathcal X_H$ such that $L_{\underline{H}} \circ \chi = \psi \circ L_{\underline{G}}$.
\end{lemma}
\begin{proof} As observed above, $\phi'(\bigcup_{C \in \mathcal C_G} X_C) = \bigcup_{C \in \mathcal C_H} X_C$. It follows therefore from Corollary \ref{30-01-26c} that 
$$
\psi \circ L_{\underline{G}}\left( \bigcup_{C \in \mathcal C_G} X_C\right) = L_{\underline{H}}\left(\bigcup_{C \in \mathcal C_H} X_C\right) .
$$
Hence, by Lemma \ref{30-01-26}, $\psi(L_{\underline{G}}(\mathcal X_G)) = L_{\underline{H}}(\mathcal X_H)$ and we can set $\chi := L_{\underline{H}}^{-1} \circ \psi \circ L_{\underline{G}}$ on $\mathcal X_G$.
\end{proof}


Let $x \in X_{\underline{G}}$. A \emph{homogenous interval} $I$ in $x$ is an interval $I = [i,j]$ in $\mathbb Z$ such that $\# s_{\underline{G}}(x_i) = \# s_{\underline{G}}(x_k)$ for all $k \in I$. 

\begin{lemma}\label{01-02-26} Let $I$ and $J$ be homogeneous intervals in $x \in X_{\underline{G}}$ and $i < j$ integers such that $[i-d,i+d] \subseteq I, \ [j-d,j+d] \subseteq J$. There is a unique path $\gamma$ in $X_{\underline{H}}$ such that $s_{\underline{H}}(\gamma) = s_{\underline{H}}(\phi''(x_{[i-d,i+d]})_i)$,  $t_{\underline{H}}(\gamma) = t_{\underline{H}}(\phi''(x_{[j-d,j+d]})_j)$ and $L_{\underline{H}}(\gamma) = \psi(L_{\underline{G}}(x))_{[i,j]}$.
\end{lemma}
\begin{proof} Consider $s_{\underline{G}}(x_I)_i$ and let $N$ be path in $G''$ with most elements such that $N_I \subseteq x_I$ and $t_{G''}(N) = t_{\underline{G}}(x_J)$. Then $\phi''(N)_I \subseteq \phi'(x_I)$ and $L_{\underline{H}}(N) = \psi(L_{G''}(N)) = \psi(L_{G''}(x)_{[i,j]})$. Note that if we arrange (!!!! - assumptions of $x_I$) that $\phi'(x_I)$ runs in $\underline{H}$ there is a unique path $\gamma$ in $\underline{H}$ such that $s_{\underline{H}}(\gamma) = s_{\underline{H}}(\phi''(x_{[i-d,i+d]})_i)$,  and $L_{\underline{H}}(\gamma) = \psi(L_{\underline{G}}(x))_{[i,j]}$. It remains to show that $t_{\underline{H}}(\gamma) = t_{\underline{H}}(\phi''(x_{[j-d,j+d]})_j)$. Note that $\gamma$ terminates at something $A$ which contains $t_{G''}(\phi''(N))$. Let $\mu$ be the path with most elements in ${H''}$ such that $\mu_I \subseteq \phi''(x_I)$ and $\mu$ has label $\psi(L_{\underline{H}}(x_{[i-d,j+d]})$. Then $N_I \subseteq(\phi'')^{-1}(\mu)_I \subseteq x_I$.  Considerations show that  $(\phi'')^{-1}(\mu) = N$, implying that $t_{\underline{H}}(\gamma) = t_{\underline{H}}(\phi''(x_{[j-d,j+d]})_j)$. Since $L_{G''}(\mu)_{[i,j]} = L_{\underline{G}}(x))_{[i,j]}$ we must have that $N_{[i,j]} =(\phi'')^{-1}(\mu)_{[i,j]}$. ++++++++++++++++++++++++++++++
\end{proof}


\begin{lemma}\label{02-02-26}
\end{lemma}






+++++++++++++++++++++++++++++++++++
In the next lemma we extend $\phi'$ from $\overline{\Per X_{\underline{G}}}$ to $X_{\underline{G}}$; the problem being, of course, that $X_{\underline{G}}$ is not a subshift of $X_{G''}$.

\begin{lemma}\label{19-01-26} There is a sliding block code $\tilde{\phi} : \ X_{\underline{G}} \to X_{\underline{H}}$ such that
\begin{itemize}
\item[1)] $\tilde{\phi}$ extends $\phi' : \overline{\Per X_{\underline{G}}} \to \overline{\Per X_{\underline{H}}}$, and
\item[2)] $L_{\underline{H}} \circ \tilde{\phi} = \psi \circ L_{\underline{G}}$.
\end{itemize}
Subject to these two conditions, the sliding block code $\tilde{\phi}$ is unique. 
\end{lemma}
\begin{proof} Recall that $\phi'$ is the restriction of $\phi''$ to $X_{G'} \supseteq \overline{\Per X_{\underline{G}}}$ and that the restriction of $\phi''$ to $X_G \subseteq X_{G''}$ is $\phi$. Let $M$ be larger than the memory and anticipation of both $\phi'': X_{G''} \to X_{H''}$ and $\psi$, cf. Definition 1.5.1 of \cite{LM}. We shall use the notation $\phi''$ also for the block map $\mathbb W_{2M+1}(X_{G''}) \to E_{H''}$ describing $\phi''$. Thus
$$
\phi''(y)_j = \phi''(y_{[j-M,j+M]})
$$
for all $y \in X_{G''}$ and all $j \in \mathbb Z$.


 Since $\underline{G}$ is a finite oriented graph there is an $N \in \mathbb N$ such that any path in $\underline{G}$ of length at least $N$ contains a subpath of length at least $3M$ which is contained in one of the components of $\underline{G}$. We shall construct the desired sliding block code $\tilde{\phi}$ to have memory $N$ and anticipation $M$. Consider therefore an element $x \in X_{\underline{G}}$ and an $i \in \mathbb Z$. There is then a component $C$ in $\underline{G}$ and an $i' \in [i-N,i]$ such that $s_{\underline{G}}(x_j) \in C$ for all $j \in [i'-3M,i']$. Let $\tilde{x} \in X_C$ be an element such that $\tilde{x}_{[i'-3M,i']} =x_{[i'-3M,i']}$. Then
$$
\phi'(x_{[j-M,j+M]}) = \phi''(\tilde{x})_j
$$
for $j \in [i'-2M,i'-M]$. 

We want to prove that there is a vertex 
$$
v \in s_{H''}(\phi'(x_{[i'-2M,i']})) = s_{H''}(\phi''(\tilde{x})_{i'-M})
$$
and a right-infinite path $\gamma$ in $H$ such that $s_H(\gamma) = v$ and $L_H(\gamma) =  \psi(L_{\underline{G}}(x))_{[i'-M,\infty)}$. For this we choose first an element $z \in X_G$ such that $L_G(z) = L_{\underline{G}}(x)$ and $s_G(z_j) \in s_{\underline{G}}(x_j)$ for all $j \in \mathbb Z$. Since $\tilde{x} \in X_C$ and $\tilde{x}_{[i'-3M,i']} =x_{[i'-3M,i']}$ there is a unique element $p \in X_G$ such that $s_G(p_j) = s_G(z_j)$ for $j \in [i'-3M,i']$,
\begin{equation}\label{19-01-26ax}
s_G(p_j) \in s_{G''}(\tilde{x}_j)
\end{equation} 
for all $j \in \mathbb Z$ and $L_G(p) =L_{\underline{G}}(\tilde{x})$. It follows from \eqref{19-01-26ax} and the definition of $\phi''$ that $s_H(\phi(p)_j) \in s_{H''}(\phi''(\tilde{x})_j)$ for all $j \in \mathbb Z$ and hence also that 
$$
s_H(\phi(z)_j) \in s_{H''}(\phi''(\tilde{x})_j)
$$ 
for $j \in [i'-2M,i'-M]$. Since $L_H(\phi(z))_j = \psi(L_G(z))_j = \psi(L_{\underline{G}}(x))_j$ for all $j$, we see that $v:= s_H(\phi(z)_{i'-M})$ is a vertex with the properties we seek and that $\gamma:= \phi(z)_{[i'-M,\infty)}$ is the corresponding right-infinite path.

 The existence of the pair $(v,\gamma)$ and the definition of $\underline{H}$ imply that there is a unique element $y \in X_{\underline{H}}[i'-M,\infty)$ such that $s_{\underline{H}}(y_{i'-M}) = s_{\underline{H}}(\phi''(\tilde{x})_{i'-M}) =  s_{H''}(\phi'(x_{[i'-2M,i']}))$ and $L_{\underline{H}}(y) = \psi(L_{\underline{G}}(x))_{[i'-M,\infty)}$. Now define $\tilde{\phi}(x)_i$ by
$$
\tilde{\phi}(x)_i : = y_ i.
$$
With this definition $\tilde{\phi}(x)_i$ depends only on $x_{[i-N,i+M]}$ and $\tilde{\phi}$ is therefore a sliding block code $\tilde{\phi} : \ X_{\underline{G}} \to X_{\underline{H}}$ with memory $N$ and anticipation $M$. By construction $L_{\underline{H}}(\tilde{\phi}(x)_i) = L_{\underline{H}}(y)_i = \psi(L_{\underline{G}}(x))_i$, proving that $L_{\underline{H}} \circ \tilde{\phi} = \psi \circ L_{\underline{G}}$. To establish property 1), note that when $x \in X_C$ for some component $C$ in $\underline{G}$ the element $y$ is equal to $\phi'(x)_{[i'-M,\infty)}$ because $\psi(L_{\underline{G}}(x)) = L_{\underline{H}}(\phi'(x))$. Hence $\tilde{\phi}(x)_i = \phi'(x)_i$, showing that $\tilde{\phi}$ extends $\phi'$. 

Finally, to establish the uniqueness assertion, assume that $\lambda :  X_{\underline{G}} \to X_{\underline{H}}$ is a sliding block code with the two properties stipulated in 1) and 2). Let $x \in X_{\underline{G}}$. Since $x$ is backward asymptotic to an element of $\overline{\Per X_{\underline{G}}}$ it follows from 1) that there is a $K \in \mathbb Z$ such $\lambda(x)_j = \tilde{\phi}(x)_j$ for $j \leq K$. Thanks to property 2) and because $(\underline{H},L_{\underline{H}})$ is right-resolving, it follows that $\lambda(x) = \tilde{\phi}(x)$. 
\end{proof}



We are now ready for \emph{the proof of Theorem \ref{18-01-26d}:} The existence of a sliding block code $\tilde{\phi}$ such that \eqref{18-01-26e} commutes follows from Lemma \ref{19-01-26}. Applying this conclusion to $\psi^{-1}$ we get a sliding block code
$$
\widetilde{{\phi'}^{-1}} : X_{\underline{H}} \to X_{\underline{G}}
$$
extending the conjugacy ${\phi'}^{-1} : \Per (X_{\underline{H}}) \to \Per (X_{\underline{G}})$ such that $L_{\underline{G}} \circ \widetilde{{\phi'}^{-1}} = \psi^{-1} \circ L_{\underline{H}}$. Since $\widetilde{{\phi'}^{-1}} \circ \tilde{\phi'} :  X_{\underline{G}} \to X_{\underline{G}}$ is a sliding block code extending the identity map on $\overline{\Per (X_{\underline{G}})}$ it follows from the uniqueness part of Lemma \ref{19-01-26} that $\widetilde{{\phi'}^{-1}} \circ \tilde{\phi'} = \id_{X_{\underline{G}}}$. By symmetry $ \tilde{\phi'}\circ \widetilde{{\phi'}^{-1}} =\id_{X_{\underline{H}}}$ and we conclude that $\tilde{\phi'}$ is a conjugacy.







To prove uniqueness assume that the conjugacy $\lambda : X_{\underline{G}} \to X_{\underline{H}}$ can substitute for $\tilde{\phi}$ in the diagram \eqref{18-01-26e}. Let $x \in X_{\underline{G}}$. If $x \in X_C$ for some source component $C$ in $\underline{G}$ we have that $\lambda(x) \in X_{C_1}$ and $\tilde{\phi}(x) \in X_{C_2}$ for some source components $C_1, C_2$ in $\underline{H}$ and $L_{\underline{H}}(\lambda(x)) = \psi(L_{\underline{G}}(x)) = L_{\underline{H}}(\tilde{\phi}(x))$. It follows from Corollary \ref{20-01-26h} that $\lambda(x) = \tilde{\phi}(x)$. Since $(\underline{H},L_{\underline{H}})$ is right-resolving it follows readily that the same holds when $x$ is backward asymptotic to a source component in $\underline{G}$. Since every element of $X_{\underline{G}}$ can be approximated by elements that are backward asymptotic to a source component in $\underline{G}$ it follows by continuity that $\lambda = \tilde{\phi}$. 


In order to obtain \eqref{20-01-25d} and thereby complete the proof, it remains only to show that $\psi_{\mathbb K} \circ f_{\underline{G}} = f_{\underline{H}} \circ \tilde{\phi}$. For this note that it follows from Lemma \ref{09-10-25fz} that $\psi_{\mathbb K} \circ f_{\underline{G}}(x) = f_{\underline{H}} \circ \tilde{\phi}(x)$ when $x \in \overline{\Per X_{\underline{G}}}$. Note that every element of $X_{\underline{G}}$ is backward asymptotic to an element of $ \overline{\Per X_{\underline{G}}}$. Since
$$
L_{\mathbb K} \circ \psi_{\mathbb K} \circ f_{\underline{G}} = \psi \circ L_{\underline{G}} = L_{\underline{H}} \circ \tilde{\phi} = L_{\mathbb K} \circ f_{\underline{H}} \circ \tilde{\phi}, 
$$
and $(\mathbb K(Z),L_{\mathbb K})$ is right-resolving, it follows from this that $\psi_{\mathbb K} \circ f_{\underline{G}} = f_{\underline{H}} \circ \tilde{\phi}$. \qed


\bigskip

For a sofic shift $Y$, set
$$
(\underline{\mathbb K}(Y),L_{\underline{\mathbb K}(Y)})  := (\underline{\mathbb K(Y)}, L_{\underline{\mathbb K(Y)}}).
$$
By Theorem \ref{Krieger} and Theorem \ref{18-01-26d} this defines a canonical cover for sofic shifts:


\begin{thm}\label{21-01-26} Let $\psi : Y \to Z$ be a conjugacy of sofic shifts.  There is a unique conjugacy $\psi_{\underline{\mathbb K}} : X_{\underline{\mathbb K}(Y)} \to X_{\underline{\mathbb K}(Z)}$ such that
\begin{equation*}
\xymatrix{
 X_{\underline{\mathbb K}(Y)} \ar[r]^-{\psi_{\underline{\mathbb K}}} \ar[d]_{L_{\underline{\mathbb K}(Y)}}& X_{\underline{\mathbb K}(Z)} \ar[d]^-{L_{\underline{\mathbb K}}(Z)}\\
   Y \ar[r]^{\psi} & Z}
\end{equation*}
commutes. This conjugacy has the property that also
\begin{equation}\label{20-01-25dx}
\xymatrix{
 X_{\underline{\mathbb K}(Y)} \ar[r]^-{\psi_{\underline{\mathbb K}}} \ar[d]_{f_{\underline{\mathbb K(Y)}}}& X_{\underline{\mathbb K}(Z)} \ar[d]^-{f_{\underline{\mathbb K(Z)}}}\\
 X_{\mathbb K(Y)} \ar[d]_-{L_{\mathbb K(Y)}} \ar[r]^-{\psi_{\mathbb K}} &  X_{\mathbb K(Z)} \ar[d]^-{L_{\mathbb K(Z)}} \\
Y \ar[r]_-{\psi}  & Z }
\end{equation}
commutes.
\end{thm}

It is not too difficult to find varying examples to show that $\underline{\mathbb K}$ in general does not give the same canonical extensions for mixing sofic shifts as any of the constructions $\mathbb F, \ \mathbb F', \ \mathbb K$ or $\mathbb K'$, and only slightly more difficult to find a single mixing sofic shift for which all these extensions are mutually different. However, the even shift $E$ we considered in Example \ref{evenshift} will not do it this time, because $(\underline{\mathbb K}(E), L_{\underline{\mathbb K}(E)}) = (\mathbb K(E),L_{\mathbb K(E)})$. We present an example below.

\begin{ex}\label{25-10-26}
The following labeled graph presents a mixing sofic shift $D$. Since the graph is strongly connected, right-resolving and follower-separated it is the minimal right-resolving presentation of $D$.

\begin{figure}[H]
\begin{equation*}
 \begin{tikzpicture}[node distance={35mm}, thick, main/.style = {draw, circle}] 
\node[main] (1) {$a$}; 
\node[main] (2) [right of=1] {$b$}; 
\node[main] (3) [right of=2] {$c$};
\draw[->] (1) to [out=130,in=80,looseness=7] node[left ,pos=0.5] {$1$} (1);
\draw[->] (1) to [out=60,in=130,looseness=0.8] node[above ,pos=0.5] {$2$} (2);
\draw[->] (2) to [out=180,in=350,looseness=0.8] node[below ,pos=0.5] {$3$} (1);
\draw[->] (2) to [out=100,in=70,looseness=10] node[above ,pos=0.5] {$1$} (2);
\draw[->] (2) to [out=60,in=30,looseness=10] node[above ,pos=0.5] {$2$} (2);
\draw[->] (2) to [out=10,in=170,looseness=0.8] node[above ,pos=0.5] {$0$} (3);
\draw[->] (3) to [out=190,in=350,looseness=0.8] node[below ,pos=0.5] {$0$} (2);
 \end{tikzpicture} 
\end{equation*}
\caption{The minimal right-resolving presentation $(\mathbb F(D),L_{\mathbb F(D)})$ of a mixing sofic shift $D$.}
\label{F}
\end{figure}

Since there are only three vertexes in the graph of Figure \ref{F} it is not difficult to find the canonical extensions we seek to compare by hand, and we just list the results:  
\begin{figure}[H]
\begin{equation*}\label{24-01-26e}
 \begin{tikzpicture}[node distance={35mm}, thick, main/.style = {draw, circle}] 
\node[main] (1) {$a$}; 
\node[main] (2) [right of=1] {$b$}; 
\node[main] (3) [right of=2] {$c$};
\node[main] (4) [below of=1] {$\{a,b\}$};
\node[main] (5) [below of=2] {$\{b,c\}$};
\draw[->] (1) to [out=130,in=80,looseness=7] node[left ,pos=0.5] {$1$} (1);
\draw[->] (1) to [out=60,in=130,looseness=0.8] node[above ,pos=0.5] {$2$} (2);
\draw[->] (2) to [out=180,in=350,looseness=0.8] node[below ,pos=0.5] {$3$} (1);
\draw[->] (2) to [out=100,in=70,looseness=10] node[above ,pos=0.5] {$1$} (2);
\draw[->] (2) to [out=60,in=30,looseness=10] node[above ,pos=0.5] {$2$} (2);
\draw[->] (2) to [out=10,in=170,looseness=0.8] node[above ,pos=0.5] {$0$} (3);
\draw[->] (3) to [out=190,in=350,looseness=0.8] node[below ,pos=0.5] {$0$} (2);
\draw[->] (4) to [out=160,in=200,looseness=5] node[left ,pos=0.5] {$1$} (4);
%\draw[->] (4) to [out=110,in=210,looseness=1] node[left ,pos=0.5] {$3$} (1);
\draw[->] (4) to [out=80,in=210,looseness=1] node[left ,pos=0.5] {$2$} (2);
%\draw[->] (4) to [out=310,in=0,looseness=2] node[right ,pos=0.5] {$0$} (3);
%\draw[->] (5) to [out=100,in=260,looseness=1] node[left ,pos=0.5] {$1$} (2);
%\draw[->] (5) to [out=70,in=300,looseness=1] node[left ,pos=0.5] {$2$} (2);
\draw[->] (5) to [out=350,in=310,looseness=5] node[below ,pos=0.5] {$0$} (5);
%\draw[->] (5) to [out=220,in=180,looseness=3] node[left ,pos=0.5] {$3$} (1);
 \end{tikzpicture} 
\end{equation*}
\caption{$(\mathbb F'(D),L_{\mathbb F'(D)})$.}
\end{figure}

\begin{figure}[H]
\begin{equation*}\label{24-01-26ex}
 \begin{tikzpicture}[node distance={35mm}, thick, main/.style = {draw, circle}] 
\node[main] (1) {$a$}; 
\node[main] (2) [right of=1] {$b$}; 
\node[main] (3) [right of=2] {$c$};
%\node[main] (4) [below of=1] {$\{a,b\}$};
\node[main] (5) [below of=2] {$\{b,c\}$};
\draw[->] (1) to [out=130,in=80,looseness=7] node[left ,pos=0.5] {$1$} (1);
\draw[->] (1) to [out=60,in=130,looseness=0.8] node[above ,pos=0.5] {$2$} (2);
\draw[->] (2) to [out=180,in=350,looseness=0.8] node[below ,pos=0.5] {$3$} (1);
\draw[->] (2) to [out=100,in=70,looseness=10] node[above ,pos=0.5] {$1$} (2);
\draw[->] (2) to [out=60,in=30,looseness=10] node[above ,pos=0.5] {$2$} (2);
\draw[->] (2) to [out=10,in=170,looseness=0.8] node[above ,pos=0.5] {$0$} (3);
\draw[->] (3) to [out=190,in=350,looseness=0.8] node[below ,pos=0.5] {$0$} (2);
%\draw[->] (4) to [out=160,in=200,looseness=5] node[left ,pos=0.5] {$1$} (4);
%\draw[->] (4) to [out=110,in=210,looseness=1] node[left ,pos=0.5] {$3$} (1);
%\draw[->] (4) to [out=80,in=210,looseness=1] node[left ,pos=0.5] {$2$} (2);
%\draw[->] (4) to [out=310,in=0,looseness=2] node[right ,pos=0.5] {$0$} (3);
\draw[->] (5) to [out=100,in=260,looseness=1] node[left ,pos=0.5] {$1$} (2);
\draw[->] (5) to [out=70,in=300,looseness=1] node[left ,pos=0.5] {$2$} (2);
\draw[->] (5) to [out=350,in=310,looseness=5] node[below ,pos=0.5] {$0$} (5);
\draw[->] (5) to [out=130,in=270,looseness=0.5] node[left ,pos=0.5] {$3$} (1);
 \end{tikzpicture} 
\end{equation*}
\caption{$(\mathbb K(D),L_{\mathbb K(D)})$.}
\end{figure}

\begin{figure}[H]
\begin{equation*}
 \begin{tikzpicture}[node distance={35mm}, thick, main/.style = {draw, circle}] 
\node[main] (1) at (-6,3.5) {$a$}; 
\node[main] (2) at (-4,3.5) {$b$}; 
\node[main] (3) at (-2,3.5) {$c$};
\node[main] (4) at (-7,0) {$\{a,b\}$};
\node[main] (6) at (-1,0) {$\{b,d\}$};
\node[main] (7) at (2,0) {$\{c,d\}$};
\node[main] (8) at (-4,6) {$\{c,b,d\}$};
\draw[->] (1) to [out=175,in=125,looseness=7] node[left ,pos=0.5] {$1$} (1);
\draw[->] (1) to [out=60,in=150,looseness=0.8] node[above ,pos=0.5] {$2$} (2);
\draw[->] (2) to [out=180,in=350,looseness=0.8] node[above ,pos=0.5] {$3$} (1);
\draw[->] (2) to [out=90,in=60,looseness=10] node[above ,pos=0.5] {$1$} (2);
\draw[->] (2) to [out=50,in=20,looseness=10] node[above ,pos=0.5] {$2$} (2);
\draw[->] (2) to [out=10,in=170,looseness=0.8] node[above ,pos=0.5] {$0$} (3);
\draw[->] (3) to [out=190,in=350,looseness=0.8] node[below ,pos=0.5] {$0$} (2);
\draw[->] (4) to [out=250,in=300,looseness=6] node[above ,pos=0.5] {$1$} (4);
\draw[->] (6) to [out=40,in=150,looseness=1] node[above ,pos=0.5] {$0$} (7);
\draw[->] (7) to [out=190,in=350,looseness=1] node[below ,pos=0.5] {$0$} (6);
\draw[->] (4) to [out=70,in=190,looseness=1] node[left ,pos=0.3] {$2$} (2);
\draw[->] (6) to [out=130,in=270,looseness=1] node[left ,pos=0.4] {$1$} (2);
\draw[->] (6) to [out=110,in=290,looseness=1] node[left ,pos=0.4] {$2$} (2);
\draw[->] (6) to [out=230,in=200,looseness=3] node[left ,pos=0.7] {$3$} (1);
\draw[->] (8) to [out=0,in=40,looseness=4] node[right ,pos=0.4] {$0$} (8);
\end{tikzpicture} 
\end{equation*}
\caption{$(\mathbb K'(D),L_{\mathbb K'(D)})$.}
\label{11-09-25c}
\end{figure}

\begin{figure}[H]
\begin{equation*}\label{24-01-26ey}
 \begin{tikzpicture}[node distance={35mm}, thick, main/.style = {draw, circle}] 
\node[main] (1) {$a$}; 
\node[main] (2) [right of=1] {$b$}; 
\node[main] (3) [right of=2] {$c$};
\node[main] (4) [below of=1] {$\{a,b\}$};
\node[main] (5) [below of=2] {$\{b,c\}$};
\draw[->] (1) to [out=130,in=80,looseness=7] node[left ,pos=0.5] {$1$} (1);
\draw[->] (1) to [out=60,in=130,looseness=0.8] node[above ,pos=0.5] {$2$} (2);
\draw[->] (2) to [out=180,in=350,looseness=0.8] node[below ,pos=0.5] {$3$} (1);
\draw[->] (2) to [out=100,in=70,looseness=10] node[above ,pos=0.5] {$1$} (2);
\draw[->] (2) to [out=60,in=30,looseness=10] node[above ,pos=0.5] {$2$} (2);
\draw[->] (2) to [out=10,in=170,looseness=0.8] node[above ,pos=0.5] {$0$} (3);
\draw[->] (3) to [out=190,in=350,looseness=0.8] node[below ,pos=0.5] {$0$} (2);
\draw[->] (4) to [out=160,in=200,looseness=5] node[left ,pos=0.5] {$1$} (4);
\draw[->] (4) to [out=110,in=210,looseness=1] node[left ,pos=0.5] {$3$} (1);
\draw[->] (4) to [out=80,in=210,looseness=1] node[left ,pos=0.5] {$2$} (2);
\draw[->] (4) to [out=310,in=0,looseness=1.9] node[right ,pos=0.5] {$0$} (3);
\draw[->] (5) to [out=100,in=260,looseness=1] node[left ,pos=0.5] {$1$} (2);
\draw[->] (5) to [out=70,in=300,looseness=1] node[left ,pos=0.5] {$2$} (2);
\draw[->] (5) to [out=350,in=310,looseness=5] node[below ,pos=0.5] {$0$} (5);
\draw[->] (5) to [out=220,in=180,looseness=3] node[left ,pos=0.5] {$3$} (1);
 \end{tikzpicture} 
\end{equation*}
\caption{$(\underline{\mathbb K}(D),L_{\underline{\mathbb K}(D)})$.}\label{underF}
\end{figure}

We have not mentioned the extension $(\underline{\mathbb F}(D),L_{\underline{\mathbb F}(D)})$, and for good reasons: It is the same as $(\underline{\mathbb K}(D),L_{\underline{\mathbb K}(D)})$ and we conjecture that this is so for all irreducible sofic shifts.

\end{ex}







\end{document}


\begin{ex}\label{25-10-26}
The following labeled graph presents a mixing sofic shift $D$. Since the graph is strongly connected, right-resolving and follower-separated it is the minimal right-resolving presentation $D$.

\begin{figure}[H]
\begin{equation*}\label{24-01-26d}
 \begin{tikzpicture}[node distance={35mm}, thick, main/.style = {draw, circle}] 
\node[main] (1) {$a$}; 
\node[main] (2) [right of=1] {$b$}; 
\node[main] (3) [right of=2] {$c$};
\draw[->] (1) to [out=130,in=80,looseness=7] node[left ,pos=0.5] {$1$} (1);
\draw[->] (1) to [out=60,in=130,looseness=0.8] node[above ,pos=0.5] {$2$} (2);
\draw[->] (2) to [out=180,in=350,looseness=0.8] node[below ,pos=0.5] {$3$} (1);
\draw[->] (2) to [out=100,in=70,looseness=10] node[above ,pos=0.5] {$1$} (2);
\draw[->] (2) to [out=60,in=30,looseness=10] node[above ,pos=0.5] {$2$} (2);
\draw[->] (2) to [out=10,in=170,looseness=0.8] node[above ,pos=0.5] {$0$} (3);
\draw[->] (3) to [out=190,in=350,looseness=0.8] node[below ,pos=0.5] {$0$} (2);
 \end{tikzpicture} 
\end{equation*}
\caption{The minimal right-resolving presentation $(\mathbb F(D),L_{\mathbb F(D)})$ of a mixing sofic shift $D$.}
\label{F}
\end{figure}

Since there are only three vertexes in the graph of Figure \ref{F} it is not difficult to find $(\underline{F}(D),L_{\underline{\mathbb F}(D)})$ and $(\mathbb F'(D),L_{\mathbb F'(D)})$. The results are 
\begin{figure}[H]
\begin{equation*}\label{24-01-26e}
 \begin{tikzpicture}[node distance={35mm}, thick, main/.style = {draw, circle}] 
\node[main] (1) {$a$}; 
\node[main] (2) [right of=1] {$b$}; 
\node[main] (3) [right of=2] {$c$};
\node[main] (4) [below of=1] {$\{a,b\}$};
\node[main] (5) [below of=2] {$\{b,c\}$};
\draw[->] (1) to [out=130,in=80,looseness=7] node[left ,pos=0.5] {$1$} (1);
\draw[->] (1) to [out=60,in=130,looseness=0.8] node[above ,pos=0.5] {$2$} (2);
\draw[->] (2) to [out=180,in=350,looseness=0.8] node[below ,pos=0.5] {$3$} (1);
\draw[->] (2) to [out=100,in=70,looseness=10] node[above ,pos=0.5] {$1$} (2);
\draw[->] (2) to [out=60,in=30,looseness=10] node[above ,pos=0.5] {$2$} (2);
\draw[->] (2) to [out=10,in=170,looseness=0.8] node[above ,pos=0.5] {$0$} (3);
\draw[->] (3) to [out=190,in=350,looseness=0.8] node[below ,pos=0.5] {$0$} (2);
\draw[->] (4) to [out=160,in=200,looseness=5] node[left ,pos=0.5] {$1$} (4);
\draw[->] (4) to [out=110,in=210,looseness=1] node[left ,pos=0.5] {$3$} (1);
\draw[->] (4) to [out=80,in=210,looseness=1] node[left ,pos=0.5] {$2$} (2);
\draw[->] (4) to [out=310,in=0,looseness=1.9] node[right ,pos=0.5] {$0$} (3);
\draw[->] (5) to [out=100,in=260,looseness=1] node[left ,pos=0.5] {$1$} (2);
\draw[->] (5) to [out=70,in=300,looseness=1] node[left ,pos=0.5] {$2$} (2);
\draw[->] (5) to [out=350,in=310,looseness=5] node[below ,pos=0.5] {$0$} (5);
\draw[->] (5) to [out=220,in=180,looseness=3] node[left ,pos=0.5] {$3$} (1);
 \end{tikzpicture} 
\end{equation*}
\caption{$(\underline{\mathbb F}(D),L_{\underline{\mathbb F}(D)})$.}\label{underF}
\end{figure}

and
\begin{figure}[H]
\begin{equation*}\label{24-01-26e}
 \begin{tikzpicture}[node distance={35mm}, thick, main/.style = {draw, circle}] 
\node[main] (1) {$a$}; 
\node[main] (2) [right of=1] {$b$}; 
\node[main] (3) [right of=2] {$c$};
\node[main] (4) [below of=1] {$\{a,b\}$};
\node[main] (5) [below of=2] {$\{b,c\}$};
\draw[->] (1) to [out=130,in=80,looseness=7] node[left ,pos=0.5] {$1$} (1);
\draw[->] (1) to [out=60,in=130,looseness=0.8] node[above ,pos=0.5] {$2$} (2);
\draw[->] (2) to [out=180,in=350,looseness=0.8] node[below ,pos=0.5] {$3$} (1);
\draw[->] (2) to [out=100,in=70,looseness=10] node[above ,pos=0.5] {$1$} (2);
\draw[->] (2) to [out=60,in=30,looseness=10] node[above ,pos=0.5] {$2$} (2);
\draw[->] (2) to [out=10,in=170,looseness=0.8] node[above ,pos=0.5] {$0$} (3);
\draw[->] (3) to [out=190,in=350,looseness=0.8] node[below ,pos=0.5] {$0$} (2);
\draw[->] (4) to [out=160,in=200,looseness=5] node[left ,pos=0.5] {$1$} (4);
%\draw[->] (4) to [out=110,in=210,looseness=1] node[left ,pos=0.5] {$3$} (1);
\draw[->] (4) to [out=80,in=210,looseness=1] node[left ,pos=0.5] {$2$} (2);
%\draw[->] (4) to [out=310,in=0,looseness=2] node[right ,pos=0.5] {$0$} (3);
%\draw[->] (5) to [out=100,in=260,looseness=1] node[left ,pos=0.5] {$1$} (2);
%\draw[->] (5) to [out=70,in=300,looseness=1] node[left ,pos=0.5] {$2$} (2);
\draw[->] (5) to [out=350,in=310,looseness=5] node[below ,pos=0.5] {$0$} (5);
%\draw[->] (5) to [out=220,in=180,looseness=3] node[left ,pos=0.5] {$3$} (1);
 \end{tikzpicture} 
\end{equation*}
\caption{$(\mathbb F'(D),L_{\mathbb F'(D)})$.}
\end{figure}
The graph $(\mathbb K(D),L_{\mathbb K(D)})$ is the merged graph of that in Figure \ref{underF}: 



\begin{figure}[H]
\begin{equation*}\label{24-01-26e}
 \begin{tikzpicture}[node distance={35mm}, thick, main/.style = {draw, circle}] 
\node[main] (1) {$a$}; 
\node[main] (2) [right of=1] {$b$}; 
\node[main] (3) [right of=2] {$c$};
%\node[main] (4) [below of=1] {$\{a,b\}$};
\node[main] (5) [below of=2] {$\{b,c\}$};
\draw[->] (1) to [out=130,in=80,looseness=7] node[left ,pos=0.5] {$1$} (1);
\draw[->] (1) to [out=60,in=130,looseness=0.8] node[above ,pos=0.5] {$2$} (2);
\draw[->] (2) to [out=180,in=350,looseness=0.8] node[below ,pos=0.5] {$3$} (1);
\draw[->] (2) to [out=100,in=70,looseness=10] node[above ,pos=0.5] {$1$} (2);
\draw[->] (2) to [out=60,in=30,looseness=10] node[above ,pos=0.5] {$2$} (2);
\draw[->] (2) to [out=10,in=170,looseness=0.8] node[above ,pos=0.5] {$0$} (3);
\draw[->] (3) to [out=190,in=350,looseness=0.8] node[below ,pos=0.5] {$0$} (2);
%\draw[->] (4) to [out=160,in=200,looseness=5] node[left ,pos=0.5] {$1$} (4);
%\draw[->] (4) to [out=110,in=210,looseness=1] node[left ,pos=0.5] {$3$} (1);
%\draw[->] (4) to [out=80,in=210,looseness=1] node[left ,pos=0.5] {$2$} (2);
%\draw[->] (4) to [out=310,in=0,looseness=2] node[right ,pos=0.5] {$0$} (3);
\draw[->] (5) to [out=100,in=260,looseness=1] node[left ,pos=0.5] {$1$} (2);
\draw[->] (5) to [out=70,in=300,looseness=1] node[left ,pos=0.5] {$2$} (2);
\draw[->] (5) to [out=350,in=310,looseness=5] node[below ,pos=0.5] {$0$} (5);
\draw[->] (5) to [out=130,in=270,looseness=0.5] node[left ,pos=0.5] {$3$} (1);
 \end{tikzpicture} 
\end{equation*}
\caption{$(\mathbb K(D),L_{\mathbb K(D)})$.}
\end{figure}
The presence of the additional vertex $\{b,c\}$ which we now rename to $d$ makes the size of the graphs to be considered considerably. The largest is $(\overline{\mathbb K(D)},L_{\overline{\mathbb K}(D)})$: 


\begin{figure}[H]
\begin{equation*}
 \begin{tikzpicture}[node distance={35mm}, thick, main/.style = {draw, circle}] 
\node[main] (1) at (-6,3.5) {$a$}; 
\node[main] (2) at (-4,3.5) {$b$}; 
\node[main] (3) at (-2,3.5) {$c$};
\node[main] (4) at (-7,0) {$\{a,b\}$};
\node[main] (5) at (-4,0) {$\{b,c\}$};
\node[main] (6) at (-1,0) {$\{b,d\}$};
\node[main] (7) at (2,0) {$\{c,d\}$};
\node[main] (8) at (-4,8) {$\{c,b,d\}$};
\draw[->] (1) to [out=175,in=125,looseness=7] node[left ,pos=0.5] {$1$} (1);
\draw[->] (1) to [out=60,in=150,looseness=0.8] node[above ,pos=0.5] {$2$} (2);
\draw[->] (2) to [out=180,in=350,looseness=0.8] node[above ,pos=0.5] {$3$} (1);
\draw[->] (2) to [out=90,in=60,looseness=10] node[above ,pos=0.5] {$1$} (2);
\draw[->] (2) to [out=50,in=20,looseness=10] node[above ,pos=0.5] {$2$} (2);
\draw[->] (2) to [out=10,in=170,looseness=0.8] node[above ,pos=0.5] {$0$} (3);
\draw[->] (3) to [out=190,in=350,looseness=0.8] node[below ,pos=0.5] {$0$} (2);
\draw[->] (4) to [out=250,in=300,looseness=6] node[above ,pos=0.5] {$1$} (4);
\draw[->] (5) to [out=250,in=300,looseness=6] node[above ,pos=0.5] {$0$} (5);
\draw[->] (6) to [out=40,in=150,looseness=1] node[above ,pos=0.5] {$0$} (7);
\draw[->] (7) to [out=190,in=350,looseness=1] node[below ,pos=0.5] {$0$} (6);
\draw[->] (4) to [out=100,in=260,looseness=0.5] node[left ,pos=0.3] {$3$} (1);
\draw[->] (4) to [out=70,in=190,looseness=1] node[left ,pos=0.3] {$2$} (2);
\draw[->] (4) to [out=50,in=230,looseness=1] node[left ,pos=0.2] {$0$} (3);
\draw[->] (5) to [out=230,in=230,looseness=3] node[left ,pos=0.5] {$3$} (1);
\draw[->] (5) to [out=120,in=230,looseness=1] node[left ,pos=0.2] {$1$} (2);
\draw[->] (5) to [out=100,in=250,looseness=1] node[left ,pos=0.2] {$2$} (2);
\draw[->] (6) to [out=130,in=270,looseness=1] node[left ,pos=0.4] {$1$} (2);
\draw[->] (6) to [out=110,in=290,looseness=1] node[left ,pos=0.4] {$2$} (2);
\draw[->] (6) to [out=230,in=200,looseness=3] node[left ,pos=0.7] {$3$} (1);
\draw[->] (7) to [out=130,in=310,looseness=1] node[right ,pos=0.4] {$1$} (2);
\draw[->] (7) to [out=100,in=330,looseness=1] node[right ,pos=0.4] {$2$} (2);
\draw[->] (7) to [out=80,in=80,looseness=1] node[right ,pos=0.4] {$3$} (1);
\draw[->] (8) to [out=0,in=40,looseness=4] node[right ,pos=0.4] {$0$} (8);
\draw[->] (8) to [out=220,in=100,looseness=1] node[left ,pos=0.5] {$3$} (1);
\draw[->] (8) to [out=250,in=130,looseness=1] node[left ,pos=0.3] {$1$} (2);
\draw[->] (8) to [out=260,in=110,looseness=1] node[right ,pos=0.3] {$2$} (2);
\end{tikzpicture} 
\end{equation*}
\caption{A presentation of $(\overline{\mathbb K(D)},L_{\overline{\mathbb K(D)}})$.}
\label{11-09-25c}
\end{figure}

The graphs $(\mathbb K(D)'',L_{\mathbb K(D)''}), \ (\mathbb K'(D),L_{\mathbb K'(D)})$ and $(\underline{\mathbb K}(D),L_{\underline{K}(D)})$ are all subgraph of the graph above.


\begin{figure}[H]
\begin{equation*}
 \begin{tikzpicture}[node distance={35mm}, thick, main/.style = {draw, circle}] 
\node[main] (1) at (-6,3.5) {$a$}; 
\node[main] (2) at (-4,3.5) {$b$}; 
\node[main] (3) at (-2,3.5) {$c$};
\node[main] (4) at (-7,0) {$\{a,b\}$};
\node[main] (5) at (-4,0) {$\{b,c\}$};
\node[main] (6) at (-1,0) {$\{b,d\}$};
\node[main] (7) at (2,0) {$\{c,d\}$};
\node[main] (8) at (-4,7) {$\{c,b,d\}$};
\draw[->] (1) to [out=175,in=125,looseness=7] node[left ,pos=0.5] {$1$} (1);
\draw[->] (1) to [out=60,in=150,looseness=0.8] node[above ,pos=0.5] {$2$} (2);
\draw[->] (2) to [out=180,in=350,looseness=0.8] node[above ,pos=0.5] {$3$} (1);
\draw[->] (2) to [out=90,in=60,looseness=10] node[above ,pos=0.5] {$1$} (2);
\draw[->] (2) to [out=50,in=20,looseness=10] node[above ,pos=0.5] {$2$} (2);
\draw[->] (2) to [out=10,in=170,looseness=0.8] node[above ,pos=0.5] {$0$} (3);
\draw[->] (3) to [out=190,in=350,looseness=0.8] node[below ,pos=0.5] {$0$} (2);
\draw[->] (4) to [out=250,in=300,looseness=6] node[above ,pos=0.5] {$1$} (4);
\draw[->] (6) to [out=40,in=150,looseness=1] node[above ,pos=0.5] {$0$} (7);
\draw[->] (7) to [out=190,in=350,looseness=1] node[below ,pos=0.5] {$0$} (6);
\draw[->] (4) to [out=70,in=190,looseness=1] node[left ,pos=0.3] {$2$} (2);
\draw[->] (6) to [out=130,in=270,looseness=1] node[left ,pos=0.4] {$1$} (2);
\draw[->] (6) to [out=110,in=290,looseness=1] node[left ,pos=0.4] {$2$} (2);
\draw[->] (6) to [out=230,in=200,looseness=3] node[left ,pos=0.7] {$3$} (1);
\draw[->] (8) to [out=0,in=40,looseness=4] node[right ,pos=0.4] {$0$} (8);
\draw[->] (5) to [out=250,in=300,looseness=6] node[above ,pos=0.5] {$0$} (5);
\end{tikzpicture} 
\end{equation*}
\caption{$(\mathbb K(D)'',L_{\mathbb K''(D)})$.}
\label{11-09-25c}
\end{figure}

While this labeled graph presents a canonical extension of $D$ by Corollary \ref{17-11-25}, this example illustrates it does in general give a strongly canonical extension: The automorphism $\gamma$ of $X_{\mathbb K(D)''}$ obtained by permuting the vertexes $\{b,c,d\}$ and $\{b,c\}$ projects to the identity map on $D$ via the labeling.

The graphs $(\mathbb K'(D),L_{\mathbb K'(D)})$ and $(\underline{K}(D), L_{\underline{K}(D)})$ are


\newpage


\begin{figure}[H]
\begin{equation*}
 \begin{tikzpicture}[node distance={35mm}, thick, main/.style = {draw, circle}] 
\node[main] (1) at (-6,3.5) {$a$}; 
\node[main] (2) at (-4,3.5) {$b$}; 
\node[main] (3) at (-2,3.5) {$c$};
\node[main] (4) at (-7,0) {$\{a,b\}$};
\node[main] (6) at (-1,0) {$\{b,d\}$};
\node[main] (7) at (2,0) {$\{c,d\}$};
\node[main] (8) at (-4,7) {$\{c,b,d\}$};
\draw[->] (1) to [out=175,in=125,looseness=7] node[left ,pos=0.5] {$1$} (1);
\draw[->] (1) to [out=60,in=150,looseness=0.8] node[above ,pos=0.5] {$2$} (2);
\draw[->] (2) to [out=180,in=350,looseness=0.8] node[above ,pos=0.5] {$3$} (1);
\draw[->] (2) to [out=90,in=60,looseness=10] node[above ,pos=0.5] {$1$} (2);
\draw[->] (2) to [out=50,in=20,looseness=10] node[above ,pos=0.5] {$2$} (2);
\draw[->] (2) to [out=10,in=170,looseness=0.8] node[above ,pos=0.5] {$0$} (3);
\draw[->] (3) to [out=190,in=350,looseness=0.8] node[below ,pos=0.5] {$0$} (2);
\draw[->] (4) to [out=250,in=300,looseness=6] node[above ,pos=0.5] {$1$} (4);
\draw[->] (6) to [out=40,in=150,looseness=1] node[above ,pos=0.5] {$0$} (7);
\draw[->] (7) to [out=190,in=350,looseness=1] node[below ,pos=0.5] {$0$} (6);
\draw[->] (4) to [out=70,in=190,looseness=1] node[left ,pos=0.3] {$2$} (2);
\draw[->] (6) to [out=130,in=270,looseness=1] node[left ,pos=0.4] {$1$} (2);
\draw[->] (6) to [out=110,in=290,looseness=1] node[left ,pos=0.4] {$2$} (2);
\draw[->] (6) to [out=230,in=200,looseness=3] node[left ,pos=0.7] {$3$} (1);
\draw[->] (8) to [out=0,in=40,looseness=4] node[right ,pos=0.4] {$0$} (8);
\end{tikzpicture} 
\end{equation*}
\caption{$(\mathbb K'(D),L_{\mathbb K'(D)})$.}
\label{11-09-25c}
\end{figure}

\newpage

\begin{figure}[H]
\begin{equation*}
 \begin{tikzpicture}[node distance={35mm}, thick, main/.style = {draw, circle}] 
\node[main] (1) at (-6,3.5) {$a$}; 
\node[main] (2) at (-4,3.5) {$b$}; 
\node[main] (3) at (-2,3.5) {$c$};
\node[main] (4) at (-7,0) {$\{a,b\}$};
\node[main] (8) at (-4,8) {$\{c,b,d\}$};
\draw[->] (1) to [out=175,in=125,looseness=7] node[left ,pos=0.5] {$1$} (1);
\draw[->] (1) to [out=60,in=150,looseness=0.8] node[above ,pos=0.5] {$2$} (2);
\draw[->] (2) to [out=180,in=350,looseness=0.8] node[above ,pos=0.5] {$3$} (1);
\draw[->] (2) to [out=90,in=60,looseness=10] node[above ,pos=0.5] {$1$} (2);
\draw[->] (2) to [out=50,in=20,looseness=10] node[above ,pos=0.5] {$2$} (2);
\draw[->] (2) to [out=10,in=170,looseness=0.8] node[above ,pos=0.5] {$0$} (3);
\draw[->] (3) to [out=190,in=350,looseness=0.8] node[below ,pos=0.5] {$0$} (2);
\draw[->] (4) to [out=250,in=300,looseness=6] node[above ,pos=0.5] {$1$} (4);
\draw[->] (4) to [out=100,in=260,looseness=0.5] node[left ,pos=0.3] {$3$} (1);
\draw[->] (4) to [out=70,in=190,looseness=1] node[left ,pos=0.3] {$2$} (2);
\draw[->] (4) to [out=50,in=230,looseness=1] node[left ,pos=0.2] {$0$} (3);
\draw[->] (8) to [out=0,in=40,looseness=4] node[right ,pos=0.4] {$0$} (8);
\draw[->] (8) to [out=220,in=100,looseness=1] node[left ,pos=0.5] {$3$} (1);
\draw[->] (8) to [out=250,in=130,looseness=1] node[left ,pos=0.3] {$1$} (2);
\draw[->] (8) to [out=260,in=110,looseness=1] node[right ,pos=0.3] {$2$} (2);
\end{tikzpicture} 
\end{equation*}
\caption{$(\underline{\mathbb K}(D),L_{\underline{\mathbb K}(D)})$.}
\label{11-09-25c}
\end{figure}
\end{ex}

For a sofic shift $Y$ presented by a labeled graph with few vertexes, it is relatively easy to find labeled graphs that depict the canonical covers $\mathbb K(Y), \mathbb K'(Y)$ and $\underline{\mathbb K}(Y)$. When the graph is strongly connected so that $Y$ is irreducible it is also relatively easy to describe $\mathbb F(Y), \mathbb F'(Y), \underline{\mathbb F}(Y)$. The graphs in the Example \ref{evenshift} and Example \ref{02-11-25} have been chosen because they are economical with respect to space, but suffice to show that in general the graphs $\mathbb K(Y), \mathbb K'(Y), \underline{\mathbb K}(Y), \mathbb F(Y)$ and $\mathbb F'(Y)$ are different for mixing sofic shifts. Concerning $\underline{\mathbb F}(Y)$ we conjecture that $(\underline{\mathbb F}(Y), L_{\underline{\mathbb F}(Y)}) =  (\underline{\mathbb K}(Y), L_{\underline{\mathbb K}(Y)})$ when $Y$ is irreducible. 


\begin{lemma}\label{24-01-26} $(\underline{G}, L_{\underline{G}})$ is regular and $\alpha_G(Y) \subseteq \mathcal R\left(L_{\underline{G}}\right)$.
\end{lemma}
\begin{proof} Let $y \in Y$ and note that $D^y = s_{\underline{G}}(\alpha_G(y)_0)$. It follows therefore from Lemma \ref{12-09-25b} that it suffices here to show that $\alpha_G(y) \in X_{\underline{G}}$ is regular; i.e. it suffices to show that 
$$
F(\sigma^{k+1}(y)) = F\left(\sigma^{k+1}(L_{\underline{G}}(\alpha_G(y)))\right) \subseteq f_{\underline{G}}(t_{\underline{G}}(\alpha_G(y)_k))
$$
for $k \in \mathbb Z$. Let $ z \in F(\sigma^{k+1}(y))$. Define $w \in Y$ such that
$$
w_i = \begin{cases} y_i,\ i \leq k \\ z_{i-k-1}, \ i \geq k+1. \end{cases}
$$
Then $\alpha_G(w) \in X_{\underline{G}}$ and $L_{\underline{G}}(\alpha_G(w)_{[k+1,\infty)}) = z$. Hence $z \in f_{\underline{G}}(t_{\underline{G}}(\alpha_G(y)_k))$.
\end{proof}

\begin{lemma}\label{24-01-26a} The diagram
\begin{equation*}
\xymatrix{
 X_{\underline{G}} \ar[r]^-{\tilde{\phi}} &  X_{\underline{H}} \\
Y \ar[r]_-\psi \ar[u]^-{\alpha_G} & Z \ar[u]_-{\alpha_{H}} }
\end{equation*}
commutes.
\end{lemma}
\begin{proof} Let $y \in Y$. Since $\tilde{\phi}$ is a conjugacy and \eqref{18-01-26e} commutes, $\tilde{\phi}(\mathcal R(L_{\underline{G}})) =\mathcal R(L_{\underline{H}})$. Hence $\alpha_H(y)$ and $\tilde{\phi}(\alpha_G(y))$ are both elements of $\mathcal R(L_{\underline{H}})$ by Lemma \ref{24-01-26}. Since $L_{\underline{H}}(\tilde{\phi}(\alpha_G(y))) = L_{\underline{G}}(\alpha_G(y)) = y =  L_{\underline{H}}(\alpha_H(y))$, it follows from?????
\end{proof}

\begin{ex}\label{02-11-25}


The following labeled graph is irreducible, right-resolving and follower-separated and hence the minimal right resolving graph $(\mathbb F(Y),L_{\mathbb F(Y)})$ of the sofic shift $Y$ it represents. Note that $Y$ is mixing.

\begin{equation}\label{02-11-25d}
 \begin{tikzpicture}[node distance={35mm}, thick, main/.style = {draw, circle}] 
\node[main] (1) {$a$}; 
\node[main] (2) [right of=1] {$b$}; 
\draw[->] (1) to [out=170,in=200,looseness=20] node[left ,pos=0.5] {$1$} (1);
\draw[->] (1) to [out=250,in=290,looseness=20] node[below ,pos=0.5] {$2$} (1);
\draw[->] (1) to [out=300,in=200,looseness=1] node[above ,pos=0.5] {$0$} (2);
\draw[->] (2) to [out=120,in=70,looseness=1] node[above ,pos=0.5] {$1$} (1);
\draw[->] (2) to [out=250,in=290,looseness=20] node[below ,pos=0.5] {$2$} (2);
 \end{tikzpicture} 
\end{equation}
The graph $(\underline{\mathbb F}(Y),L_{\underline{\mathbb F}(Y)})$ is the  following.



\begin{equation*}\label{02-11-24e}
 \begin{tikzpicture}[node distance={35mm}, thick, main/.style = {draw, circle}] 
\node[main] (1) {$\{a\}$}; 
\node[main] (2) [right of=1] {$\{b\}$};
\node[main] (3) [below of=1] {$\{a,b\}$}; 
\draw[->] (1) to [out=170,in=200,looseness=20] node[left ,pos=0.5] {$1$} (1);
\draw[->] (1) to [out=250,in=290,looseness=10] node[below ,pos=0.5] {$2$} (1);
\draw[->] (1) to [out=300,in=200,looseness=1] node[above ,pos=0.5] {$0$} (2);
\draw[->] (2) to [out=120,in=70,looseness=1] node[above ,pos=0.5] {$1$} (1);
\draw[->] (2) to [out=250,in=290,looseness=10] node[below ,pos=0.5] {$2$} (2);
\draw[->] (3) to [out=250,in=290,looseness=10] node[below ,pos=0.5] {$2$} (3);
\draw[->] (3) to [out=110,in=220,looseness=1] node[left ,pos=0.5] {$1$} (1);
\draw[->] (3) to [out=50,in=220,looseness=1] node[left ,pos=0.5] {$0$} (2);
 \end{tikzpicture} 
\end{equation*}


Here the vertex $\{a,b\}$ has the same follower set as $\{a\}$ and the merged graph, which is the future cover $(\mathbb K(Y),L_{\mathbb K(Y)})$ by Corollary \ref{04-09-25d}, is the graph \eqref{02-11-25d} we started with. Thus $(\mathbb K(Y),L_{\mathbb K(Y)}) = (\mathbb F(Y),L_{\mathbb F(Y)})$ and $(\underline{\mathbb K}(Y),L_{\underline{\mathbb K}(Y)}) = (\underline{\mathbb F}(Y),L_{\underline{\mathbb F}(Y)})$ in this case. The labeled graph $(\mathbb K'(Y),L_{\mathbb K'(Y)})$ is smaller than $(\underline{\mathbb K}(Y),L_{\underline{\mathbb K}(Y)})$. It is obtained from the graph depicted above by removing the arrow from $\{a,b\}$ to $\{b\}$.
\end{ex} 

\begin{figure}
\begin{equation*}\label{24-01-26e}
 \begin{tikzpicture}[node distance={35mm}, thick, main/.style = {draw, circle}] 
\node[main] (1) {$a$}; 
\node[main] (2) [right of=1] {$b$}; 
\node[main] (3) [right of=2] {$c$};
\node[main] (4) [below of=1] {$\{a,b\}$};
\node[main] (5) [below of=2] {$\{b,d\}$};
\node[main] (6) [below of=3] {$\{c,d\}$};
\draw[->] (1) to [out=130,in=80,looseness=7] node[left ,pos=0.5] {$1$} (1);
\draw[->] (1) to [out=60,in=130,looseness=0.8] node[above ,pos=0.5] {$2$} (2);
\draw[->] (2) to [out=180,in=350,looseness=0.8] node[below ,pos=0.5] {$3$} (1);
\draw[->] (2) to [out=100,in=70,looseness=10] node[above ,pos=0.5] {$1$} (2);
\draw[->] (2) to [out=60,in=30,looseness=10] node[above ,pos=0.5] {$2$} (2);
\draw[->] (2) to [out=10,in=170,looseness=0.8] node[above ,pos=0.5] {$0$} (3);
\draw[->] (3) to [out=190,in=350,looseness=0.8] node[below ,pos=0.5] {$0$} (2);
\draw[->] (4) to [out=160,in=200,looseness=5] node[left ,pos=0.5] {$1$} (4);
\draw[->] (4) to [out=110,in=210,looseness=1] node[left ,pos=0.5] {$3$} (1);
\draw[->] (4) to [out=80,in=210,looseness=1] node[left ,pos=0.5] {$2$} (2);
\draw[->] (4) to [out=300,in=340,looseness=1.9] node[right ,pos=0.5] {$0$} (3);
\draw[->] (5) to [out=100,in=240,looseness=1] node[left ,pos=0.5] {$1$} (2);
\draw[->] (5) to [out=70,in=275,looseness=1] node[left ,pos=0.5] {$2$} (2);
\draw[->] (5) to [out=20,in=160,looseness=1] node[above ,pos=0.5] {$0$} (6);
\draw[->] (5) to [out=220,in=160,looseness=3] node[left ,pos=0.5] {$3$} (1);
\draw[->] (6) to [out=200,in=350,looseness=1] node[below ,pos=0.5] {$0$} (5);
\draw[->] (6) to [out=130,in=300,looseness=0.7] node[right ,pos=0.5] {$1$} (2);
\draw[->] (6) to [out=100,in=320,looseness=0.7] node[right ,pos=0.5] {$2$} (2);
\draw[->] (6) to [out=220,in=180,looseness=3] node[left ,pos=0.5] {$3$} (1);
 \end{tikzpicture} 
\end{equation*}
\caption{A presentation of $(\underline{\mathbb K}(D),L_{\underline{\mathbb K}(D)})$.}
\end{figure}
+++++++++++++++++++++++


+++++++++++++++++++++++++++++++++
Hence, in this example $(\underline{\mathbb K}(Y),L_{\underline{\mathbb K}(Y)})$ is different from $(\mathbb K'(Y),L_{\mathbb K'(Y)}) = (\mathbb F'(Y),L_{\mathbb F'(Y)})$ and $(\mathbb K(Y),L_{\mathbb K'(Y)}) = (\mathbb F'(Y),L_{\mathbb F'(Y)})$.


. In this case the labeled graph $G'$ is 

\begin{equation*}
\begin{tikzpicture}[node distance={35mm}, thick, main/.style = {draw, circle}] 
\node[main] (1) {$\{a\}$}; 
\node[main] (2) [right of=1] {$\{b\}$};
\node[main] (3) [below of=1] {$\{a,b\}$}; 
\draw[->] (1) to [out=170,in=200,looseness=20] node[left ,pos=0.5] {$1$} (1);
\draw[->] (1) to [out=250,in=290,looseness=10] node[below ,pos=0.5] {$2$} (1);
\draw[->] (1) to [out=300,in=200,looseness=1] node[above ,pos=0.5] {$0$} (2);
\draw[->] (2) to [out=120,in=70,looseness=1] node[above ,pos=0.5] {$1$} (1);
\draw[->] (2) to [out=250,in=290,looseness=10] node[below ,pos=0.5] {$2$} (2);
\draw[->] (3) to [out=250,in=290,looseness=10] node[below ,pos=0.5] {$2$} (3);
\draw[->] (3) to [out=110,in=220,looseness=1] node[left ,pos=0.5] {$1$} (1);
%\draw[->] (3) to [out=50,in=220,looseness=1] node[left ,pos=0.5] {$0$} (2);
 \end{tikzpicture} 
\end{equation*}
Thus $(G',L_{G'})$ is a proper subgraph of $(\underline{G},L_{\underline{G}})$ in this case.



++++++++++++++++++++++++++++++++++++++

++++++++++++++++++++++++++++++++++++++++





We end the paper by giving two examples where the future cover of a mixing sofic shift is the same as the minimal right-resolving cover, while the canonical cover from Theorem \ref{31-10-25hx} applied to the minimal right-resolving cover is different. The examples also show that in some cases the labeled graph $(\underline{G},L_{\underline{G}})$ from Section \ref{consec} is the same as $(G',L_{G'})$ and in others it isn't.
An important point concerning the construction is that it produces new strongly canonical covers of sofic shifts. Up to now only two strongly canonical covers for sofic shifts have been known; at least when we for simplicity ignore the doubling coming from the passage from right-resolving to left-resolving covers that can be achieved by 'time-reversal'. Besides the future cover it is the minimal right-resolving cover, and only for irreducible sofic shifts. It is therefore interesting to compare the strongly canonical covers that can be produced using Theorem \ref{31-10-25hx} with the future cover and, in the irreducible case, the minimal right-resolving cover. For this we introduce a method to find the future cover of a sofic shift  starting from an arbitrary given presentation. This is based on a variation of the so-called subset construction described in \S 3.3 of \cite{LM}. More specifically, it is shown that a restriction to a subset of the vertexes in the subset construction, applied to an arbitrary presentation of the sofic shift, gives rise to a labeled graph whose merged graph is isomorphic to the future cover. Thus the future cover is the factor obtained from a subgraph of the subset construction by identifying vertexes with the same follower set. This fact is used to produce the examples occurring in the paper, and makes it easy to show by examples that the new strongly canonical covers that can be obtained from Theorem \ref{31-10-25hx} in general are different both from the future cover, and in the irreducible case, also the minimal right-resolving cover.  
 
 
  In the form we use it here, the subset construction works well to find the future cover regardless of which presentation of the sofic shift we start with. In the last section of the paper we show that when the construction is applied to a weakly canonical cover it gives rise to a strongly canonical cover of the future cover and hence also of the sofic shift itself. See Theorem \ref{18-01-26d}. We show by examples that in general the strongly canonical covers coming from Theorem \ref{31-10-25hx} and Theorem \ref{18-01-26d} are different.
 
 


\begin{thm}\label{20-01-26} Assume that $\psi : Y \to Z$ is a conjugacy of sofic shifts and that $(G,L_G)$ and $(H,L_H)$ are right-resolving presentations of $Y$ and $Z$, respectively, such that
\begin{equation}\label{20-01-26b}
\xymatrix{
 X_G\ar[r]^-{\phi} \ar[d]_{L_G}& X_{H} \ar[d]^-{L_H}\\
   Y \ar[r]^{\psi} & Z}
\end{equation}
commutes. It follows that there is a unique conjugacy $\tilde{\phi} : X_{\underline{G}} \to X_{\underline{H}}$ such that
\begin{equation}\label{20-01-26c}
\xymatrix{
 X_{\underline{G}}\ar[r]^-{\tilde{\phi}} \ar[d]_{L_{\underline{G}}}& X_{\underline{H}} \ar[d]^-{L_{\underline{H}}}\\
   Y \ar[r]^{\psi} & Z}
\end{equation}
commutes. This conjugacy has the property that also
\begin{equation}\label{20-01-25d}
\xymatrix{\\
X_{\underline{G}} \ar[r]^-{\tilde{\phi}} \ar[d]_-{f_{\underline{G}}} & X_{\underline{H}}  \ar[d]^-{f_{\underline{H}}}\\
 X_{\mathbb K(Y)} \ar[d]_-{L_{\mathbb K(Y)}} \ar[r]^-{\psi_{\mathbb K}} &  X_{\mathbb K(Z)} \ar[d]^-{L_{\mathbb K(Z)}} \\
Y \ar[r]_-{\psi}  & Z }
\end{equation}
commutes.

\end{thm}
\begin{color}{red} 

1. step: For any $M \in \mathbb N$ there is an $N \in \mathbb N$ such that when $x \in X_{\underline{G}}$ and $i \in \mathbb Z$, there is an $i' \in [i-N,i]$ and a component $C \in\underline{G}$ such that $x_j \in C$ for all $j \in [i'-M,i']$.

\bigskip

2. step: Let $\tilde{x} \in X_C$ be a periodic element such that $\tilde{x}_{[i'-M,i']} =x_{[i'-M,i']}$. There is periodic point $p$ in $X_{G''}$ such that $r(x)_{[i'-M,i']} = p_{[i'-M,i']}$ and $s_{G''}(p_j) \subseteq s_{G''}(\tilde{x}_j)$ for all $j \in \mathbb Z$. 

\bigskip

3. step: Let $d$ be the width of the window of $\phi''$. Then
$$
\phi'(x)_j = \phi''(\tilde{x})_j \supseteq \phi''(p)_j = \phi''(r(x))_j
$$
for all $j \in [i'-M+d,i'-d]$.

\bigskip

4. step: There is a unique element $y \in X_{\underline{H}}[i'-d,\infty)$ such that $s_{\underline{H}}(y_{i'-d}) = s_{\underline{H}}(\phi'(x)_{i'-d})$ and $L_{\underline{H}}(y) = \psi(L_{\underline{G}}(x))_{[i'-d,\infty)}$.

In fact, $\phi'(x)_{i'-d} = \phi''(\tilde{x})_{i'-d}$ and $s_{G''}(\phi''(\tilde{x})_{i'-d}) \supseteq s_{G''}(\phi''(r(x))_{i'-d})$. Since $L_{H''}(\phi''(r(x)))_j = \psi(L_{G''}(r(x))_j = \psi(L_{\underline{G}}(x))_j$, the existence of $y$ follows from the definition of $\underline{H}$. Uniqueness follows because $(\underline{H},L_{\underline{H}})$ is right-resolving.
\bigskip

5. step: Define $\tilde{\phi} : X_{\underline{G}} \to X_{\underline{H}}$ such that
$$
\tilde{\phi}(x)_i = y_i.
$$
Note that $\tilde{\phi}$ is continuous and extends $\phi'|_{\overline{\Per X_{\underline{G}}}}$. 
\end{color}


+++++++++++++++++++++++++++++++++++++++++++++++


+++++++++++++++++++++++++++++++++++++++++++++++++

(division by Jan. 18)

++++++++++++++++++++++++++++++++++++++++++++++++++++


\section{A canonical cover of the non-wandering part of the future cover}\label{cancover}

Let $(G,L_G)$ be a right-resolving presentation of the sofic shift $Y$. The labeled graphs $(G',L_{G'})$ and $(\underline{G},L_{\underline G})$ are both subgraphs of $(\overline{G},L_{\overline{G}})$. While $(\underline{G},L_{\underline{G}})$ is obtained from $(\overline{G},L_{\overline{G}})$ by restricting to a subset of the vertexes in $\overline{G}$, $(G',L_{G'})$ comes from a restriction both to a subset of vertexes and arrows in $\overline{G}$. In this section we compare the two subgraphs and show that their intersection gives rise to a  canonical cover of the non-wandering part of the future cover of $Y$.








\begin{lemma}\label{08-10-25bxy} For all $y \in Y$ the element $\beta_G(y)$ is forward asymptotic to $\alpha_G(y)$.
\end{lemma}
\begin{proof} Note that
$$
s_{\underline{G}}(\alpha_G(y)_n) = \left\{t_G(x): \ x \in X_G(-\infty,n-1], \ L_G(x) = y_{(-\infty,n-1]}\right\} \supseteq  s_{G''}(\beta_G(y)_n) 
$$
for all $n \in \mathbb Z$. A compactness argument shows that there is an $N \in \mathbb Z$ such that
$$
\left\{t_G(x): \ x \in X_G(-\infty,n-1], \ L_G(x) = y_{(-\infty,n-1]}\right\} = s_{G''}(\beta_G(y)_n) 
$$
for all $n \geq N$. It follows that $\alpha_G(y)_j = \beta_G(y)_j$ for all $j \geq N$.
\end{proof}

When $y$ is periodic, both $\alpha_G(y)$ and $\beta_G(y)$ are periodic and hence Lemma \ref{08-10-25bxy} has the following

\begin{corollary}\label{09-10-25b} Let $p \in \Per (Y)$. Then $\beta_G(p) = \alpha_G(p)$.
\end{corollary}


\begin{lemma}\label{10-02-25dxxy} Let $F \in V_{\underline{G}}$. There is a periodic point $p \in \Per (Y)$ and an element $y' \in X_{\underline{G}}$ backward asymptotic to $\alpha_G(p)$ such that $t_{\underline{G}}(y'_{(-\infty,-1]}) = F$.
\end{lemma}
\begin{proof} Since $F \in V_{\underline{G}}$ there is an element $y \in Y$ such that $F = s_{\underline{G}}(\alpha_G(y)_0)$. Set $H_j := s_{\underline{G}}(\alpha(y)_j), \ j \leq 0$. Since $(G,L_G)$ is right-resolving there is an $N <0$ such that 
\begin{equation}\label{17-04-25bxy}
\# H_j = \# H_N
\end{equation} 
for all $j \leq N$.
 By the pigeon hole principle there is a sequence $\cdots < n_3 < n_2 <n_1 < n_0 \leq N$ and a vertex $H \in V_{\underline{G}}$ such that $H =s_{\underline{G}}(\alpha_G(y)_{n_j})$ for all $j \in \mathbb N \cup \{0\}$. Consider the loops in $\underline{G}$:
$$
\gamma_j := H \overset{y_{n_j}}{\to} H_{n_j+1}  \overset{y_{n_j+1}}{\to} \cdots \cdots \overset{y_{n_0-1}}{\to} H .
$$
If $\# L_G^{-1}(L_{\underline{G}}(\gamma_j)^\infty) > \# H$ for all $j$ it follows that there is an element $x' \in X_{\overline{G}}(-\infty,{n_0}-1]$ such that $H \subsetneq t_{\overline{G}}(x')$ and $L_{\overline{G}}(x') = y_{(-\infty,n_0-1]}$. By definition of $\alpha_G(y)$ this implies that $\# s_{\underline{G}}(\alpha_G(y)_j)  = \# H_j > \# H =\# H_N $ for all $j \leq n_0$, contradicting \eqref{17-04-25bxy}. This shows that there is a $j$ such that $q:= L_{\underline{G}}((\gamma_j)^\infty) \in Y$ is periodic and $\# L_G^{-1}(q) = \#H$. Thus $\alpha_G(q) = (\gamma_j)^\infty$ and 
$$
y' := \alpha_G(\sigma^{n_0}(q))_{(-\infty,n_0-1]}\alpha_G(y)_{[n_0,\infty)}
$$
has the stated properties.


\end{proof}


\begin{lemma}\label{11-11-24axy} Let $x,y \in X_{\underline{G}}$, and assume that $x = \alpha_G(L_{\underline{G}}(x))$. If there is an $N \in \mathbb Z$ such that $x_i = y_i$ for $i \leq N$, then $y = \alpha_G(L_{\underline{G}}(y))$.
\end{lemma}
\begin{proof} Since $L_{\underline{G}}(x)_i = L_{\underline{G}}(y)_i$ for $i \leq N$, it follows from the definition of 
$\alpha_G$ that 
$$
\alpha_G(L_{\underline{G}}(x))_j = \alpha_G(L_{\underline{G}}(y))_j
$$ 
for $j \leq N-1$. Hence 
$$
t_{\underline{G}}(y_{N-1}) =t_{\underline{G}}(x_{N-1}) = t_{\underline{G}}\left(\alpha_G(L_{\underline{G}}(x))_{N-1}\right) =  t_{\underline{G}}\left(\alpha_G(L_{\underline{G}}(y))_{N-1}\right).
$$
Since $L_{\underline{G}}(y_{[N,\infty)}) = L_{\underline{G}}\left(\alpha_G(L_{\underline{G}}(y))_{[N,\infty)}\right)$ it follows that $y_{[N,\infty)} = \alpha_G(L_{\underline{G}}(y))_{[N,\infty)}$ because $(\underline{G},L_{\underline{G}})$ is right resolving.  We know already that $y_{(-\infty,N-1]} = x_{(-\infty,N-1]} = \alpha_G(L_{\underline{G}}(x))_{(-\infty,N-1]} =  \alpha(L_{\underline{G}}(y))_{(-\infty,N-1]}$, and conclude therefore that $y = \alpha_G(L_{\underline{G}}(y))$.
\end{proof}


\begin{lemma}\label{03-07-25axy} Let $x \in X_{\underline{G}}$. There is an element $y \in Y$ such that $\beta_G(y)$ is forward asymptotic to $x$.
\end{lemma}
\begin{proof} By
Lemma \ref{10-02-25dxxy} there is an element $y' \in X_{\underline{G}}$ and a periodic point $p\in Y$ such that $y'$ is backward asymptotic to $\alpha_G(p)$ and forward asymptotic to $x$. Since $\alpha_G(p) =  \alpha_G(L_{\underline{G}}(\alpha_G(p)))$, it follows from Lemma \ref{11-11-24axy} that $y' = \alpha_G(L_{\underline{G}}(y'))$. Set $y := L_{\underline{G}}(y')$ and note that $\beta_G(y)$ is forward asymptotic to $\alpha_G(y) = y'$ by Lemma \ref{08-10-25bxy}. Hence $\beta_G(y)$ is also forward asymptotic to $x$.
\end{proof}


\begin{lemma}\label{08-10-25ay}
$$
\overline{\Per X_{\underline{G}}}
 = \bigcup_{y \in Y} \bigcap_{N \geq 1} \overline{\left\{\sigma^n(\beta_G(y)) : \ n \geq N \right\}} .
$$
\end{lemma}
\begin{proof} Let $y \in Y$. It follows from Lemma \ref{08-10-25bxy} that
$$
\bigcap_{N \geq 1} \overline{\left\{\sigma^n(\beta_G(y)) : \ n \geq N \right\}} = \bigcap_{N \geq 1} \overline{\left\{\sigma^n(\alpha_G(y)) : \ n \geq N \right\}}.
$$
Since there is a component of $X_{\underline{G}}$ which contains $\sigma^n(\alpha_G(y))$ for all sufficiently large $n$, it follows that 
$$
 \bigcap_{N \geq 1} \overline{\left\{\sigma^n(\beta_G(y)) : \ n \geq N \right\}} \subseteq \overline{\Per X_{\underline{G}}} .
 $$
This proves the inclusion $\supseteq$. The obtain the reversed inclusion, let $C$ be a component in $X_{\underline{G}}$. Choose an element $x \in X_C$ which is forward transitive in $X_C$; i.e. every word from $\mathbb W(X_C)$ occurs infinitely many times to the right in $x$. It follows from Lemma \ref{03-07-25axy} that there is an element $y \in Y$ such that 
\begin{equation}\label{08-10-25dx}
\lim_{n \to \infty} d\left(\sigma^n(\beta_G(y)),\sigma^n(x)\right) = 0.
\end{equation}
For each $z \in X_C$ there is a strictly increasing sequence $\{n_i\}$ in $\mathbb N$ such that
$\lim_{i \to \infty} d(\sigma^{n_i}(x),z) = 0$.
Thanks to \eqref{08-10-25dx}, by passing to a subsequence we can then arrange that
$$
\lim_{i \to \infty} d(\sigma^{n_i}(\beta_G(y)),z) = 0.
$$
It follows that $z \in  \bigcap_{N \geq 1} \overline{\left\{\sigma^n(\beta_G(y)) : \ n \geq N \right\}}$.
\end{proof}


\begin{corollary}\label{01-11-25}  $\overline{\Per X_{\underline{G}}} \subseteq X_{G'} \cap X_{\underline{G}}$
\end{corollary}




Let $\psi : Y \to Z$ be a conjugacy of sofic shifts. Let $(G,L_G)$ and $(H,L_H)$ be right-resolving presentations of $Y$ and $Z$, respectively. We assume that there is a conjugacy $\phi : X_G \to X_H$ such that
\begin{equation}\label{09-10-25ax}
\xymatrix{
 X_G\ar[r]^-{\phi} \ar[d]_{L_G}& X_{H} \ar[d]^-{L_H}\\
   Y \ar[r]^{\psi} & Z}
\end{equation}
commutes.



\begin{lemma}\label{09-10-25f} 
$$
\phi'( \overline{\Per X_{\underline{G}}}) = \overline{\Per X_{\underline{H}}} .
$$
and 
$$
f_{\underline{H}} \circ \phi' = \phi_{\mathbb K} \circ f_{\underline{G}}
$$
on $\overline{\Per X_{\underline{G}}}$.
\end{lemma}
\begin{proof} The first equality follows from Lemma \ref{08-10-25ay} and \eqref{30-10-25d}. Let $z \in \overline{\Per X_{\underline{G}}}$. By Lemma \ref{08-10-25ay} there is an $y \in Y$ such that $z = \lim_{i \to \infty} \sigma^{n_i}(\beta_{G}(y)) $ for some sequence $\{n_i\}$ in $\mathbb N$.
Note that 
$$
\lim_{i \to \infty} \sigma^{n_i}(\beta_{H}(y)) = \lim_{i \to \infty} \sigma^{n_i}(\alpha_{H}(y))
$$
by Lemma \ref{08-10-25bxy}. From \eqref{30-10-25d} we see that 
$\phi' \circ \beta_G = \beta_H \circ \psi$. 
Hence
\begin{align*}
&\phi'( \lim_{i \to \infty} \sigma^{n_i}(\beta_{G}(y))) = \lim_{i \to \infty} \sigma^{n_i}(\beta_{H}(\psi(y)) =  \lim_{i \to \infty} \sigma^{n_i}(\alpha_{H}(\psi(y)).
\end{align*}
Note that $f_{\underline{H}} \circ \alpha_H \circ \psi = \alpha_{Z} \circ \psi$ by Lemma \ref{11-09-25ex}. Since $\alpha_{Z}\circ\psi = \phi_{\mathbb K} \circ \alpha_Y$ by \eqref{24-11-24x} we conclude that
\begin{align*}
&f_{\underline{H}} \circ \phi'(z) = f_{\underline{H}} \circ \phi'( \lim_{i \to \infty} \sigma^{n_i}(\beta_{G}(y))) = \lim_{i \to \infty} f_{\underline{H}}(\sigma^{n_i}(\alpha_{H}(\psi(y))) \\
&  = \lim_{i \to \infty} \sigma^{n_i}\left( f_{\underline{H}} \circ \alpha_H \circ \psi(y)\right)= \lim_{i \to \infty}  \sigma^{n_i}(\alpha_{Z}(\psi(y)))= \lim_{i \to \infty}  \sigma^{n_i}(\phi_{\mathbb K}\circ\alpha_Y(y)) \\
& =  \phi_{\mathbb K}\left(\lim_{i \to \infty} \sigma^{n_i}(\alpha_{Y}(y))\right) =  \phi_{\mathbb K}\left(\lim_{i \to \infty} \sigma^{n_i}(f_{\underline{G}} \circ \alpha_{G}(y))\right) = \phi_{\mathbb K}\circ f_{\underline{G}}(z).
\end{align*}
%For comparison 
%\begin{align*}
%&\phi \circ f_{G}\left( \lim_{i \to \infty} \sigma^{n_i}(\beta_{G}(y))\right) \\
%& = \phi\left(  \lim_{i \to \infty} \sigma^{n_i}(f_{G} \circ \beta_{G}(y))\right)\\
%& = \phi \left( \lim_{i \to \infty} \sigma^{n_i}(\alpha_Y(y))\right)\\
%&  = \phi \left( \lim_{i \to \infty} \sigma^{n_i}(f_G\circ \alpha_G(y))\right)\\
%&= \phi \circ f_G(z) 
%\end{align*}
\end{proof}


This leads to the following:

\begin{thm}\label{09-10-25g} Let $\psi : Y \to Z$ be a conjugacy of sofic shifts. Assume that $(G,L_G)$ and $(H,L_H)$ are right-resolving presentations of $Y$ and $Z$, respectively. Assume that there is a conjugacy $\phi : X_G \to X_H$ such that \eqref{09-10-25ax} commutes. Then 
$$
\phi' : \ \overline{\Per X_{\underline{G}}} \to \overline{\Per X_{\underline{H}}}
$$
is a conjugacy and the diagram 
\begin{equation}\label{09-10-25x}
\xymatrix{\\
\overline{\Per (X_{\underline{G}})} \ar[r]^-{\phi'} \ar[d]_-{f_{\underline{G}}} & \overline{\Per (X_{\underline{H}})}  \ar[d]^-{f_{\underline{H}}}\\
 \overline{\Per X_{\mathbb K(Y)}} \ar[d]_-{L_{\mathbb K(Y)}} \ar[r]^-{\phi_{\mathbb K}} &  \overline{\Per X_{\mathbb K(Z)}} \ar[d]^-{L_{\mathbb K(Z)}} \\
\overline{\Per Y} \ar[r]_-{\psi}  & \overline{\Per Z }}
\end{equation}
commutes.


\end{thm}

In general the conjugacies $\phi_{\mathbb K}$ and $\phi'$ in the diagram \eqref{09-10-25x} are not unique.


Note that $\overline{\Per(X_G)} = X_{\hat{G}}$ where 
$$
\hat{G} := \bigsqcup_{C \in V_{\underline{G}}^r/\sim} G_C \subseteq \underline{G}.
$$
It it therefore natural to ask

\begin{question}\label{13-10-25} Does the conjugacy $\phi'$ of Theorem \ref{09-10-25g} extend to a conjugacy $\lambda : X_{\underline{G}} \to X_{\underline{H}}$ such that
\begin{equation*}\label{13-10-25a}
\xymatrix{\\
X_{\underline{G}} \ar[r]^-{\lambda} \ar[d]_-{f_{\underline{G}}} & X_{\underline{H}}  \ar[d]^-{f_{\underline{H}}}\\
 X_{\mathbb K(Y)} \ar[d]_-{L_{\mathbb K(Y)}} \ar[r]^-{\phi_{\mathbb K}} &  X_{\mathbb K(Z)} \ar[d]^-{L_{\mathbb K(Z)}} \\
Y \ar[r]_-{\psi}  & Z }
\end{equation*}
commutes ?
\end{question}

I think the answer is affirmative and that the following holds:

\emph{Conjecture} (17-01-26) There is a unique sliding block code 
$$
\tilde{\phi} : X_{\underline{G}} \to X_{\underline{H}}
$$
such that $L_{\underline{H}} \circ \tilde{\phi} = \psi \circ L_{\underline{G}}$ and $\tilde{\phi}(x) = \phi'(x)$ for $x \in \overline{\Per X_{\underline{G}}}$.


Let $C_1,C_2$ be components in $G''$. We write $C_1 \prec C_2$ when the following holds: For each vertex $F_2 \in C_2$ there is a vertex $F_1 \in C_1$ such that $F_1 \subseteq F_2$.

\begin{lemma}\label{17-01-26} Let $x \in X_{\underline{G}}$ and let $C$ be the component in $\underline{G}$ such that $x$ is backward asymptotic to $X_C$. There is an element $y \in X_{G''}$ such that $L_{\underline{G}}(x) = L_{G''}(y)$ and such that $y$ is backward asymptotic to a component $C'$ in $G''$ with the property that $C' \prec C$.
\end{lemma}

\begin{lemma}\label{18-01-26} Let $x \in X_{\underline{G}}$ and let $C$ be a component in $\underline{G}$. Let $N \in \mathbb N \cup \{\infty\}$. Assume that for some $i \in \mathbb Z$ we have that $x_j \in C$ for all $i-N < j \leq i$. For every $k > i$ there is a unique element $y \in X_{G''}[i-N+1,k]$ such that $L_{\underline{G}}(x_{[i-N+1,k]}) = L_{G''}(y)$, $s_{G''}(y_j) \subseteq s_{\underline{G}}( x_j)$ when $i-N + 1 \leq j \leq k$, and such that $t_{G''}(y_{k}) = t_{\underline{G}}(x_{k})$.
\end{lemma}

\begin{lemma}\label{18-01-26} There is a unique shift-commuting map $r: X_{\overline{G}} \to X_{G''}$ such that $x$ and $r(x)$ are forward asymptotic, $s_{G''}(r(x)_j) \subseteq s_{\overline{G}}(x_j)$ for all $j \in \mathbb Z$ and $L_{G''}(r(x)) = L_{\overline{G}}(x)$ for all $x \in X_{\overline{G}}$.
\end{lemma}
\begin{proof} Let $i \in \mathbb Z$. For each $N \in \mathbb N$ and $x \in X_{\overline{G}}$ there is a unique element $y^N \in X_{G''}(-\infty, i+N]$ such that $s_{G''}(y^N_j) \subseteq s_{\overline{G}}(x_j)$ for all $j \leq i+N$, $L_{G''}(y^N) = L_{\overline{G}}(x_{(-\infty,i+N]})$ and $t_{G''}(y^N_{i+N}) = t_{\overline{G}}(x_{i+N})$. Note that
$$
s_{G''}(y^N_i)
$$
is the set of vertexes $v \in s_{\overline{G}}(x_i)$ with the property that there is a path $\gamma$ in $G$ with $L_G(\gamma) = x_{[i,i+N]}$ and $s_G(\gamma) = v$. In particular $s_{G''}(y^{N+1}_i) \subseteq s_{G''}(y^N_i)$, implying that
$$
\lim_{N \to \infty} s_{G''}(y^N_i)
$$
exists. It follows that there is a unique element $r(x) \in X_{G''}$ such that $L_{G''}(r(x)) = L_{\overline{G}}(x)$ and $s_{G''}(r(x)_i) = \lim_{N \to \infty} s_{G''}(y^N_i)$.
\end{proof}

Note that $r(x) = x$ when $x \in X_{G''}$.

Let $M \in \mathbb N$.  Let $d \in \mathbb N$ be bigger than the width of the block code defining $\phi''$ and $\psi$. Consider $x \in X_{\underline{G}}$ and $i \in \mathbb Z$. Let $i' \leq i$ be the largest integer such that $x_{[i'-M,i']}$ is contained in a component $C$ of $\underline{G}$. There is then a component $C'$ in $\underline{H}$ such that $\phi'(x)_j \in C'$ for $j \in [i'-M+d, i'-d]$. It follows from Lemma \ref{17-01-26} that there is a component $C_1$ in $G''$ such that $C_1 \prec C$ and an element $y \in X_{\underline{G}}$ such that $L_{\underline{G}}(y)_{[i'-M,\infty)} = L_{\underline{G}}(x)_{[i'-M,\infty)}$ and $y_{[i'-M,i']} \in C_1$. Since $\phi'$ is the restriction of $\phi''$, it follows that because $C_1 \prec C$ we have also that $C'_1 := \phi''(C_1) \prec C'$.  (Check this!) 

Note that $L_{\underline{H}}(\phi''(y))_{[i'-M+d,i'-d]} = L_{\underline{H}}(\phi'(x))_{[i'-M+d,i'-d]} = \psi\left( L_{\underline{G}}(x)\right)_{[i'-M+d,i'-d]}$. 


It follows therefore that there is a unique $z \in X_{\underline{H}}([i'-M+d,\infty)$ such that $z_ {[i'-M+d,i'-d]} =\phi'(x)_{[i'-M+d,i'-d]}$ and $L_{\underline{H}}(z) = \psi\left( L_{\underline{G}}(x)\right)_{[i'-M+d,\infty}$. Set
$$
\tilde{\phi}(x)_i := z_i .
$$


 \begin{thm}\label{18-01-26b} The extension $ L_{\underline{\mathbb K(Y)}} : \underline{\mathbb K(Y)}\to Y$ is strongly canonical: When $\psi : Y \to Z$ is a conjugacy of sofic shifts there is a unique conjugacy 
 $$
 \psi_{\underline{\mathbb K}} : \ X_{\underline{\mathbb K(Y)}} \to  X_{\underline{\mathbb K(Z)}}
 $$
 such that
 \begin{equation*}\label{13-10-25a}
\xymatrix{\\
X_{\underline{\mathbb K(Y)}} \ar[r]^-{\psi_{\underline{\mathbb K}}} \ar[d]_-{f_{\underline{\mathbb K(Y)}}} & X_{\underline{\mathbb K(Z)}}  \ar[d]^-{f_{\underline{\mathbb K(Z)}}}\\
 X_{\mathbb K(Y)} \ar[d]_-{L_{\mathbb K(Y)}} \ar[r]^-{\psi_{\mathbb K}} &  X_{\mathbb K(Z)} \ar[d]^-{L_{\mathbb K(Z)}} \\
Y \ar[r]_-{\psi}  & Z }
\end{equation*}
commutes.
 \end{thm}




\begin{color}{red} I think there is a sliding block code $r : X_{\overline{G}} \to X_{G''}$ defined from the observation that for any $M \in \mathbb N$ there is an $N \in \mathbb N$, such that for $x \in X_{\overline{G}}$ and $i \in \mathbb N$ there are components $C_-$ and $C_+$ in $\overline{G}$ such that $x_{[i- N,i-N+M]} \in C_-$ and $ x_{[i+N-M,i+N]} \in C_+$. Note that $C_\pm$ are also components in $G''$ and that there is a unique path $y_{[i-N,i+N]}$ in $G''$ such that $y_{[i+N-M,i+N]} = x_{[i+N-M,i+N]}$. Define $r(x)_i$ to be $y_i$. Note that $r(x) = x$ when $x \in G''$ and hence that $G''$ is a retract of $\overline{G}$.
\end{color}
\begin{color}{red} 
\begin{lemma}\label{foerstejuledag} Let $C$ be a component in $H''$ and $z \in X_{H''}$ an element backward transitive to $X_C$. Then $z= \beta_H(L_{H''}(z))$ if $C$ has the following property:
\begin{itemize}
\item[(*)] If $C'$ is a component in $H''$ and $L_{H''}(X_C) \subseteq L_{H''}(X_{C'})$, then $M(C') \subseteq M(C)$.
\end{itemize}
\end{lemma}
\begin{proof} Note that $\beta_H(L_{H''}(z))$ is backward asymptotic to some component $X_{C'}$. Since $L_{H''}(z)$ is backward transitive to $L_{H''}(X_C)$, it follows that $L_{H''}(X_C) \subseteq L_{H''}(X_{C'})$. Since $C$ has property (*), it follows that $M(C') \leq M(C)$. Since $s_{H''}(z_j) \subseteq s_{H''}(\beta_H(L_{H''}(z))_j)$ for all $j$, it follows first that $M(C) = M(C')$ and then that $z= \beta_H(L_{H''}(z))$.

%'only if': Assume $z= \beta_H(L_{H''}(z))$ and that $L_{H''}(X_C) \subseteq L_{H''}(X_{C'})$. Then $L_{H''}(z) = L_{H''}(z')$ for some $z' \in X_{C'}$ (NO NO NO). This implies that $s_{H''}(z'_j) \subseteq s_{H''}(\beta_H(L_{H''}(z))_j) = s_{H''}(z_j)$ for all $j$ 
\end{proof} 

\bigskip


1. Argument: Let $y \in Y$ and choose a component $C$ in $H''$ and $N \in \mathbb Z$ such that $\beta_H(y)_j \in C$ for all $j \leq N$. It follows that there is an $ M < N$ such that if $z \in X_{H''}$ and $L_{H''}(z)_{[M,\infty)} = y_{[M,\infty)}$, then $s_{H''}(z_N) \subseteq s_{H''}(\beta_H(y)_N)$. (If not, a compactness argument would produce an element $z' \in X_{H''}$ such that $L_{H''}(z') = y$ and $\# s_{H''}(z'_N) > \# s_{H''}(\beta_H(y)_N)$. A contradiction.) Let $C'$ be component which connects to $C$ and has the largest multiplicity among the components that connects to $C$. Let $\gamma$ be a finite path in $H_C$. There is an element $z'  \in X_{H''}$ with the following properties:
\begin{itemize}
\item[i)] $z'$ is backward asymptotic to $C'$,
\item[ii)] there is a $M' < M$ such that $z'_{[M',M]}$ is a path in $C$ with $\gamma \subseteq z'_{[M',M]}$, and
\item[iii)] $z'_{[M,\infty)} = \beta_H(y)_{[M,\infty)}$.
\end{itemize}
We claim that 
\begin{equation}\label{foerstejuledag1}
z' = \beta_H(L_{H''}(z')).
\end{equation}
To see this let $C''$ be the component with the property that $\beta_H(L_{H''}(z'))$ is backward asymptotic to $X_{C''}$. Note that 
\begin{equation}\label{foerstejuledag2} 
\beta_H(L_{H''}(z'))_N =  \beta_H(y)_N
 \end{equation}
 since $L_{H''}(\beta_H(L_{H''}(z'))_{[M,\infty)} = y_{[M,\infty)}$ by iii), which implies that 
 $$
 s_{H''}(\beta_H(L_{H''}(z'))_N) \subseteq  s_{H''}(\beta_H(y)_N)
 $$ 
 by choice of $M$ while the opposite inclusion follows from 
 $$
 s_{H''}(\beta_H(y)_N) = s_{H''}(z'_N) \subseteq s_{H''}(\beta_H(L_{H''}(z'))_N).
 $$ 
 Since $s_{H''}(z'_j) \subseteq s_{H''}(\beta_H(L_{H''}(z'))_j)$ for all $j$ we have that $M(C')  \leq M(C'')$. But it follows from \eqref{foerstejuledag2} that $C''$ connects to $C$ and hence $M(C'') \leq M(C')$ by definition of $C'$ and we find therefore that $M(C'') = M(C')$. Since $s_{H''}(z'_j) \subseteq s_{H''}(\beta_H(L_{H''}(z'))_j)$ for all $j$, it follows from this identity that $s_{H''}(z'_j) = s_{H''}(\beta_H(L_{H''}(z'))_j)$ for all $j$ 'close to $-\infty$'; a conclusion which implies the claim; \eqref{foerstejuledag1}.



\begin{cor}\label{28-12-25} $X_C \subseteq \overline{\beta_H(Y)}$.
\end{cor}
\begin{proof} This follows because the $\gamma$ is arbitrary and $\beta_H(Y)$ is $\sigma$-invariant.
\end{proof}

In the setting of the proof of Theorem \ref{31-10-25h}: The first reference to Lemma 4.8 must be replaced by the observation that because $\phi''(\beta_G(y)) = \beta_H(\phi_H(y))$ implies that when $C_1$ is a component in $G''$ such that $\beta_G(y)$ is backward asymptotic to $X_{C_1}$, then $\phi''(X_{C_1}) \subseteq X_{C_2}$ where $C_2$ is a component in $H''$ such that $\beta_H(\psi(y))$ is backward asymptotic to $X_{C_2}$. By symmetry $\phi''(X_{C_1}) = X_{C_2}$. Two relatively easy arguments then imply that $\phi''(X_{G''}) = X_{H''}$. 

The second reference to Lemma 4.8 can be replaced by the observation that $\beta_H(Y) \subseteq X_{H'}$. 

The third should be replaced by the following: It follows from Corollary \ref{28-12-25} that $\phi_1(x) = \phi_2(x)$ when $x \in X_C$ when $C$ is a component in $G''$ such that $\beta_G(y)$ is backward asymptotic to $X_{C}$. Assume then that $x$ is backward asymptotic to $X_C$ for such a $C$. Then $\phi_1(x)_{(-\infty, N]} = \phi_2(x)_{(-\infty,N]}$ for some $N \in \mathbb Z$ since $\phi_1$ and $\phi_2$ agree on $X_C$. Since $L_{H''}(\phi_1(x)) = \psi(L_{G''}(x)) = L_{H''}(\phi_2(x))$ we find that $\phi_1(x) = \phi_2(x)$ because $L_{H''}$ is right-resolving. 

\bigskip

2. Argument: Let $z \in X_{H''}$ and assume that $z$ is backward asymptotic to $\beta_H(p)$ where $p \in \Per(Y)$. Choose $N \in \mathbb Z$ such that $z_j = \beta_H(p)_j, \ j \leq N$. Since $L_{H''}(\beta_H(L_{H''}(z)))_j = L_{H''}(z)_j = p_j$ for $j \leq N$, the periodicity implies that $\beta_H(L_{H''}(z))_j = q_j, \ j \leq N$, where $q \in X_{H''}$ is periodic, necessarily with $L_{H''}(q) =p$. (This follows because $L_{H''}$ is right-resolving.) It follows that 
$$
s_{H''}(\beta_H(L_{H''}(z))_j) = s_{H''}(q_j) \subseteq s_{H''}(\beta_H(p))
$$
for $j \leq N$. On this other hand
$$
s_{H''}(q_j) = s_{H''}(\beta_H(L_{H''}(z))_j) \supseteq s_{H''}(z_j) = s_{H''}(\beta_H(p)_j),
$$
also for $j \leq N$. Combined this gives first that $s_{H''}(\beta_H(L_H(z))_j) = s_{H''}(z_j)
$ for $j \leq N$ and then that
$$
\beta_H(L_H(z)) = z.
$$


\end{color}

\begin{lemma}\label{31-10-25c} Let $C$ be a component in $H''$ such that there is an element $y \in Y$ with $\beta_H(y)$ backward asymptotic to $X_C$, and let $z \in X_{H''} $ be an element backward transitive to $X_C$. Then $z = \beta_{H} (L_{H''} (z))$.
\end{lemma}

\begin{proof} Set $y' := L_{H''}(z)$. Then $y'$ is backward transitive to $L_{H''}(X_C)$. Let $C'$ be the component in $H''$ with the property that $\beta_H(y')$ is backward asymptotic to $X_{C'}$. Note that the number
$$
M(C):= \# s_{H''} (x_j)
$$
is the same for all elements $x \in X_C$ and all $j \in \mathbb Z$. Since $L_{H''} (z) = y'$ we have that 
\begin{equation}\label{31-10-25d}
s_{H''} (z_j) \subseteq s_{H''} (\beta_H(y')_j)
\end{equation} 
for all $j \in \mathbb Z$, implying that $M(C) \leq M(C')$ since $z$ is backward asymptotic to $X_C$ and $\beta_H(y')$ to $X_{C'}$. On the other hand, if $w$ is a word in $L_{H''} (X_C)$ it occurs infinitely often to the left in $y'$. Since $L_{H''} ^{-1}(w)$ is a finite set and $L_{H''} (\beta_H(y')) = y'$ it follows that $w = L_{H''}(w')$ for some word $w'$ appearing infinitely often to the left in $\beta_H(y')$. Then $w' \in \mathbb W(X_{C'})$ and we conclude therefore that $L_{H''} (X_{C}) \subseteq L_{H''} (X_{C'})$. It follows that there is a $z' \in X_{C'}$ such that $y$ is backward asymptotic to $L_{H''}(z')$, and there is therefore a $N \in \mathbb Z$ such that $s_{H''} (z'_j) \subseteq s_{H''} (\beta_H(y)_j)$ for all $j \leq N$. It follows that $M(C') \leq M(C)$ and therefore also that $M(C) = M(C')$. Then \eqref{31-10-25d} implies that $z = \beta_H(y') = \beta_H(L_{H''} (z))$.
\end{proof}


\begin{definition}\label{30-10-25} We denote by $H'$ the hereditary subgraph of $H''$ generated by the components $C$ of $H''$ with the property that there is a point $y \in Y$ such that $\beta_H(y)$ is backward asymptotic to $X_C$. 
\end{definition}

To simplify the notation, set $L_{H'} := L_{H''}|_{H'}$.

\begin{lemma}\label{31-10-25f} $X_{H'} = \overline{\beta_H(Y)}$. In particular, $(X_{H'},L_{H'})$ is a right-resolving presentation of $Y$.
\end{lemma}
\begin{proof} By definition of $H'$ every element of $X_{H'}$ can be approximated by elements $z$ with the property that $z$ is backward transitive to $X_C$, where $C$ is the component in $H''$ with the property that $\beta_H(y)$ is backward asymptotic to $X_C$ for some $y \in Y$. Since $z \in \beta_H(Y)$ by Lemma \ref{31-10-25c}, it follows that $X_{H'} \subseteq \overline{\beta_H(Y)}$. On the other hand, every element of $\beta_H(Y)$ is backward asymptotic to some $X_C$, where $C$, by definition, is a component in $H'$. Hence $\beta_H(Y) \subseteq X_{H'}$.
\end{proof}


\begin{thm}\label{31-10-25h} Let $(G,L_G)$ and $(H,L_H)$ be right-resolving presentations of the sofic shifts $Y$ and $Z$, respectively. Assume that there are conjugacies $\psi : Y \to Z$ and $\phi : X_G \to X_H$ such that 
\begin{equation*}\label{30-10-25bx}
\xymatrix{
 X_G \ar[r]^-{\phi} \ar[d]_{L_G}&  X_H \ar[d]^{L_H} \\
 Y  \ar[r]^-\psi & Z}
\end{equation*}
commutes. There is a unique conjugacy $\phi' : X_{G'} \to X_{H'}$ such that
\begin{equation*}\label{30-10-25cx}
\xymatrix{
 X_{G'} \ar[r]^-{\phi'} \ar[d]_{L_{G'}}&  X_{H'} \ar[d]^{L_{H'}} \\
Y   \ar[r]^-\psi & Z}
\end{equation*}
commutes.
\end{thm}
\begin{proof} It follows from Corollary \ref{17-11-25} that there is a conjugacy $\phi'' : X_{G''} \to X_{H''}$ such that
\begin{equation*}\label{22-11-25a}
\xymatrix{
 X_{G''} \ar[r]^-{\phi''} \ar[d]_{L_{G''}}&  X_{H''} \ar[d]^{L_{H''}} \\
 Y  \ar[r]^-\psi & Z}
\end{equation*}
commutes. When we use the same symbol $\kappa$ for the two conjugacies $\kappa : X_G^{L_G} \to X_{G''}$ and $\kappa : X_H^{L_H} \to X_{H''}$ from Lemma \ref{23-06-25dyx}, we note that $\phi'' = \kappa \circ \phi^\# \circ \kappa^{-1}$. Since $\phi^\#(L_G^{-1}(y)) = L_H^{-1}(\psi(y))$ it follows that $\phi''(\beta_G(y)) = \beta_H(\psi(y))$ for all $y \in Y$. By Lemma \ref{31-10-25f} this implies that $\phi''(X_{G'}) = X_{H'}$, and we set $\phi':= \phi''|_{X_{G'}}$. This proves the existence part. To prove uniqueness, consider a conjugacy $\lambda : X_{G'} \to X_{H'}$ such that
\begin{equation*}\label{30-10-25cy}
\xymatrix{
 X_{G'} \ar[r]^-{\lambda} \ar[d]_{L_{G'}}&  X_{H'} \ar[d]^{L_{H'}} \\
Y   \ar[r]^-\psi & Z}
\end{equation*}
commutes. Set 
$$
\lambda' := \kappa^{-1}\circ \lambda \circ \kappa|_{\kappa^{-1}(X_{G'})}.
$$ 
By Lemma \ref{23-06-25dyx} also
\begin{equation*}\label{30-10-25cxx}
\xymatrix{
 X_{G}^{L_G}\cap \kappa^{-1}(X_{G'}) \ar[r]^-{\lambda'} \ar[d]_{L_{G}^{\#}}&  X_{H}^{L_H}\cap \kappa^{-1}(X_{H'}) \ar[d]^{L_{H}^{\#}} \\
Y   \ar[r]^-\psi & Z}
\end{equation*}
commutes. Note that $L_G^{-1}(y) \in X_G^{L_G} \cap \kappa^{-1}(X_{G'})$ by Lemma \ref{31-10-25f}. It follows that $\lambda'(L_G^{-1}(y)) \subseteq L_H^{-1}(\psi(y))$ and ${\lambda'}^{-1}(L_H^{-1}(\psi(y))) \subseteq L_G^{-1}(y)$, implying that $\lambda'(L_G^{-1}(y)) =L_H^{-1}(\psi(y))$ for all $y \in Y$ and hence $\lambda(\beta_G(y)) = \beta_H(\psi(y)) = \phi'(\beta_G(y))$ for all $y \in Y$. By Lemma \ref{31-10-25f} this implies that $\lambda = \phi'$.


\end{proof}

\begin{ex} \label{08-01-25a}


The following labeled graph $(G,L_G)$ is irreducible, right resolving and follower-separated and hence the minimal right resolving graph of the sofic shift $Y$ it represents. Note that $Y$ is mixing.


\begin{equation*}
 \begin{tikzpicture}[node distance={35mm}, thick, main/.style = {draw, circle}] 
\node[main] (1) {$a$}; 
\node[main] (2) [right of=1] {$b$}; 
\node[main] (3) [right of=2] {$c$};
\draw[->] (1) to [out=170,in=200,looseness=20] node[left ,pos=0.5] {$2$} (1); 
\draw[->] (1) to [out=80,in=100,looseness=2] node[above ,pos=0.5] {$1$} (2); 
\draw[->] (1) to [out=45,in=150,looseness=1.5] node[above ,pos=0.5] {$0$} (2); 
\draw[->] (2) to [out=260,in=280,looseness=2] node[above ,pos=0.5] {$1$} (1);
\draw[->] (2) to [out=230,in=300,looseness=1.5] node[above ,pos=0.5] {$0$} (1); 
\draw[->] (2) to [out=45,in=170,looseness=1] node[above ,pos=0.5] {$2$} (3); 
\draw[->] (3) to [out=200,in=350,looseness=1] node[below ,pos=0.5] {$3$} (2);
\draw[->] (3) to [out=45,in=350,looseness=10] node[above,pos=0.5] {$0$} (3);  
\draw[->] (3) to [out=70,in=290,looseness=30] node[right,pos=0.5] {$1$} (3); 
\end{tikzpicture} 
\end{equation*}
This labeled graph is irreducible, right resolving and follower-separated and hence the minimal right resolving graph of the sofic shift $Y$ it represents. Note that $Y$ is mixing.

After removal of inessential vertexes the graph $(\overline{G},L_{\overline{G}})$ becomes the following.

\begin{equation}
 \begin{tikzpicture}[node distance={35mm}, thick, main/.style = {draw, circle}] 
\node[main] (1) {$\{a\}$}; 
\node[main] (2) [right of=1] {$\{b\}$}; 
\node[main] (3) [right of=2] {$\{c\}$};
\node[main] (4) [below of=2] {$\{a,c\}$};
\node[main] (5) [right of=4] {$\{b,c\}$};
\node[main] (6) [below of=5] {$\{a,b,c\}$};
\node[main] (7) [below of=4] {$\{a,b\}$};
\draw[->] (1) to [out=170,in=200,looseness=20] node[left ,pos=0.5] {$2$} (1); 
\draw[->] (1) to [out=80,in=100,looseness=2] node[above ,pos=0.5] {$1$} (2); 
\draw[->] (1) to [out=45,in=150,looseness=1.5] node[above ,pos=0.5] {$0$} (2); 
\draw[->] (2) to [out=260,in=280,looseness=2] node[above ,pos=0.5] {$1$} (1);
\draw[->] (2) to [out=230,in=300,looseness=1.5] node[above ,pos=0.5] {$0$} (1); 
\draw[->] (2) to [out=45,in=170,looseness=1] node[above ,pos=0.5] {$2$} (3); 
\draw[->] (3) to [out=200,in=350,looseness=1] node[above ,pos=0.5] {$3$} (2);
\draw[->] (3) to [out=45,in=350,looseness=10] node[above,pos=0.5] {$0$} (3);  
\draw[->] (3) to [out=60,in=290,looseness=15] node[right,pos=0.5] {$1$} (3); 
\draw[->] (6) to [out=260,in=290,looseness=10] node[below,pos=0.5] {$0$} (6); 
\draw[->] (6) to [out=10,in=350,looseness=10] node[right,pos=0.5] {$1$} (6); 
\draw[->] (6) to [out=100,in=300,looseness=0.5] node[left,pos=0.5] {$2$} (4); 
\draw[->] (6) to [out=50,in=60,looseness=2.5] node[right,pos=0.5] {$3$} (2); 
\draw[->] (4) to [out=40,in=150,looseness=1.5] node[above,pos=0.5] {$0$} (5); 
\draw[->] (4) to [out=20,in=170,looseness=1] node[above,pos=0.5] {$1$} (5); 
\draw[->] (5) to [out=190,in=350,looseness=0.5] node[below,pos=0.5] {$0$} (4); 
\draw[->] (5) to [out=220,in=320,looseness=1] node[below,pos=0.3] {$1$} (4); 
\draw[->] (4) to [out=170,in=210,looseness=1] node[below,pos=0.3] {$2$} (1); 
\draw[->] (7) to [out=90,in=270,looseness=1] node[left,pos=0.5] {$2$} (4); 
\draw[->] (7) to [out=260,in=290,looseness=10] node[below,pos=0.5] {$0$} (7);
\draw[->] (7) to [out=170,in=200,looseness=10] node[left,pos=0.5] {$1$} (7);
\draw[->] (4) to [out=90,in=270,looseness=1] node[left,pos=0.4] {$3$} (2);
\draw[->] (5) to [out=90,in=270,looseness=1] node[left,pos=0.5] {$2$} (3);
\draw[->] (5) to [out=100,in=300,looseness=1] node[left,pos=0.5] {$3$} (2);
\end{tikzpicture} 
\end{equation}
The graph $(\underline{G},L_{\underline{G}})$ is the following labeled subgraph of  $(\overline{G},L_{\overline{G}})$. We note that it is follower-separated and hence also (isomorphic to) the future cover of $Y$.

\begin{equation*}
 \begin{tikzpicture}[node distance={35mm}, thick, main/.style = {draw, circle}] 
\node[main] (1) {$\{a\}$}; 
\node[main] (2) [right of=1] {$\{b\}$}; 
\node[main] (3) [right of=2] {$\{c\}$};
\node[main] (4) [below of=2] {$\{a,c\}$};
\node[main] (5) [right of=4] {$\{b,c\}$};
\node[main] (6) [below of=5] {$\{a,b,c\}$};
%\node[main] (7) [below of=4] {$\{a,b\}$};
\draw[->] (1) to [out=170,in=200,looseness=20] node[left ,pos=0.5] {$2$} (1); 
\draw[->] (1) to [out=80,in=100,looseness=2] node[above ,pos=0.5] {$1$} (2); 
\draw[->] (1) to [out=45,in=150,looseness=1.5] node[above ,pos=0.5] {$0$} (2); 
\draw[->] (2) to [out=260,in=280,looseness=2] node[above ,pos=0.5] {$1$} (1);
\draw[->] (2) to [out=230,in=300,looseness=1.5] node[above ,pos=0.5] {$0$} (1); 
\draw[->] (2) to [out=45,in=170,looseness=1] node[above ,pos=0.5] {$2$} (3); 
\draw[->] (3) to [out=200,in=350,looseness=1] node[above ,pos=0.5] {$3$} (2);
\draw[->] (3) to [out=45,in=350,looseness=10] node[above,pos=0.5] {$0$} (3);  
\draw[->] (3) to [out=60,in=290,looseness=18] node[right,pos=0.5] {$1$} (3); 
\draw[->] (6) to [out=260,in=290,looseness=10] node[below,pos=0.5] {$0$} (6); 
\draw[->] (6) to [out=10,in=350,looseness=10] node[right,pos=0.5] {$1$} (6); 
\draw[->] (6) to [out=100,in=300,looseness=0.5] node[left,pos=0.5] {$2$} (4); 
\draw[->] (6) to [out=50,in=60,looseness=3] node[right,pos=0.5] {$3$} (2); 
\draw[->] (4) to [out=40,in=150,looseness=1.5] node[above,pos=0.5] {$0$} (5); 
\draw[->] (4) to [out=20,in=170,looseness=1] node[above,pos=0.5] {$1$} (5); 
\draw[->] (5) to [out=190,in=350,looseness=0.5] node[below,pos=0.5] {$0$} (4); 
\draw[->] (5) to [out=220,in=320,looseness=1] node[below,pos=0.3] {$1$} (4); 
\draw[->] (4) to [out=170,in=210,looseness=1] node[below,pos=0.3] {$2$} (1); 
%\draw[->] (7) to [out=90,in=270,looseness=1] node[left,pos=0.5] {$2$} (4); 
%\draw[->] (7) to [out=260,in=290,looseness=10] node[below,pos=0.5] {$0$} (7);
%\draw[->] (7) to [out=170,in=200,looseness=10] node[left,pos=0.5] {$1$} (7);
\draw[->] (4) to [out=90,in=270,looseness=1] node[left,pos=0.4] {$3$} (2);
\draw[->] (5) to [out=90,in=270,looseness=1] node[left,pos=0.5] {$2$} (3);
\draw[->] (5) to [out=100,in=300,looseness=1] node[left,pos=0.5] {$3$} (2);
\end{tikzpicture} 
\end{equation*}

It may be interesting to observe that in this example there are components in future cover where the restriction of the labeling from the future cover is not follower-separated.


In this example $(G',L_{G'})$ is the following labeled graph.

\begin{equation*}
 \begin{tikzpicture}[node distance={35mm}, thick, main/.style = {draw, circle}] 
\node[main] (1) {$\{a\}$}; 
\node[main] (2) [right of=1] {$\{b\}$}; 
\node[main] (3) [right of=2] {$\{c\}$};
\node[main] (4) [below of=2] {$\{a,c\}$};
\node[main] (5) [right of=4] {$\{b,c\}$};
\node[main] (6) [below of=5] {$\{a,b,c\}$};
\node[main] (7) [below of=4] {$\{a,b\}$};
\draw[->] (1) to [out=170,in=200,looseness=20] node[left ,pos=0.5] {$2$} (1); 
\draw[->] (1) to [out=80,in=100,looseness=2] node[above ,pos=0.5] {$1$} (2); 
\draw[->] (1) to [out=45,in=150,looseness=1.5] node[above ,pos=0.5] {$0$} (2); 
\draw[->] (2) to [out=260,in=280,looseness=2] node[above ,pos=0.5] {$1$} (1);
\draw[->] (2) to [out=230,in=300,looseness=1.5] node[above ,pos=0.5] {$0$} (1); 
\draw[->] (2) to [out=45,in=170,looseness=1] node[above ,pos=0.5] {$2$} (3); 
\draw[->] (3) to [out=200,in=350,looseness=1] node[above ,pos=0.5] {$3$} (2);
\draw[->] (3) to [out=45,in=350,looseness=10] node[above,pos=0.5] {$0$} (3);  
\draw[->] (3) to [out=70,in=290,looseness=30] node[right,pos=0.5] {$1$} (3); 
\draw[->] (6) to [out=260,in=290,looseness=10] node[below,pos=0.5] {$0$} (6); 
\draw[->] (6) to [out=10,in=350,looseness=10] node[right,pos=0.5] {$1$} (6); 
%\draw[->] (6) to [out=100,in=300,looseness=0.5] node[left,pos=0.5] {$2$} (4); 
%\draw[->] (6) to [out=70,in=60,looseness=3.5] node[left,pos=0.5] {$3$} (2); 
\draw[->] (4) to [out=40,in=150,looseness=1.5] node[above,pos=0.5] {$0$} (5); 
\draw[->] (4) to [out=20,in=170,looseness=1] node[above,pos=0.5] {$1$} (5); 
\draw[->] (5) to [out=190,in=350,looseness=0.5] node[below,pos=0.5] {$0$} (4); 
\draw[->] (5) to [out=220,in=320,looseness=1] node[below,pos=0.3] {$1$} (4); 
%\draw[->] (4) to [out=170,in=210,looseness=1] node[below,pos=0.3] {$2$} (1); 
\draw[->] (7) to [out=90,in=270,looseness=1] node[left,pos=0.5] {$2$} (4); 
\draw[->] (7) to [out=260,in=290,looseness=10] node[below,pos=0.5] {$0$} (7);
\draw[->] (7) to [out=170,in=200,looseness=10] node[left,pos=0.5] {$1$} (7);
%\draw[->] (4) to [out=90,in=270,looseness=1] node[left,pos=0.4] {$3$} (2);
%\draw[->] (5) to [out=90,in=270,looseness=1] node[left,pos=0.5] {$2$} (3);
%\draw[->] (5) to [out=100,in=300,looseness=1] node[left,pos=0.5] {$3$} (2);
\end{tikzpicture} 
\end{equation*}

Hence none of the two graphs $\underline{G}$ or $G'$ is contained in the other, but in this case $(\underline{G},L_{\underline G})$ is isomorphic to the future cover.
\end{ex}


%In conclusion, I suppose it is clear that the new canonical covers from Theorem \ref{31-10-25h} and Theorem \ref{09-10-25g} give rise to conjugacy invariants for sofic shifts that seem not to have been investigated before. The same is to some extend true for the invariants coming from the part of the future cover of an irreducible sofic shift that is not contained in the terminal component. It seems worthwhile to investigate all of this in more detail and to compare with the conjugacy invariants from \cite{T}. With a bit of luck this will be the focus in a forthcoming publication.  
++++++

Up to now there 



Perhaps the most important new information coming out of the paper is that the future cover, and for irreducible sofic shifts also the minimal right-resolving cover, are not the only canonical covers for a sofic shift. %In \cite{Kr1} Krieger observes about the work on covers for sofic shifts that pre-dates his work, that 'it seems not to have noticed that some of these constructions are canonical'. It is tempting to add now that it seems not to have been noticed that they are not the only canonical covers. 
 
 
The first section below, the second section of the paper, contains no new results; only new proofs of the canonicity of the future cover and, for irreducible sofic shifts, the canonicity of the cover given by the minimal right-resolving presentation, \cite{LM}. These results were obtained by Krieger in the two papers \cite{Kr1}, \cite{Kr2}. While Krieger's proof is a tour de force in handling sliding block codes, the proof given here is more conceptual and leads to an abstract characterization of the future cover analogous to the abstract characterization of the second canonical cover of an irreducible sofic shift as the minimal right-resolving presentation. Alternative proofs of parts of Krieger's theorem have been given before; first by Boyle, Kitchens and Marcus in \cite{BKM} and later by Nasu in \cite{N}. The proofs presented here are also different from those.

The third section contains a study of the universal properties of the future cover and its relation to the follower set graph. This material is closely related to things from \cite{BKM} and \cite{LM}, and parts of it is at least implicitly contained in these references. 

In the fourth section it is shown how one canonical cover can be used to get new canonical covers. This is based on a variation of the so-called subset construction described in \S 3.3 of \cite{LM} in which only some of the arrows are allowed. In Section 5 it is shown that a restriction to a subset of the vertexes in the subset construction, applied to an arbitrary presentation of the sofic shift, gives rise to a labeled graph whose merged graph is isomorphic to the future cover. Thus the future cover is the factor obtained from a subgraph of the subset construction by identifying vertexes with the same follower set. This fact is used to produce the examples occurring in the paper. 

The paper closes with three examples where the presentations of the covers obtained in Sections 4 and 5 are described starting from the minimal right-resolving presentations of three mixing sofic shifts. In this way it is shown that the new canonical cover in general is different from the future cover. 
\section{A canonical cover of the non-wandering part of the future cover}\label{cancover}

Let $(G,L_G)$ be a right-resolving presentation of the sofic shift $Y$. The labeled graphs $(G',L_{G'})$ and $(\underline{G},L_{\underline G})$ are both subgraphs of $(\overline{G},L_{\overline{G}})$. While $(\underline{G},L_{\underline{G}})$ is obtained from $(\overline{G},L_{\overline{G}})$ by restricting to a subset of the vertexes in $\overline{G}$, $(G',L_{G'})$ comes from a restriction both to a subset of vertexes and arrows in $\overline{G}$. In this section we compare the two subgraphs and show that their intersection gives rise to a  canonical cover of the non-wandering part of the future cover of $Y$.








\begin{lemma}\label{08-10-25bxy} For all $y \in Y$ the element $\beta_G(y)$ is forward asymptotic to $\alpha_G(y)$.
\end{lemma}
\begin{proof} Note that
$$
s_{\underline{G}}(\alpha_G(y)_n) = \left\{t_G(x): \ x \in X_G(-\infty,n-1], \ L_G(x) = y_{(-\infty,n-1]}\right\} \supseteq  s_{G''}(\beta_G(y)_n) 
$$
for all $n \in \mathbb Z$. A compactness argument shows that there is an $N \in \mathbb Z$ such that
$$
\left\{t_G(x): \ x \in X_G(-\infty,n-1], \ L_G(x) = y_{(-\infty,n-1]}\right\} = s_{G''}(\beta_G(y)_n) 
$$
for all $n \geq N$. It follows that $\alpha_G(y)_j = \beta_G(y)_j$ for all $j \geq N$.
\end{proof}

When $y$ is periodic, both $\alpha_G(y)$ and $\beta_G(y)$ are periodic and hence Lemma \ref{08-10-25bxy} has the following

\begin{corollary}\label{09-10-25b} Let $p \in \Per (Y)$. Then $\beta_G(p) = \alpha_G(p)$.
\end{corollary}


\begin{lemma}\label{10-02-25dxxy} Let $F \in V_{\underline{G}}$. There is a periodic point $p \in \Per (Y)$ and an element $y' \in X_{\underline{G}}$ backward asymptotic to $\alpha_G(p)$ such that $t_{\underline{G}}(y'_{(-\infty,-1]}) = F$.
\end{lemma}
\begin{proof} Since $F \in V_{\underline{G}}$ there is an element $y \in Y$ such that $F = s_{\underline{G}}(\alpha_G(y)_0)$. Set $H_j := s_{\underline{G}}(\alpha(y)_j), \ j \leq 0$. Since $(G,L_G)$ is right-resolving there is an $N <0$ such that 
\begin{equation}\label{17-04-25bxy}
\# H_j = \# H_N
\end{equation} 
for all $j \leq N$.
 By the pigeon hole principle there is a sequence $\cdots < n_3 < n_2 <n_1 < n_0 \leq N$ and a vertex $H \in V_{\underline{G}}$ such that $H =s_{\underline{G}}(\alpha_G(y)_{n_j})$ for all $j \in \mathbb N \cup \{0\}$. Consider the loops in $\underline{G}$:
$$
\gamma_j := H \overset{y_{n_j}}{\to} H_{n_j+1}  \overset{y_{n_j+1}}{\to} \cdots \cdots \overset{y_{n_0-1}}{\to} H .
$$
If $\# L_G^{-1}(L_{\underline{G}}(\gamma_j)^\infty) > \# H$ for all $j$ it follows that there is an element $x' \in X_{\overline{G}}(-\infty,{n_0}-1]$ such that $H \subsetneq t_{\overline{G}}(x')$ and $L_{\overline{G}}(x') = y_{(-\infty,n_0-1]}$. By definition of $\alpha_G(y)$ this implies that $\# s_{\underline{G}}(\alpha_G(y)_j)  = \# H_j > \# H =\# H_N $ for all $j \leq n_0$, contradicting \eqref{17-04-25bxy}. This shows that there is a $j$ such that $q:= L_{\underline{G}}((\gamma_j)^\infty) \in Y$ is periodic and $\# L_G^{-1}(q) = \#H$. Thus $\alpha_G(q) = (\gamma_j)^\infty$ and 
$$
y' := \alpha_G(\sigma^{n_0}(q))_{(-\infty,n_0-1]}\alpha_G(y)_{[n_0,\infty)}
$$
has the stated properties.


\end{proof}


\begin{lemma}\label{11-11-24axy} Let $x,y \in X_{\underline{G}}$, and assume that $x = \alpha_G(L_{\underline{G}}(x))$. If there is an $N \in \mathbb Z$ such that $x_i = y_i$ for $i \leq N$, then $y = \alpha_G(L_{\underline{G}}(y))$.
\end{lemma}
\begin{proof} Since $L_{\underline{G}}(x)_i = L_{\underline{G}}(y)_i$ for $i \leq N$, it follows from the definition of 
$\alpha_G$ that 
$$
\alpha_G(L_{\underline{G}}(x))_j = \alpha_G(L_{\underline{G}}(y))_j
$$ 
for $j \leq N-1$. Hence 
$$
t_{\underline{G}}(y_{N-1}) =t_{\underline{G}}(x_{N-1}) = t_{\underline{G}}\left(\alpha_G(L_{\underline{G}}(x))_{N-1}\right) =  t_{\underline{G}}\left(\alpha_G(L_{\underline{G}}(y))_{N-1}\right).
$$
Since $L_{\underline{G}}(y_{[N,\infty)}) = L_{\underline{G}}\left(\alpha_G(L_{\underline{G}}(y))_{[N,\infty)}\right)$ it follows that $y_{[N,\infty)} = \alpha_G(L_{\underline{G}}(y))_{[N,\infty)}$ because $(\underline{G},L_{\underline{G}})$ is right resolving.  We know already that $y_{(-\infty,N-1]} = x_{(-\infty,N-1]} = \alpha_G(L_{\underline{G}}(x))_{(-\infty,N-1]} =  \alpha(L_{\underline{G}}(y))_{(-\infty,N-1]}$, and conclude therefore that $y = \alpha_G(L_{\underline{G}}(y))$.
\end{proof}


\begin{lemma}\label{03-07-25axy} Let $x \in X_{\underline{G}}$. There is an element $y \in Y$ such that $\beta_G(y)$ is forward asymptotic to $x$.
\end{lemma}
\begin{proof} By
Lemma \ref{10-02-25dxxy} there is an element $y' \in X_{\underline{G}}$ and a periodic point $p\in Y$ such that $y'$ is backward asymptotic to $\alpha_G(p)$ and forward asymptotic to $x$. Since $\alpha_G(p) =  \alpha_G(L_{\underline{G}}(\alpha_G(p)))$, it follows from Lemma \ref{11-11-24axy} that $y' = \alpha_G(L_{\underline{G}}(y'))$. Set $y := L_{\underline{G}}(y')$ and note that $\beta_G(y)$ is forward asymptotic to $\alpha_G(y) = y'$ by Lemma \ref{08-10-25bxy}. Hence $\beta_G(y)$ is also forward asymptotic to $x$.
\end{proof}


\begin{lemma}\label{08-10-25ay}
$$
\overline{\Per X_{\underline{G}}}
 = \bigcup_{y \in Y} \bigcap_{N \geq 1} \overline{\left\{\sigma^n(\beta_G(y)) : \ n \geq N \right\}} .
$$
\end{lemma}
\begin{proof} Let $y \in Y$. It follows from Lemma \ref{08-10-25bxy} that
$$
\bigcap_{N \geq 1} \overline{\left\{\sigma^n(\beta_G(y)) : \ n \geq N \right\}} = \bigcap_{N \geq 1} \overline{\left\{\sigma^n(\alpha_G(y)) : \ n \geq N \right\}}.
$$
Since there is a component of $X_{\underline{G}}$ which contains $\sigma^n(\alpha_G(y))$ for all sufficiently large $n$, it follows that 
$$
 \bigcap_{N \geq 1} \overline{\left\{\sigma^n(\beta_G(y)) : \ n \geq N \right\}} \subseteq \overline{\Per X_{\underline{G}}} .
 $$
This proves the inclusion $\supseteq$. The obtain the reversed inclusion, let $C$ be a component in $X_{\underline{G}}$. Choose an element $x \in X_C$ which is forward transitive in $X_C$; i.e. every word from $\mathbb W(X_C)$ occurs infinitely many times to the right in $x$. It follows from Lemma \ref{03-07-25axy} that there is an element $y \in Y$ such that 
\begin{equation}\label{08-10-25dx}
\lim_{n \to \infty} d\left(\sigma^n(\beta_G(y)),\sigma^n(x)\right) = 0.
\end{equation}
For each $z \in X_C$ there is a strictly increasing sequence $\{n_i\}$ in $\mathbb N$ such that
$\lim_{i \to \infty} d(\sigma^{n_i}(x),z) = 0$.
Thanks to \eqref{08-10-25dx}, by passing to a subsequence we can then arrange that
$$
\lim_{i \to \infty} d(\sigma^{n_i}(\beta_G(y)),z) = 0.
$$
It follows that $z \in  \bigcap_{N \geq 1} \overline{\left\{\sigma^n(\beta_G(y)) : \ n \geq N \right\}}$.
\end{proof}


\begin{corollary}\label{01-11-25}  $\overline{\Per X_{\underline{G}}} \subseteq X_{G'} \cap X_{\underline{G}}$
\end{corollary}




Let $\psi : Y \to Z$ be a conjugacy of sofic shifts. Let $(G,L_G)$ and $(H,L_H)$ be right-resolving presentations of $Y$ and $Z$, respectively. We assume that there is a conjugacy $\phi : X_G \to X_H$ such that
\begin{equation}\label{09-10-25ax}
\xymatrix{
 X_G\ar[r]^-{\phi} \ar[d]_{L_G}& X_{H} \ar[d]^-{L_H}\\
   Y \ar[r]^{\psi} & Z}
\end{equation}
commutes.



\begin{lemma}\label{09-10-25f} 
$$
\phi'( \overline{\Per X_{\underline{G}}}) = \overline{\Per X_{\underline{H}}} .
$$
and 
$$
f_{\underline{H}} \circ \phi' = \phi_{\mathbb K} \circ f_{\underline{G}}
$$
on $\overline{\Per X_{\underline{G}}}$.
\end{lemma}
\begin{proof} The first equality follows from Lemma \ref{08-10-25ay} and \eqref{30-10-25d}. Let $z \in \overline{\Per X_{\underline{G}}}$. By Lemma \ref{08-10-25ay} there is an $y \in Y$ such that $z = \lim_{i \to \infty} \sigma^{n_i}(\beta_{G}(y)) $ for some sequence $\{n_i\}$ in $\mathbb N$.
Note that 
$$
\lim_{i \to \infty} \sigma^{n_i}(\beta_{H}(y)) = \lim_{i \to \infty} \sigma^{n_i}(\alpha_{H}(y))
$$
by Lemma \ref{08-10-25bxy}. From \eqref{30-10-25d} we see that 
$\phi' \circ \beta_G = \beta_H \circ \psi$. 
Hence
\begin{align*}
&\phi'( \lim_{i \to \infty} \sigma^{n_i}(\beta_{G}(y))) = \lim_{i \to \infty} \sigma^{n_i}(\beta_{H}(\psi(y)) =  \lim_{i \to \infty} \sigma^{n_i}(\alpha_{H}(\psi(y)).
\end{align*}
Note that $f_{\underline{H}} \circ \alpha_H \circ \psi = \alpha_{Z} \circ \psi$ by Lemma \ref{11-09-25ex}. Since $\alpha_{Z}\circ\psi = \phi_{\mathbb K} \circ \alpha_Y$ by \eqref{24-11-24x} we conclude that
\begin{align*}
&f_{\underline{H}} \circ \phi'(z) = f_{\underline{H}} \circ \phi'( \lim_{i \to \infty} \sigma^{n_i}(\beta_{G}(y))) = \lim_{i \to \infty} f_{\underline{H}}(\sigma^{n_i}(\alpha_{H}(\psi(y))) \\
&  = \lim_{i \to \infty} \sigma^{n_i}\left( f_{\underline{H}} \circ \alpha_H \circ \psi(y)\right)= \lim_{i \to \infty}  \sigma^{n_i}(\alpha_{Z}(\psi(y)))= \lim_{i \to \infty}  \sigma^{n_i}(\phi_{\mathbb K}\circ\alpha_Y(y)) \\
& =  \phi_{\mathbb K}\left(\lim_{i \to \infty} \sigma^{n_i}(\alpha_{Y}(y))\right) =  \phi_{\mathbb K}\left(\lim_{i \to \infty} \sigma^{n_i}(f_{\underline{G}} \circ \alpha_{G}(y))\right) = \phi_{\mathbb K}\circ f_{\underline{G}}(z).
\end{align*}
%For comparison 
%\begin{align*}
%&\phi \circ f_{G}\left( \lim_{i \to \infty} \sigma^{n_i}(\beta_{G}(y))\right) \\
%& = \phi\left(  \lim_{i \to \infty} \sigma^{n_i}(f_{G} \circ \beta_{G}(y))\right)\\
%& = \phi \left( \lim_{i \to \infty} \sigma^{n_i}(\alpha_Y(y))\right)\\
%&  = \phi \left( \lim_{i \to \infty} \sigma^{n_i}(f_G\circ \alpha_G(y))\right)\\
%&= \phi \circ f_G(z) 
%\end{align*}
\end{proof}


This leads to the following:

\begin{thm}\label{09-10-25g} Let $\psi : Y \to Z$ be a conjugacy of sofic shifts. Assume that $(G,L_G)$ and $(H,L_H)$ are right-resolving presentations of $Y$ and $Z$, respectively. Assume that there is a conjugacy $\phi : X_G \to X_H$ such that \eqref{09-10-25ax} commutes. Then 
$$
\phi' : \ \overline{\Per X_{\underline{G}}} \to \overline{\Per X_{\underline{H}}}
$$
is a conjugacy and the diagram 
\begin{equation}\label{09-10-25x}
\xymatrix{\\
\overline{\Per (X_{\underline{G}})} \ar[r]^-{\phi'} \ar[d]_-{f_{\underline{G}}} & \overline{\Per (X_{\underline{H}})}  \ar[d]^-{f_{\underline{H}}}\\
 \overline{\Per X_{\mathbb K(Y)}} \ar[d]_-{L_{\mathbb K(Y)}} \ar[r]^-{\phi_{\mathbb K}} &  \overline{\Per X_{\mathbb K(Z)}} \ar[d]^-{L_{\mathbb K(Z)}} \\
\overline{\Per Y} \ar[r]_-{\psi}  & \overline{\Per Z }}
\end{equation}
commutes.


\end{thm}

In general the conjugacies $\phi_{\mathbb K}$ and $\phi'$ in the diagram \eqref{09-10-25x} are not unique.


Note that $\overline{\Per(X_G)} = X_{\hat{G}}$ where 
$$
\hat{G} := \bigsqcup_{C \in V_{\underline{G}}^r/\sim} G_C \subseteq \underline{G}.
$$
It it therefore natural to ask

\begin{question}\label{13-10-25} Does the conjugacy $\phi'$ of Theorem \ref{09-10-25g} extend to a conjugacy $\lambda : X_{\underline{G}} \to X_{\underline{H}}$ such that
\begin{equation*}\label{13-10-25a}
\xymatrix{\\
X_{\underline{G}} \ar[r]^-{\lambda} \ar[d]_-{f_{\underline{G}}} & X_{\underline{H}}  \ar[d]^-{f_{\underline{H}}}\\
 X_{\mathbb K(Y)} \ar[d]_-{L_{\mathbb K(Y)}} \ar[r]^-{\phi_{\mathbb K}} &  X_{\mathbb K(Z)} \ar[d]^-{L_{\mathbb K(Z)}} \\
Y \ar[r]_-{\psi}  & Z }
\end{equation*}
commutes ?
\end{question}

+++++
 It follows from Lemma \ref{25-06-25ax} that 
\begin{equation*}\label{30-10-25bxx}
\xymatrix{
 X_G^{L_G} \ar[r]^-{\phi^{\#}} \ar[d]_{L_G^{\#}}&  X_H^{L_H} \ar[d]^{L_H^{\#}} \\
 Y  \ar[r]^-\psi & Z}
\end{equation*}
commutes. When we use the same symbol $\kappa$ for the two conjugacies $\kappa : X_G^{L_G} \to X_{G''}$ and $\kappa : X_H^{L_H} \to X_{H''}$ from Lemma \ref{23-06-25dyx} and set $\phi'' = \kappa \circ \phi^\# \circ \kappa^{-1}$, it follows from Lemma \ref{23-06-25dyx} that $\phi'' : X_{G''} \to X_{H''}$ is a conjugacy such that
\begin{equation*}\label{30-10-25bxxx}
\xymatrix{
 X_{G''} \ar[r]^-{\phi''} \ar[d]_{L_{G''}}&  X_{H''} \ar[d]^{L_{H''}} \\
 Y  \ar[r]^-\psi & Z}
\end{equation*}
commutes. Since $\phi^\#(L_G^{-1}(y)) = L_H^{-1}(\psi(y))$ it follows that $\phi''(\beta_G(y)) = \beta_H(\psi(y))$ for all $y \in Y$. By Lemma \ref{31-10-25f} this implies that $\phi''(X_{G'}) = X_{H'}$, and we set $\phi':= \phi''|_{X_{G'}}$. This proves the existence part. To prove uniqueness, consider a conjugacy $\lambda : X_{G'} \to X_{H'}$ such that
\begin{equation*}\label{30-10-25cy}
\xymatrix{
 X_{G'} \ar[r]^-{\lambda} \ar[d]_{L_{G'}}&  X_{H'} \ar[d]^{L_{H'}} \\
Y   \ar[r]^-\psi & Z}
\end{equation*}
commutes. Set 
$$
\lambda' := \kappa^{-1}\circ \lambda \circ \kappa|_{\kappa^{-1}(X_{G'})}.
$$ 
By Lemma \ref{23-06-25dyx} also
\begin{equation*}\label{30-10-25cxx}
\xymatrix{
 X_{G}^{L_G}\cap \kappa^{-1}(X_{G'}) \ar[r]^-{\lambda'} \ar[d]_{L_{G}^{\#}}&  X_{H}^{L_H}\cap \kappa^{-1}(X_{H'}) \ar[d]^{L_{H}^{\#}} \\
Y   \ar[r]^-\psi & Z}
\end{equation*}
commutes. 


Note that $L_G^{-1}(y) \in X_G^{L_G} \cap \kappa^{-1}(X_{G'})$ by Lemma \ref{31-10-25f}. It follows that $\lambda'(L_G^{-1}(y)) \subseteq L_H^{-1}(\psi(y))$ and ${\lambda'}^{-1}(L_H^{-1}(\psi(y))) \subseteq L_G^{-1}(y)$, implying that $\lambda'(L_G^{-1}(y)) =L_H^{-1}(\psi(y))$ for all $y \in Y$ and hence $\lambda(\beta_G(y)) = \beta_H(\psi(y)) = \phi'(\beta_G(y))$ for all $y \in Y$. By Lemma \ref{31-10-25f} this implies that $\lambda = \phi'$.

\begin{ex}

The following labeled graph $(G,L_G)$ is irreducible, right resolving and follower-separated and hence the minimal right resolving graph of the sofic shift $Y$ it represents. Note that $Y$ is mixing.

\begin{equation}\label{05-11-25}
 \begin{tikzpicture}[node distance={35mm}, thick, main/.style = {draw, circle}] 
\node[main] (1) {$a$}; 
\node[main] (2) [right of=1] {$b$};
%\node[main] (3) [below of=1] {$\{a,b\}$}; 
\draw[->] (1) to [out=90,in=170,looseness=5] node[above ,pos=0.5] {$0$} (1);
\draw[->] (1) to [out=190,in=270,looseness=5] node[below ,pos=0.5] {$2$} (1);
\draw[->] (2) to [out=10,in=90,looseness=5] node[above ,pos=0.5] {$0$} (2);
\draw[->] (2) to [out=0,in=270,looseness=5] node[below ,pos=0.5] {$1$} (2);
\draw[->] (1) to [out=300,in=230,looseness=1] node[above ,pos=0.5] {$1$} (2);
\draw[->] (2) to [out=120,in=70,looseness=1] node[above ,pos=0.5] {$2$} (1);
 \end{tikzpicture} 
\end{equation}


In this case the graphs $(\underline{G},L_{\underline{G}})$ and $(G',L_{G'})$ are the same:



\begin{equation*}
 \begin{tikzpicture}[node distance={35mm}, thick, main/.style = {draw, circle}] 
\node[main] (1) {$\{a\}$}; 
\node[main] (2) [right of=1] {$\{b\}$};
\node[main] (3) [below of=1] {$\{a,b\}$}; 
\draw[->] (1) to [out=80,in=170,looseness=3] node[above ,pos=0.5] {$0$} (1);
\draw[->] (1) to [out=190,in=270,looseness=3] node[below ,pos=0.5] {$2$} (1);
\draw[->] (2) to [out=10,in=90,looseness=3] node[above ,pos=0.5] {$0$} (2);
\draw[->] (2) to [out=0,in=290,looseness=3] node[below ,pos=0.5] {$1$} (2);
\draw[->] (1) to [out=310,in=230,looseness=1] node[above ,pos=0.5] {$1$} (2);
\draw[->] (2) to [out=120,in=70,looseness=1] node[above ,pos=0.5] {$2$} (1);
\draw[->] (3) to [out=230,in=310,looseness=4] node[below ,pos=0.5] {$0$} (3);
\draw[->] (3) to [out=90,in=290,looseness=1] node[left ,pos=0.5] {$2$} (1);
\draw[->] (3) to [out=60,in=260,looseness=1] node[right ,pos=0.5] {$1$} (2);
 \end{tikzpicture} 
\end{equation*}
\end{ex}

Since the vertex $\{a,b\}$ has the same follower set as $\{a\}$, the future cover is again the labeled graph \eqref{05-11-25} we started with,  and hence equal to the minimal right-resolving cover, while the canonical cover $(G',L_{G'})$ is bigger.




++++++++++++++++++

+++++++++++++++++++++++


++++++++

 For this note that there is a periodic point $p$ in $X_G$ such that \begin{equation}\label{19-01-26a}
s_G(p_j) \in s_{G''}(\tilde{x}_j)
\end{equation} 
for all $j \in \mathbb Z$ and an element $z \in X_G$ such that $L_G(z) = L_{\underline{G}}(x)$ and $z_j = p_j$ for $j \in [i'-3M,i']$. 



It follows from \eqref{19-01-26a} and the definition of $\phi''$ that $s_H(\phi(p)_j) \in s_{H''}(\phi''(\tilde{x})_j)$ for all $j \in \mathbb Z$ and hence also that $s_H(\phi(z)_j) \in s_{H''}(\phi''(\tilde{x})_j)$ for $j \in [i'-2M,i'-M]$. Since $L_H(\phi(z))_j = \psi(L_G(z))_j = \psi(L_{\underline{G}}(x))_j$ for all $j$, we see that $v:= s_H(\phi(z)_{i'-M})$ is vertex with properties we seek and that $\gamma:= \phi(z)_{[i'-M,\infty)}$ is the corresponding right-infinite path. The existence of this pair and the definition $\underline{H}$ imply that there is a unique element $y \in X_{\underline{H}}[i'-M,\infty)$ such that $s_{\underline{H}}(y_{i'-M}) = s_{\underline{H}}(\phi''(\tilde{x})_{i'-M}) =  s_{H''}(\phi'(x_{[i'-2M,i']}))$ and $L_{\underline{H}}(y) = \psi(L_{\underline{G}}(x))_{[i'-M,\infty)}$ . Now define $\tilde{\phi}(x)_i$ by
$$
\tilde{\phi}(x)_i : = y_ i.
$$
With this definition $\tilde{\phi}(x)_i$ depends only on $x_{[i-N,i+M]}$ and gives us therefore a sliding block code $\tilde{\phi} : \ X_{\underline{G}} \to X_{\underline{H}}$ with memory $N$ and anticipation $M$. By construction $L_{\underline{H}}(\tilde{\phi}(x)_i) = L_{\underline{H}}(y)_i = \psi(L_{\underline{G}}(x))_i$, proving that $L_{\underline{H}} \circ \tilde{\phi} = \psi \circ L_{\underline{G}}$. Finally, note that when $x \in X_C$ for some component $C$ in $\underline{H}$ the element $y$ is equal to $\phi(x)_{[i'-M,\infty)}$ and hence $\tilde{\phi}(x)_i = \phi(x)_i$, showing that $\tilde{\phi}$ extends $\phi$.

\end{document}
